\documentclass[11pt,letterpaper]{article}
\usepackage[T1]{fontenc}
\usepackage[utf8]{inputenc}
\usepackage{lmodern}
\usepackage[margin=1in,headheight=15pt]{geometry}
\usepackage{amsmath,amssymb,amsthm,mathtools,mathrsfs}
\usepackage{microtype,booktabs,longtable,array,enumitem}
\usepackage{ragged2e}
\usepackage{xcolor,fancyhdr}
\definecolor{ink}{HTML}{283A35}
\definecolor{accent}{HTML}{49675A}
\usepackage[colorlinks=true,linkcolor=ink,citecolor=accent,urlcolor=accent]{hyperref}
\usepackage{aliascnt}
\usepackage[nameinlink,noabbrev]{cleveref}
\hypersetup{%
  pdftitle={Surquaternions: Algebra, Infinitesimal Geometry, and Support-Controlled Analysis},
  pdfauthor={Merged research report prepared for Vladimir Reshetnikov},
  pdfsubject={Hamilton quaternions over the surreal numbers}}
\pagestyle{fancy}\fancyhf{}
\fancyhead[L]{\small\textcolor{ink}{Surquaternions}}
\fancyhead[R]{\small\thepage}
\renewcommand{\headrulewidth}{0.3pt}
\setlength{\parskip}{0.3em}
\setlength{\emergencystretch}{3em}
\setlist{itemsep=0.15em,topsep=0.3em}
\numberwithin{equation}{section}
\newtheorem{theorem}{Theorem}[section]
\newaliascnt{proposition}{theorem}
\newtheorem{proposition}[proposition]{Proposition}
\aliascntresetthe{proposition}
\crefname{proposition}{Proposition}{Propositions}
\newaliascnt{lemma}{theorem}
\newtheorem{lemma}[lemma]{Lemma}
\aliascntresetthe{lemma}
\crefname{lemma}{Lemma}{Lemmas}
\newaliascnt{corollary}{theorem}
\newtheorem{corollary}[corollary]{Corollary}
\aliascntresetthe{corollary}
\crefname{corollary}{Corollary}{Corollaries}
\theoremstyle{definition}
\newaliascnt{definition}{theorem}
\newtheorem{definition}[definition]{Definition}
\aliascntresetthe{definition}
\crefname{definition}{Definition}{Definitions}
\newaliascnt{example}{theorem}
\newtheorem{example}[example]{Example}
\aliascntresetthe{example}
\crefname{example}{Example}{Examples}
\newaliascnt{convention}{theorem}
\newtheorem{convention}[convention]{Convention}
\aliascntresetthe{convention}
\crefname{convention}{Convention}{Conventions}
\theoremstyle{remark}
\newaliascnt{remark}{theorem}
\newtheorem{remark}[remark]{Remark}
\aliascntresetthe{remark}
\crefname{remark}{Remark}{Remarks}
\newaliascnt{warning}{theorem}
\newtheorem{warning}[warning]{Warning}
\aliascntresetthe{warning}
\crefname{warning}{Warning}{Warnings}
\newcommand{\No}{\mathbf{No}}
\newcommand{\SQ}{\mathbb H_{\No}}
\newcommand{\HF}{\mathbb H_F}
\newcommand{\HR}{\mathbb H_{\mathbb R}}
\newcommand{\R}{\mathbb R}
\newcommand{\C}{\mathbb C}
\newcommand{\Q}{\mathbb Q}
\newcommand{\Z}{\mathbb Z}
\newcommand{\N}{\mathbb N}
\newcommand{\OO}{\mathcal O}
\newcommand{\mm}{\mathfrak m}
\newcommand{\Oz}{\mathbf{Oz}}
\newcommand{\eps}{\varepsilon}
\newcommand{\ii}{\mathbf i}
\newcommand{\jj}{\mathbf j}
\newcommand{\kk}{\mathbf k}
\newcommand{\Rea}{\operatorname{Re}}
\newcommand{\Ima}{\operatorname{Im}}
\newcommand{\st}{\operatorname{st}}
\newcommand{\supp}{\operatorname{supp}}
\newcommand{\lc}{\operatorname{lc}}
\newcommand{\fin}{\operatorname{fin}}
\newcommand{\tr}{\operatorname{tr}}
\newcommand{\Aut}{\operatorname{Aut}}
\newcommand{\GL}{\operatorname{GL}}
\newcommand{\SO}{\operatorname{SO}}
\newcommand{\SU}{\operatorname{SU}}
\newcommand{\Sp}{\operatorname{Sp}}
\newcommand{\ad}{\operatorname{ad}}
\newcommand{\diag}{\operatorname{diag}}
\newcommand{\ord}{\operatorname{ord}}
\newcommand{\Exp}{\operatorname{Exp}}
\newcommand{\Log}{\operatorname{Log}}
\newcommand{\Expm}{\Exp_{\mm}}
\newcommand{\Logm}{\Log_{\mm}}
\newcommand{\Expt}{\Exp_{\mathcal T}}
\newcommand{\Expr}{\Exp_{\mathrm{rad}}}
\newcommand{\Sin}{\operatorname{Sin}}
\newcommand{\Cos}{\operatorname{Cos}}
\newcommand{\BCH}{\operatorname{BCH}}
\newcommand{\Span}{\operatorname{span}}
\newcommand{\abs}[1]{\left|#1\right|}
\newcommand{\norm}[1]{\left\|#1\right\|}
\newcommand{\angles}[1]{\left\langle#1\right\rangle}
\newcommand{\DD}{\mathcal D}
\newcommand{\dsl}{\partial_{\mathrm{sl}}}
\newcommand{\dH}{\partial_{\mathbb H}}
\newcommand{\delH}{\delta_{\mathbb H}}
\newcommand{\id}{\operatorname{id}}
\newcolumntype{L}[1]{>{\RaggedRight\arraybackslash}p{#1}}
\begin{document}
\hypersetup{pageanchor=false}
\pagenumbering{roman}
\begin{titlepage}
\centering
\vspace*{0.5cm}
{\large\scshape Algebra at every surreal scale\par}
\vspace{0.7cm}
{\Huge\bfseries\textcolor{ink}{Surquaternions}\par}
\vspace{0.35cm}
{\Large Algebra, Infinitesimal Geometry,\\[0.2em]and Support-Controlled Analysis\par}
\vspace{0.7cm}
{\large September 21, 2026\par}
\vspace{0.6cm}
\begin{minipage}{0.95\textwidth}
\begin{abstract}
A surquaternion is an element of
$\SQ=\No\oplus\No\ii\oplus\No\jj\oplus\No\kk$ with Hamilton's multiplication. The
definition is elementary; a usable theory must also say which parts of quaternionic
geometry survive non-Archimedean scales, what an infinite sum denotes, and which
notions of differentiation are compatible with the surreal field. The algebra is
noncommutative, so factor order must be retained in every extension. Commutative
results apply within explicitly identified central fields or complex slices; the
quaternionic extensions require the additional arguments given here.

We develop the division algebra and its ordered-field-valued modulus, algebraically
closed complex slices, conjugacy spheres, the algebraic rotation and spin description,
one-sided polynomials with a constructive fundamental theorem and a complete
zero-class algorithm, and finite-dimensional Hermitian spectral theory with a
projector estimate at arbitrary surreal gap scale. We identify surquaternions with
set-supported Hahn series having ordinary quaternion coefficients, prove the valuation,
leading-coefficient and standard-part formulas with the \emph{noncommutative} residue
algebra $\HR$, and separate that summation theory from the exceptionally fine topology
of the full surreal class. Positive support then supplies an exact infinitesimal
exponential and logarithm, Baker--Campbell--Hausdorff at positive valuation, a filtered
nonabelian rotation group with a graded bracket, a support-controlled multivariable
implicit lifting theorem, and an exact square-root conditioning analysis. For analysis
we keep slice regularity, Fueter regularity and field derivations apart, and prove
coefficientwise representation, Cauchy and Fueter results on explicitly named
coefficient spaces.

Two results are negative and sharp. The componentwise Berarducci--Mantova derivation
admits no nonzero solution of $\dH y=\ii y$, so no global quaternionic exponential can
satisfy the usual commuting chain rule for that derivation at $\omega\ii$. And the
global radial exponential built from the Ehrlich--Kaplan scalar normalization, whose
complete logarithm fibres and four-coordinate differential we compute --- finding
rank-two critical spheres at radii in $\pi\Oz$, including the infinite radius
$\omega$ --- is not an additive-to-multiplicative homomorphism on noncommuting
arguments. Three exponentials and four derivative-like operators are therefore fixed by
name and symbol in \cref{squat:sec:ledger} and never identified.
\end{abstract}
\end{minipage}
\vfill
\begin{minipage}{0.95\textwidth}\small
\textbf{Status.} A self-contained mathematical development, with cited foundational
inputs and proofs of the stated extensions. This is not a refereed paper, a priority
claim, or a claim to solve a named published open problem. The accompanying symbolic
checks verify finite identities, finite formal expansions and worked examples, not the
general theorems. No Lean verification of this article is claimed.

\textbf{Provenance.} This is a merged report. It is the union of two independently
written manuscripts, \emph{noncommutative-hahn-obstructions} and
\emph{radial-exponential-support-lifting}, both dated September 21, 2026 and each
delivered as a standalone archive with its own symbolic suite. Neither manuscript is
reproduced here. Every result either
manuscript reached is printed here; where the two proved the same statement by
genuinely different routes, both routes are kept and marked; every honest limitation
from either manuscript survives in \cref{squat:sec:nonclaims}.

\textbf{Repository baseline.} The source repository is Vladimir Reshetnikov's
\texttt{Surreal}, read at commit
\texttt{39f2be6667ade51bca2b45daa47e289d69c09764}, accessed September 21, 2026. This
article is a new standalone companion, not a modification of that repository.
\end{minipage}
\end{titlepage}
\tableofcontents
\clearpage
\pagenumbering{arabic}
\hypersetup{pageanchor=true}

\section{Scope, provenance, and the choice of analytic structure}\label{squat:sec:scope}

\subsection{What is being generalized}\label{squat:sub:whatgen}
Hamilton's quaternions do not require real completeness for their multiplication. They
require an ordered field to make a sum of four squares positive, and a real closed field
to obtain the square roots and complex factorizations used in much of their geometry.
The surreal numbers supply both, while also supplying infinitely large and infinitesimal
coefficients.

The basic object is therefore
\begin{equation}\label{squat:eq:defn}
 \SQ=\{a+b\ii+c\jj+d\kk:a,b,c,d\in\No\},\qquad
 \ii^2=\jj^2=\kk^2=\ii\jj\kk=-1 .
\end{equation}
We use \emph{surquaternion} descriptively, without asserting that the word or the
construction is new. A surquaternion is an ordered quadruple of surreals, not a new
Conway cut with quaternionic left and right options. Throughout, the four coefficients
are \emph{central}: they commute with $\ii,\jj,\kk$ and with every surquaternion. This
is not an algebra of games with a new multiplication, nor an attempt to order the
quaternion algebra. The order used to construct surreal numbers is an order on scalar
components only; it does not become a compatible ordering of the division algebra
(\cref{squat:prop:center}).

The one structural fact that governs the whole report is that $\SQ$ is
\emph{noncommutative}. Conjugation, multiplicativity of the modulus, the meaning of
``valuation'', leading coefficients, polynomial evaluation, matrix adjoints, exponential
laws and every chain rule require separate treatment from the commutative surcomplex
case. Where a companion report in the same repository proves a commutative statement, we
identify its scalar or slice hypotheses before using it, and prove the additional
noncommutative assertions over $\mathbb H$.

\subsection{Three levels of scalars}\label{squat:sub:levels}
Three levels organize the development, and a theorem at one level is never silently
promoted to another.

\begin{center}\small
\begin{tabular}{L{0.16\textwidth}L{0.36\textwidth}L{0.40\textwidth}}
\toprule
Level & Available structure & What must not be inferred\\
\midrule
$\HF$, $F$ real closed & Finite algebra, ordered-field-valued modulus, polynomial and
matrix geometry & Real compactness, completeness, or unrestricted analytic transfer\\
$D_\Gamma$, a set-sized Hahn workspace & Well-ordered set-sized supports, valuation,
exact Hahn evaluation & That every power series is evaluable everywhere, or that the
workspace is closed under $\exp_{\No}$ or under a chosen derivation\\
$\SQ$, the full class & All surreal scales and Conway normal forms & That ordinary
convergent sequences model fine limits, or that class-indexed summation is available\\
\bottomrule
\end{tabular}
\end{center}

Level~1 is finite algebra valid over any real closed field $F$; this is the right level
of generality, since the proofs then apply both to $\No$ and to suitable set-sized
subfields. Level~2 adds well-ordered, set-sized supports and exact infinite summation.
Level~3 consists of class statements. Analytic constructions additionally specify a
domain, a support condition, and which differentiation is meant.

\subsection{Relation to companion reports in the repository}\label{squat:sub:companions}
Three committed reports in the same repository carry material that this one extends. We
inherit and cite them rather than restating them, and we print only what
noncommutativity changes.

\emph{Surcomplex Spectral Theory: Singular Values, Infinitesimal Rank, and Multiscale
Stability}, \texttt{docs/surcomplex/spectral-theory/article.tex} \cite{RepoSpec},
already proves Hermitian and normal diagonalization, positive square roots, the singular
value decomposition, polar decomposition and the variational principles with attained
extrema over a real closed field and its complexification, on the stated ground that
these need real closedness rather than Archimedean completeness; over a Hahn workspace it
already has gap-sensitive spectral projectors and subspace perturbation, with an explicit
warning about Gram squaring. \Cref{squat:sec:spectral} keeps the quaternionic case
separate --- the right-module structure, the adjoint $A^*=\bar A^{\,T}$ and the reality
of eigenvalues of quaternionic Hermitian matrices are genuinely different --- and cites
that report for the complex case. One hard constraint is recorded at
\cref{squat:warn:nodet}: that report's determinantal singular-scale theorem and its
positive Cauchy--Binet argument use commutative minors and cannot be transferred
unchanged to quaternionic matrices.

\emph{Trigonometry on the Surcomplex Plane: Canonical Phases, Arbitrary-Scale Geometry,
and Infinitesimal Degeneration}, \texttt{docs/surcomplex/trigonometry/article.tex}
\cite{RepoTrig}, establishes that every surcomplex direction has a \emph{finite} surreal
angle, and already contains the Ehrlich--Kaplan integer-part normalization with the
omnific integers $\Oz$, presenting it as one member of a character family and explicitly
not claiming it as new. \Cref{squat:sec:angles,squat:sec:exponentials} cite that report
for the scalar normalization and its character-family caveat, and print only the
conjugation-equivariant quaternionic radial extension, its full logarithm fibres, its
four-coordinate differential and its rank-two critical points.

\emph{Differential Algebra and Differential Equations over the Surreal and Surcomplex
Numbers}, \texttt{docs/surcomplex/differential-equations/article.tex} \cite{RepoDiff},
proves ``the missing oscillator'' three independent ways, with the consequence that
$\partial^2y+y=0$ has only $y=0$ and that $\sin\omega$ is not a Berarducci--Mantova
differential element. \Cref{squat:thm:nooscillation} is the componentwise quaternionic
analogue and is stated as an extension of that result, not as an independent discovery.
That report also fixes a three-symbol phase convention precisely because three different
objects were being called ``phase'': the \emph{finite angle} $\theta$, the
\emph{infinitesimal phase residue} $\eta$, and the \emph{accumulated phase} $B$. We
reuse those three symbols in exactly those senses and introduce no fourth sense of the
word; \cref{squat:sec:ledger} is written as an extension of that table.

A fourth companion, the physics report \emph{Surreal Scalars and Spacetime}
\cite{RepoPhys}, asserts the group-homomorphism property of the \emph{commutative}
surcomplex Ehrlich--Kaplan exponential with kernel $2\pi\ii\Oz$.
\Cref{squat:sub:consistency} explains why that claim and the failure of the homomorphism
law recorded in \cref{squat:thm:globalexp} are both true, of different objects.


{\raggedright The report \path{docs/surreal/euclidean-three-space/} treats $\mathrm{SO}(3,\No)$
as $3\times3$ matrices: the split extension (\texttt{e3:so:thm:split}), the
perfect uniquely divisible kernel (\texttt{e3:so:thm:perfect}), finite-subgroup
conjugacy (\texttt{e3:so:thm:finiterigidity}) and the rotation exponential
(\texttt{e3:so:thm:expdifferential}), whose critical spheres have rank one,
unlike the rank-two spheres of \cref{squat:thm:exp-derivative}. It cites the
spin double cover from here rather than reproving it.\par}
\subsection{Foundational inputs and set-sized workspaces}\label{squat:sub:foundations}
We use the established facts that $\No$ is a real closed ordered class field, that every
surreal has a unique Conway normal form with real coefficients and set-sized
reverse-well-ordered exponents, and that Gonshor's exponential is an ordered
exponential-field operation \cite{Conway,Gonshor,vdDE,EK}. No proof here uses a sum
indexed by a proper class.

Class language can be interpreted in a class theory such as NBG; in ZFC the same
notation denotes definable class schemes. In a dependent type theory an ordinary small
type cannot contain every ordinal and hence every surreal, so one works in a fixed
universe or in a set-sized workspace and records the embeddings between workspaces. For
the full surreal class, topological statements use a displayed ball basis, not a
postulated set of all open subclasses. Algebraic group points are finite coordinate
tuples. Statements about arbitrary class automorphisms or derivations concern individual
class maps; we do not postulate a class whose elements are all class functions.

Write
\begin{equation}\label{squat:eq:monomial}
 t^\gamma:=\omega^{-\gamma},
\end{equation}
the \emph{Conway monomial map}. It must not be replaced, for arbitrary surreal $\gamma$,
by a presumed identity with $\exp(-\gamma\log\omega)$ involving Gonshor's exponential.
The two notational uses of powers are kept distinct throughout.

For a set-sized ordered subgroup $\Gamma\subset(\No,+)$ let
\[
 K_\Gamma=\R((t^\Gamma)),\qquad D_\Gamma=\mathbb H_{K_\Gamma}.
\]
Here $K_\Gamma$ denotes the \emph{real} Hahn field (the workspace denoted
$F_\Gamma$ in the repository notation guide), not its complexification. A Hahn series
has a well-ordered support in $\Gamma$, with our increasing-exponent convention.
If $\Gamma$ is divisible, the standard real-closed Hahn-field theorem makes
$K_\Gamma$ real closed \cite{Alling,EK}. That theorem, and the surreal normal-form
theorem that embeds these series in $\No$, are foundational inputs, not consequences of
quaternion multiplication.

\begin{theorem}[Set localization]\label{squat:thm:localization}
Every set $A$ of surquaternions is contained in some $D_\Gamma$ with $\Gamma$ a
divisible, set-sized ordered subgroup of $(\No,+)$. Its elements are represented by their
actual Conway normal forms, not merely by an abstract field embedding.
\end{theorem}
\begin{proof}
Take the union of the supports of all four coordinates of all elements of $A$. This is a
set of surreal exponents. Its rational linear span is a set-sized divisible subgroup
$\Gamma$. All the coordinate series belong to $K_\Gamma$, so all the original
quaternions belong to $D_\Gamma$. The normal-form addition and multiplication identify
the operations there with those in $\SQ$.
\end{proof}

This localization does \emph{not} assert that $K_\Gamma$ is stable under the
exponential, under a chosen surreal derivation, or under every analytic operation later
considered; specified images can be placed in a larger workspace, and such closure is an
additional hypothesis wherever those operations are used. Nor does localization identify
the workspace's intrinsic topology with its subspace topology in the full surreal class:
by \cref{squat:thm:discrete} these are dramatically different.

For fixed finite matrix sizes and polynomial degrees, algebraic statements about $\HF$
can be translated into ordered-field statements about four coordinates per variable, so
first-order real-closed-field transfer is available for such statements, subject to
first-order expressibility \cite{Swan}. It is not a license to transfer compactness,
infinite summation, arbitrary function spaces, or the proper-class assertions about
$\No$. Most proofs below avoid even this shortcut.

\subsection{What is classical and what is developed here}\label{squat:sub:classical}
Much of \cref{squat:sec:algebra,squat:sec:slices,squat:sec:polynomials} is classical
quaternion algebra over a real closed field \cite{Voight}, written out to make its
hypotheses visible; its specialization to $\No$ does not make the division formula, the
conjugacy classification or the spectral theorem new. The polynomial star product,
spherical zero sets, the stem-function construction and the Fueter transformation belong
to the established theory of quaternionic slice regularity \cite{GP}. The scalar global
phase normalization of \cref{squat:sub:ekphase} is due to Ehrlich and Kaplan \cite{EK}
and is not introduced here.

The genuinely surreal content is the class/set boundary, exact support control,
multiscale estimates, and the incompatibilities between otherwise natural analytic
structures. Particularly useful combined results are set localization, the
noncommutative residue calculus, the filtered rotation calculus, the complete zero-class
algorithm, the spectral-projector estimate at arbitrary surreal gap scale,
support-certified implicit lifting, the square-root conditioning formula, the
differential and logarithm fibres of the global radial exponential, and the
no-oscillation theorem with its chain-rule corollary. These are presented as results
proved in this article, not as exhaustively verified claims of first publication. The
literature review is targeted, not a complete bibliographic classification.

\section{The ledger: three exponentials, four derivatives, one name each}\label{squat:sec:ledger}

Read this section before any later one. Three exponentials and four derivative-like
operators in this subject carry overlapping names in the literature and in the two
source manuscripts, and each of them satisfies a different law on a different domain.
Printing one ``exponential'' and one ``derivative'' would manufacture contradictions
between statements that are all true. This section fixes one name and one symbol per
object and records, for each law, exactly which object satisfies it.

\begin{convention}[The three exponentials]\label{squat:conv:exps}
\begin{enumerate}[label=\textup{(\alph*)},leftmargin=2.4em]
\item The \emph{infinitesimal Hahn exponential} $\Expm$, defined on the maximal ideal
$\mm_H$ by the strong Hahn sum of $X^n/n!$, with inverse $\Logm$ on $1+\mm_H$
(\cref{squat:thm:principal-log}). It is a bijection, it commutes with conjugation and
with inner automorphisms, and it is additive \emph{only on commuting pairs}; on
noncommuting pairs the group law is $\BCH$ (\cref{squat:prop:bch}). Source~10 writes
$\Exp_{\mm}$ and source~11 writes $\exp_{\mathrm H}$ for this same object.
\item The \emph{finite-angle tube exponential} $\Expt$, defined on the tube
$\mathcal T=\{a+v:\abs{v}\text{ finite}\}$ as Gonshor's real exponential times canonical
finite-angle trigonometry (\cref{squat:eq:tubeexp}). Its domain contains every finite
quaternion and every central surreal but \emph{excludes} $\omega\ii$. It is
\emph{not} the strong Hahn sum of $q^n/n!$, for two distinct reasons
(\cref{squat:warn:notaylor}), and it is not $\exp(q_0)\Expm(\eps)$ in general. Source~10
writes $\Exp_0$; source~11 does not isolate this intermediate object.
\item The \emph{global radial exponential} $\Expr$, defined on all of $\SQ$ by
\cref{squat:eq:globalexp} using the Ehrlich--Kaplan initial-integer-part normalization
of the scalar phase, applied to the \emph{actual modulus and direction} of the imaginary
part. It is surjective onto $\SQ^\times$, conjugation-equivariant, and obeys the
addition law on commuting pairs; it is \emph{not} an everywhere Hahn-summed power series and
\emph{not} an additive-to-multiplicative homomorphism on the whole noncommutative
algebra, so its fibre over $1$ must not be called a kernel
(\cref{squat:thm:globalexp,squat:thm:logfibres}). Source~10 writes $\Exp_{\mathrm{EK}}$
and source~11 writes $\Exp_{\mathrm{can}}$ for this same object.
\end{enumerate}
The scalar Gonshor exponential on $\No$ is written $\exp_{\No}$, with inverse
$\log_{\No}$; it does obey the unrestricted additive-to-multiplicative law, on the
commutative field $\No$.
\end{convention}

\begin{convention}[The four derivative-like operators]\label{squat:conv:derivs}
\begin{enumerate}[label=\textup{(\alph*)},leftmargin=2.4em]
\item $\dH$, the \emph{componentwise Berarducci--Mantova derivation}: the scalar
derivation $\partial$ on $\No$, normalized by $\partial\omega=1$ and extended
coordinatewise with $\dH\ii=\dH\jj=\dH\kk=0$, so that it differentiates the surreal
coefficients themselves. Its
constant algebra is exactly $\HR$ and it is surjective
(\cref{squat:cor:bmconstants}). Source~10 writes $D$ and source~11 writes $\partial_H$.
More generally $\delH$ denotes the componentwise extension of an \emph{arbitrary} scalar
field derivation $\delta$ (\cref{squat:thm:derivations}), and $\ad_w$ the inner part of a
general derivation.
\item The \emph{coefficientwise slice derivative} $\dsl$, acting on a right-coefficient
series by $\dsl\bigl(\sum_n q^na_n\bigr)=\sum_{n\ge1}nq^{n-1}a_n$
(\cref{squat:prop:evaluation,squat:prop:representation}). This is a formal or slice
derivative in the \emph{variable}; it is not the full algebraic differential, and it is
not $\dH$.
\item The \emph{fine differential} $df_q$, the four-coordinate $\No$-linear
Fr\'echet derivative of \cref{squat:def:fine-derivative}, used to compute
$d\Expr$ in \cref{squat:thm:exp-derivative}. Even the polynomial $q^2$ has fine
differential $h\mapsto qh+hq$ rather than $h\mapsto2hq$
(\cref{squat:prop:sylvester}). A nonsingular fine differential is never used here as an
inverse function theorem.
\item The \emph{left Fueter operator} $\DD=\partial_a+\ii\partial_{x_1}+\jj\partial_{x_2}
+\kk\partial_{x_3}$, with the four-coordinate Laplacian $\Delta_4$
(\cref{squat:sec:fueter}). Here $\partial_a,\partial_{x_m},\partial_r,\partial_b$ are
\emph{variable} derivatives of ordinary coefficient functions and their Taylor
extensions; they are never $\dH$. The identity map is slice regular but has
$\DD q=-2$, so slice regularity does not imply Fueter regularity.
\end{enumerate}
\end{convention}

Reusing the three phase symbols of \cite{RepoDiff}: $\theta$ always denotes a
\emph{finite} angle (\cref{squat:prop:finiteangles}); $\eta$ always denotes an
infinitesimal phase residue, so that $\theta=\st(\theta)+\eta$; and $B$ denotes an
accumulated phase, an arbitrary surreal with $\partial B=b$, which need not be finite
and is not the angle of anything when it is infinite. No fourth sense of ``phase'' is
introduced, and $\Sin,\Cos$ of \cref{squat:eq:global-phase} are the global
integer-part-normalized functions, distinct from the finite-angle $\sin,\cos$ of
\cref{squat:eq:finite-trig}.

\begin{center}\small
\begin{tabular}{L{0.155\textwidth}L{0.215\textwidth}L{0.25\textwidth}L{0.26\textwidth}}
\toprule
Object & Domain & Laws it satisfies & Laws it does \emph{not} satisfy\\
\midrule
$\Expm$, $\Logm$ &
$\mm_H\leftrightarrow1+\mm_H$, by strong Hahn summation &
Bijection; conjugation- and inner-equivariant; $w(\Expm(x)-1)=w(x)$; additive on
commuting pairs; $\BCH$ otherwise &
Not additive on noncommuting pairs; not an additive-group isomorphism onto
$1+\mm_H$, the transported law
being noncommutative\\[3pt]
$\Expt$ &
the tube $\mathcal T$, excluding $\omega\ii$ &
Conjugation-equivariant; $\abs{\Expt(a+v)}=e^a$; addition law on commuting pairs;
surjective onto $\SQ^\times$; polar logarithms and $n$th roots &
Not a strong sum of $q^n/n!$; not $\exp(q_0)\Expm(\eps)$; no preferred logarithm on the
negative scalar axis\\[3pt]
$\Expr$ &
all of $\SQ$, onto $\SQ^\times$ &
Conjugation-equivariant and inverse-compatible, with
$\norm{\Expr q}=\exp_{\No}(\Rea q)$; agrees with $\Expm$ on $\mm_H$; on each slice a
surjective homomorphism with kernel $2\pi\Oz I$; multiplicative on commuting pairs &
Not an everywhere Hahn-summed power series; \emph{not} an
additive-to-multiplicative homomorphism on noncommuting arguments; fibre over $1$ is not
a kernel; no commuting chain rule for $\dH$\\[3pt]
$\dH$ &
all of $\SQ$, coefficientwise &
Leibniz; strongly additive; constant algebra $\HR$; surjective;
$\dH(\OO_H)\subseteq\mm_H$ &
No solution of $\dH y=cIy$ or $y cI$; hence no global commuting exponential chain rule
(\cref{squat:cor:nochain})\\[3pt]
$\delH$, $\ad_w$ &
$\HF$, any scalar derivation $\delta$ &
$D=\delH+\ad_w$ uniquely; $\delta N(q)=2\Rea(\delH(q)\bar q)$ &
$\delH(q^n)=nq^{n-1}\delH q$ only under commutation\\[3pt]
$\dsl$ &
admissible right-coefficient series &
Preserves the common-support class; coefficientwise Cauchy--Riemann identities &
Not the full differential: $d(q^2)_q\ne h\mapsto2hq$\\[3pt]
$df_q$ &
class scheme on $\SQ$ &
$\No$-linear; computes $d\Expr$ and $\det d\Expr$ &
No Banach structure; no inverse function theorem; zero differential does not imply
constancy (\cref{squat:ex:zeroderivative})\\[3pt]
$\DD$, $\Delta_4$ &
coefficient functions in four variables &
$\DD\Delta_4f=0$ for coherent slice-regular $f$ and for slice polynomials &
$\DD f=0$ does not follow from slice regularity: $\DD q=-2$\\
\bottomrule
\end{tabular}
\end{center}

\subsection{Why the homomorphism claim and its failure are both true}\label{squat:sub:consistency}
The physics companion \cite{RepoPhys} derives, for the \emph{commutative} surcomplex
field $\No[\ii]$, the identity
$\Exp_{\mathrm{EK}}(a+\ii b)=\exp_{\No}(a)(\Cos b+\ii\Sin b)$ and states that it extends
the ordinary complex exponential and \emph{has} the group-homomorphism property, with
kernel $2\pi\ii\Oz$. \Cref{squat:thm:globalexp} states that $\Expr$ is \emph{not} an
additive-to-multiplicative homomorphism on noncommuting arguments. These two statements
are consistent, and the reason is structural rather than a matter of emphasis.

Restrict $\Expr$ to one complex slice $\No+\No I$. That slice is a commutative field
isomorphic to $\No[\ii]$ (\cref{squat:thm:conjugacy}), the restriction \emph{is} the
surcomplex $\Exp_{\mathrm{EK}}$ by \cref{squat:thm:globalexp}, and on it the homomorphism
law holds with kernel exactly $2\pi\Oz I$. The physics report's claim is a claim about
that commutative object, and it is correct. The global failure is a failure at
\emph{noncommuting} arguments $p,q$: there is then no slice containing both, so
$\Expr(p+q)$ is computed from the modulus and direction of $\Ima(p+q)$, which bear no
multiplicative relation to those of $\Ima p$ and $\Ima q$.
\Cref{squat:ex:radialnotproduct} is the exact witness:
$\Expr(\omega\ii+\pi\jj)=1+\tfrac{\pi^2}{2\omega}\ii+\mathrm O(\omega^{-2})$ while
$\Expr(\omega\ii)\Expr(\pi\jj)=-1$. Both computations are exact, and the two arguments
$\omega\ii$ and $\pi\jj$ do not commute. So the homomorphism law is a theorem on each
slice and false across slices, and the additive group being abelian while the
multiplicative group is not already shows that a global homomorphism is not available.

Commutation is a sufficient condition for this law, not a necessary condition for an
individual pair: $p=6\pi\ii$, $q=8\pi\jj$ do not commute, but
$\Expr(p)=\Expr(q)=\Expr(p+q)=1$, since $\abs{p+q}=10\pi$.

A third statement must be kept out of this comparison. \Cref{squat:cor:nochain} says
that no map $E:\SQ\to\SQ^\times$ satisfies $\dH(E(q))=E(q)\dH q$ on commuting pairs
$q,\dH q$. That is a differential-algebraic statement about the specific derivation
$\dH$, not a statement about the multiplicative behaviour of any exponential. It is
compatible with $\Expr$ being a slicewise homomorphism, and it is compatible with
\cref{squat:thm:exp-derivative}, which computes the \emph{fine} differential $d\Expr$
--- a different operator again, in the sense of \cref{squat:conv:derivs}(c). All three
statements hold, of three different objects, and flattening any two of them into one
would manufacture a false contradiction across this collection.

\section{The division algebra and its norm}\label{squat:sec:algebra}

Let $F$ initially be any ordered field. Write $q=(a,u)=a+u$ with $a\in F$ and
$u\in F^3$ the pure imaginary part, identified with the span of $\ii,\jj,\kk$. The
multiplication law is
\begin{equation}\label{squat:eq:product}
 (a+u)(b+v)=ab-u\cdot v+av+bu+u\times v,
\end{equation}
where the dot and cross products are the polynomial formulas with the usual orientation
$(\ii,\jj,\kk)$. No metric completeness is hidden in this notation. Define
\[
 \bar q=a-u,\qquad T(q)=q+\bar q=2a,\qquad
 \Rea q=a,\qquad \Ima q=u,\qquad
 N(q)=q\bar q=a^2+u\cdot u .
\]
The notation $\Rea q$ means the \emph{central scalar part}; it need not be an ordinary
real number. Quaternion conjugation is distinct from a derivation, from the standard part
map of \cref{squat:thm:valuation}, and from conjugation inside a chosen complex slice.

\begin{proposition}[Associativity and change of scalars]\label{squat:prop:assoc}
$\HF$ is an associative unital $F$-algebra of dimension four. Every ordered field
embedding $F\hookrightarrow E$ induces an injective algebra homomorphism
$\HF\hookrightarrow\mathbb H_E$ preserving conjugation and norm square.
\end{proposition}
\begin{proof}
Impose Hamilton's relations in the free associative algebra and reduce every word to the
span of $1,\ii,\jj,\kk$; independence follows from the faithful matrix representation
\cref{squat:prop:matrix}. Alternatively \cref{squat:eq:product} verifies associativity
directly as a polynomial identity in twelve scalar variables, which is one of the
recorded symbolic checks. Applying the scalar embedding to each of four coordinates
proves the last assertion.
\end{proof}

\begin{theorem}[Hamilton algebra over an ordered field]\label{squat:thm:division}
The algebra $\HF$ is an associative division algebra with center exactly $F$.
Conjugation is an anti-involution, $\overline{pq}=\bar q\bar p$, and
\begin{equation}\label{squat:eq:inverse}
 q\bar q=\bar qq=N(q),\qquad N(pq)=N(p)N(q),\qquad
 q^{-1}=\frac{\bar q}{N(q)}\quad(q\ne0).
\end{equation}
Every element satisfies
\begin{equation}\label{squat:eq:quadratic}
 q^2-T(q)q+N(q)=0 .
\end{equation}
\end{theorem}
\begin{proof}
The conjugation identity follows from the relations on the basis, or from
\cref{squat:eq:product}, and gives $q\bar q=\bar qq=N(q)$. If $q\ne0$ then at least one
coordinate square is positive, and a sum of squares in an ordered field vanishes only
when every summand does, so $N(q)>0$; the inverse formula follows. Consequently
$N(pq)=pq\bar q\bar p=N(q)p\bar p=N(p)N(q)$. An element commuting with both $\ii$ and
$\jj$ has zero imaginary coordinates, so the center is exactly $F$. Expanding
$q(q+\bar q)$ proves \cref{squat:eq:quadratic}.
\end{proof}

In particular $\SQ$ has no zero divisors. Its ordered geometry comes from $F$, not from
an order on $\HF$ itself.

\begin{proposition}[Center, commutators, and the absence of a ring order]\label{squat:prop:center}
The center of $\HF$ is $F$, and
\[
 [a+u,b+v]=2\,u\times v,\qquad
 N([p,q])\le4N(\Ima p)N(\Ima q).
\]
There is no ordered-division-ring ordering on $\HF$ extending the order of $F$.
\end{proposition}
\begin{proof}
\Cref{squat:eq:product} gives the commutator, and a vector whose cross product with each
coordinate vector vanishes is zero, which is the assertion about the center. The identity
$N(u\times v)=N(u)N(v)-(u\cdot v)^2$ proves the squared-norm inequality, which
requires no square roots in $F$. In an ordered
division ring every nonzero square is positive, contradicting $\ii^2=-1$.
\end{proof}

Suppose from now on that $F$ is real closed, unless a weaker hypothesis is stated.
Define the \emph{$F$-valued modulus} and the scalar inner product by
\[
 \abs q=\sqrt{N(q)},\qquad (p,q)=\Rea(p\bar q).
\]
Taking positive square roots in \cref{squat:prop:center} now gives
$\abs{[p,q]}\le2\abs{\Ima p}\abs{\Ima q}$. The distinction matters already
over $F=\Q$: $N(1+\ii)=2$ belongs to $F$, whereas its positive square root does not.
This is not an ordinary real norm when $F$ is non-Archimedean, and an implementation must
not replace it by a machine floating-point norm.

\begin{proposition}[Euclidean identities without completeness]\label{squat:prop:norm}
For $p,q\in\HF$,
\[
 \abs{pq}=\abs p\abs q,\qquad \abs{p+q}\le\abs p+\abs q,\qquad
 \abs{(p,q)}\le\abs p\abs q,\qquad
 \max_m\abs{q_m}\le\abs q\le2\max_m\abs{q_m}.
\]
For nonzero $p,q$, equality in the triangle inequality holds exactly when $q=\lambda p$
for some $\lambda\in F_{>0}$. Moreover, for invertible $p,q$,
\begin{equation}\label{squat:eq:inverse-lipschitz}
 p^{-1}-q^{-1}=p^{-1}(q-p)q^{-1},\qquad
 \abs{p^{-1}-q^{-1}}=\frac{\abs{p-q}}{\abs p\abs q}.
\end{equation}
Finally $\HF\otimes_FF[\ii]\simeq M_2(F[\ii])$, a splitting that introduces zero divisors
only after extension.
\end{proposition}
\begin{proof}
Multiplicativity and the coordinate bounds follow from multiplicativity of $N$ and the
four coordinate squares. For Cauchy--Schwarz, the polynomial $N(p-sq)$ in the central
variable $s$ is nonnegative for every $s\in F$. The case $q=0$ is immediate; otherwise
substituting $s=(p,q)/N(q)$ gives
$(p,q)^2\le N(p)N(q)$. Therefore $N(p+q)\le(\abs p+\abs q)^2$, proving the triangle
inequality; equality forces $(p,q)=\abs p\abs q$ and equality in Cauchy--Schwarz, and the
same substitution gives $p=((p,q)/N(q))q$ with positive scalar coefficient. The inverse
identity is verified by multiplication, and its modulus is exactly multiplicative by the
first assertion. No supremum, compactness or completeness is used. The splitting is
\cref{squat:prop:matrix}.
\end{proof}

\begin{example}[Several scales in one number]\label{squat:ex:scales}
For $q=\omega+\omega^{-1}\ii+\jj+\omega^{-\omega}\kk$,
\[
 N(q)=\omega^2+1+\omega^{-2}+\omega^{-2\omega},\qquad
 q^{-1}=\frac{\omega-\omega^{-1}\ii-\jj-\omega^{-\omega}\kk}
 {\omega^2+1+\omega^{-2}+\omega^{-2\omega}} .
\]
The coordinates may have arbitrarily different orders of magnitude without compromising
associativity, division, or the norm identities.
\end{example}

\begin{remark}[Which quaternion symbol algebra this is]\label{squat:rem:symbol}
A quaternion symbol algebra $(a,b)_F$ over a real closed field is split unless both $a$
and $b$ are negative: a positive parameter is a square and produces a nontrivial
idempotent, whereas two negative parameters can be rescaled to $(-1,-1)_F$. Thus
Hamilton's algebra is the quaternion division-algebra case. This modest classification
suffices here; no classification of arbitrary class-sized division algebras is invoked.
\end{remark}

\section{Complex slices, conjugacy, rotations, and rational charts}\label{squat:sec:slices}

Define the sphere of imaginary units
\[
 \mathbb S(F)=\{I\in\HF:\Rea I=0,\ N(I)=1\},
\]
which by \cref{squat:eq:quadratic} is exactly the solution set of $I^2=-1$. For
$I\in\mathbb S(F)$ the \emph{slice} $F_I=F+FI$ is isomorphic to $F[\ii]$ and is therefore
algebraically closed when $F$ is real closed. Slices intersect in $F$ unless their units
differ by sign; their union is the whole algebra, but not a disjoint union, and a product
of elements of two different slices need not lie in either.

\begin{theorem}[Slice decomposition, centralizers, conjugacy]\label{squat:thm:conjugacy}
Let $F$ be real closed. Every noncentral $q$ has a unique expression
\begin{equation}\label{squat:eq:slice-decomp}
 q=a+rI,\qquad a=\Rea q,\quad r=\abs{\Ima q}>0,\quad
 I=\frac{\Ima q}{\abs{\Ima q}}\in\mathbb S(F),
\end{equation}
its centralizer is exactly $F_I$, and its minimal polynomial over $F$ is
$X^2-2aX+(a^2+r^2)$. Two quaternions are conjugate by an invertible quaternion if and
only if they have the same trace and norm; the conjugacy class of $a+rI$ is
\[
 [a,r]=a+r\,\mathbb S(F),
\]
and the conjugator may be taken of modulus one. Every commuting pair lies in a common
slice.
\end{theorem}
\begin{proof}
The commutator of $(a,u)$ and $(b,v)$ is $(0,2u\times v)$; for $u\ne0$ its vanishing is
equivalent to $v$ being an $F$-multiple of $u$, which proves the centralizer assertion and
the last sentence, any slice being available for two central elements. The quadratic
follows from $I^2=-1$, and no degree-one central polynomial annihilates a noncentral
element. Conjugation preserves scalar part and norm. Conversely let
$I,J\in\mathbb S(F)$. If $J\ne-I$ then $h=1-JI\ne0$ satisfies $hI=I+J=Jh$. If $J=-I$,
choose a unit imaginary $K$ perpendicular to $I$; then $KI=-IK=JK$. Normalize $h$ by its
positive modulus.
\end{proof}

A nonreal surquaternion thus has a proper-class sphere of conjugates even though its
minimal polynomial over the center has degree two. Polynomial degree bounds must count
conjugacy classes or multiplicities, never individual points on such a sphere.

\begin{proposition}[Explicit square roots]\label{squat:prop:sqrt}
For $q=a+v\ne0$ put $R=\abs q$. If $q$ is not a nonpositive central scalar, let
$s_0=\sqrt{(R+a)/2}>0$; then the two square roots of $q$ are
\begin{equation}\label{squat:eq:sqrt}
 \pm\left(s_0+\frac{v}{2s_0}\right).
\end{equation}
A negative central scalar $-\rho$, $\rho>0$, has the roots $\sqrt\rho\,I$ for every
$I\in\mathbb S(F)$. A positive central scalar has exactly its two central square roots,
and zero has only zero. No ``infinite angle'' is required to solve these equations.
\end{proposition}
\begin{proof}
If $v\ne0$ then $R>\abs a$. The square of \cref{squat:eq:sqrt} has imaginary part $v$
and scalar part $s_0^2-\abs v^2/(4s_0^2)=a$, using $\abs v^2=R^2-a^2$. Any square root
commutes with its square, so for noncentral $q$ it lies in $F_I$, where a quadratic has
only two roots. For central $q$, writing a prospective root $b+u$ gives imaginary part
$2bu$: if $u=0$ one obtains the central roots, and if $b=0$ its square is $-\abs u^2$,
giving exactly the negative-scalar sphere. Positivity treats the remaining cases.
\end{proof}

For example $2+\ii$ squares to $3+4\ii$, while $\omega^{-1}\ii+\omega^{-2}\jj$ has a
square-root scale determined by the positive square root of its norm.

\subsection{Matrix realization}\label{squat:sub:matrix}
Put $C_F=F[\ii]$ and write $q=z+w\jj$ with $z,w\in C_F$. Since $\jj z=\bar z\jj$,
\[
 (z+w\jj)(z'+w'\jj)=(zz'-w\bar w')+(zw'+w\bar z')\jj .
\]

\begin{proposition}[Matrix realization]\label{squat:prop:matrix}
The map
\begin{equation}\label{squat:eq:rho}
 \rho(z+w\jj)=\begin{pmatrix}z&w\\-\bar w&\bar z\end{pmatrix}
\end{equation}
is a faithful algebra homomorphism $\HF\to M_2(C_F)$ with
$\rho(\bar q)=\rho(q)^*$, $\det\rho(q)=N(q)$ and $\tr\rho(q)=T(q)$. The modulus-one
group is identified with $\SU(2,C_F)$, and scalar extension splits the algebra,
$\HF\otimes_FC_F\simeq M_2(C_F)$.
\end{proposition}
\begin{proof}
The displayed product law is precisely matrix multiplication in
\cref{squat:eq:rho}; the determinant, trace and star formulas follow directly. A unitary
$2\times2$ matrix of determinant one has the displayed form, by comparing its adjugate
with its conjugate transpose. The four images of the basis elements are linearly
independent over $C_F$ and span the four-dimensional matrix algebra.
\end{proof}

The image of $\HF$ is a special four-dimensional real form of the matrix algebra, not the
entire matrix algebra. Complex scalar extension splits the algebra and introduces zero
divisors; it does not exhibit zero divisors in $\HF$.

\subsection{The rotation double cover as an algebraic statement}\label{squat:sub:spin}
Let $V=\Ima\HF\simeq F^3$ and $S^3(F)=\{u\in\HF:N(u)=1\}=\Sp(1,F)$. Every $q\ne0$ has
the unique polar factorization $q=\abs qu$ with $u\in S^3(F)$, and because the positive
scalar factor is central,
\begin{equation}\label{squat:eq:polargroup}
 \HF^\times\simeq F_{>0}\times S^3(F)
\end{equation}
as groups. This is algebraic polar decomposition; it does not assert topological
compactness of $S^3(\No)$. For a quaternion $u=a+v$,
\begin{equation}\label{squat:eq:rodrigues}
 u\,x\,\bar u=(a^2-\abs v^2)x+2(v\cdot x)v+2a(v\times x)\qquad(x\in V),
\end{equation}
an exact polynomial identity at arbitrary surreal scales, requiring no unit hypothesis;
for $N(u)=1$ it equals $uxu^{-1}$. Writing $u=c+sI$ with $c^2+s^2=1$ turns
\cref{squat:eq:rodrigues} into the unnormalized Rodrigues form
$(c^2-s^2)x+2s^2(I\cdot x)I+2cs(I\times x)$.

\begin{theorem}[Rotations, spin, and inner automorphisms]\label{squat:thm:rotations}
Over a real closed field $F$, the map $u\mapsto R_u=(x\mapsto uxu^{-1})$ gives an exact
sequence of groups
\begin{equation}\label{squat:eq:spin3}
 1\longrightarrow\{1,-1\}\longrightarrow S^3(F)\longrightarrow\SO(3,F)\longrightarrow1,
\end{equation}
and
\[
 \Aut_F(\HF)\simeq\SO(3,F)\simeq\HF^\times/F^\times .
\]
In particular every center-fixing automorphism of the surquaternions is inner. Moreover
\begin{equation}\label{squat:eq:spin4}
 S^3(F)\times S^3(F)\longrightarrow\SO(4,F),\qquad (u,v):q\longmapsto uqv^{-1},
\end{equation}
is surjective with kernel $\{(1,1),(-1,-1)\}$.
\end{theorem}
\begin{proof}
Conjugation preserves $V$, its norm and its cross product, so its matrix lies in
$\SO(3,F)$; its kernel commutes with every imaginary quaternion, hence is central, and
central elements of modulus one are $\pm1$.

\emph{Surjectivity, first route.} Let $A\in\SO(3,F)$. Orthogonality and determinant one
give $\det(A-I)=-\det(A-I)$, so $A$ has a nonzero fixed vector; normalize it to $I\in V$.
On $I^\perp$ the transformation is a two-dimensional rotation with entries $c,s\in F$,
$c^2+s^2=1$, in an appropriately oriented orthonormal basis. If $c\ne-1$ take
$a=\sqrt{(1+c)/2}$ and $b=s/(2a)$, so $a^2+b^2=1$ and \cref{squat:eq:rodrigues} shows
that $a+bI$ induces $A$; if $c=-1$, use $I$ itself.

\emph{Surjectivity, second route.} The cubic characteristic polynomial of $A$ has a root
in $F$, and every orthogonal eigenvalue in $F$ is $1$ or $-1$. If only a $-1$ root is
found, its perpendicular plane has determinant $-1$ and hence eigenvalues $1,-1$, since
its characteristic polynomial has discriminant $\tr^2+4>0$. Either way a fixed vector
exists; finish as before. The two arguments are independent; the first is the shorter, the
second avoids the determinant identity and only uses that an odd-degree polynomial over a
real closed field has a root.

For the automorphism statement, an $F$-algebra automorphism preserves the coefficients of
the quadratic minimal polynomial of every noncentral element, hence scalar part and norm,
hence the imaginary dot and cross products by \cref{squat:eq:product}; its action on $V$
lies in $\SO(3,F)$. Conversely each such orthogonal transformation preserves
\cref{squat:eq:product} and extends by fixing $F$; surjectivity makes it inner.
Equivalently, an automorphism sends $\ii,\jj,\kk$ to an oriented orthonormal imaginary
triple, and every such triple preserves Hamilton multiplication.

For \cref{squat:eq:spin4}, take an orientation-preserving isometry $A$ of $\HF$ and put
$z=A(1)$. Left multiplication by $z^{-1}$ has determinant $1$ and makes the composite fix
$1$; its restriction to $V$ lies in $\SO(3,F)$, so equals conjugation by a unit $u$, and
$A(q)=zuqu^{-1}$. Here the determinant of left multiplication by $q$ is $N(q)^2$,
obtainable from its four-coordinate matrix. For the kernel, $uqv^{-1}=q$ for
every $q$ first gives $u=v$ at $q=1$, then forces $u$ to be central. A central
modulus-one scalar is $1$ or $-1$, so the kernel is exactly
$\{(1,1),(-1,-1)\}$.
\end{proof}

Equation \cref{squat:eq:spin3} is a two-to-one map on algebraic group points. It is
\emph{not}, without a separate topology and definition, a covering-space theorem for the
fine surreal topology of \cref{squat:sec:topology}. An arbitrary ring automorphism may
also act on the center: restriction to $F$ and the coefficientwise extension of a field
automorphism show that for set-sized $F$ the full automorphism group is a semidirect
product of the inner automorphisms with $\Aut(F)$. For $\No$ this is read for each
individual class automorphism, without forming a collection of all class functions.

Consequently arbitrary-scale spatial rotations need no trigonometric value at an infinite
argument: they are constructed by rational operations and positive square roots alone.
The source manuscripts give Cayley and stereographic presentations of the same rational
chart on the unit group. Both formulas are exact at every scale, including infinitesimal
ones.

\begin{proposition}[Cayley and stereographic presentations of one chart]\label{squat:prop:charts}
For $x\in V$ the \emph{Cayley chart}
\begin{equation}\label{squat:eq:cayley}
 \operatorname{Cay}(x)=(1+x)(1-x)^{-1}
\end{equation}
is defined because $N(1-x)=1+N(x)>0$, has modulus one, never equals $-1$, and is
inverted on $S^3(F)\setminus\{-1\}$ by $x=(u-1)(u+1)^{-1}$, which is imaginary.
Equivalently, its \emph{stereographic presentation}
\begin{equation}\label{squat:eq:stereographic}
 u=\frac{1-\abs x^2+2x}{1+\abs x^2},\qquad
 x=\frac{\Ima u}{1+\Rea u}\quad(u\ne-1),
\end{equation}
has positive denominators and exact modulus one.
\end{proposition}
\begin{proof}
Since $\bar x=-x$, conjugation reverses the two factors and gives
\[
 \overline{\operatorname{Cay}(x)}=(1+x)^{-1}(1-x)
 =\operatorname{Cay}(x)^{-1}.
\]
All displayed rational functions of the single element $x$ commute, and
$N(1+x)=N(1-x)=1+N(x)$, proving modulus one. The equation
$\operatorname{Cay}(x)=-1$ would give $1+x=-1+x$, which is impossible.
Conversely, for $N(u)=1$ and $u\ne-1$, let $x=(u-1)(u+1)^{-1}$.
Using $\bar u=u^{-1}$ gives
$\bar x=(u^{-1}+1)^{-1}(u^{-1}-1)=-x$; substitution recovers $u$.
Thus this inverse is imaginary and the chart is a bijection onto the indicated domain.

Finally $(1-x)^{-1}=(1+x)/(1+N(x))$ and $x^2=-N(x)$ give
\[
 \operatorname{Cay}(x)=\frac{(1+x)^2}{1+N(x)}
 =\frac{1-N(x)+2x}{1+N(x)},
\]
which identifies the two presentations. Independently checking the stereographic norm
amounts to $(1-N(x))^2+4N(x)=(1+N(x))^2$. Its inverse follows from
$\operatorname{Im}u=2x/(1+N(x))$ and $1+\operatorname{Re}u=2/(1+N(x))>0$.
Neither presentation asserts that $\Sp(1,F)$ is topologically a real three-sphere.
\end{proof}

\subsection{Standard part of the unit group}\label{squat:sub:unitreduction}
Every coordinate of a unit surquaternion has absolute value at most $1$, so every unit
quaternion is finite in the sense of \cref{squat:thm:valuation}. Standard part therefore
gives a split exact sequence
\begin{equation}\label{squat:eq:unitsreduction}
 1\longrightarrow U_1\longrightarrow S^3(\No)\xrightarrow{\ \st\ }S^3(\R)
 \longrightarrow1,\qquad U_1=S^3(\No)\cap(1+\mm_H),
\end{equation}
whose section is the inclusion of constant quaternions. The kernel $U_1$ describes
infinitesimal rotations and is nonabelian; \cref{squat:thm:principal-log} and
\cref{squat:thm:rotation-log} identify it, as a set with an explicitly transported group law,
with the infinitesimal imaginary quaternions. The corresponding statement over a
workspace $D_\Gamma$ reads $1\to U_1\to U\to\Sp(1,\R)\to1$ with
$U=\{u\in\OO_D:N(u)=1\}$.

\subsection{Quaternionic fractional linear transformations}\label{squat:sub:mobius}
A further exact geometry is the right projective line $\mathbb P^1(\HF)$. A matrix
$M=\left(\begin{smallmatrix}a&b\\c&d\end{smallmatrix}\right)\in\GL_2(\HF)$ acts by
\begin{equation}\label{squat:eq:mobius}
 f(q)=(aq+b)(cq+d)^{-1}
\end{equation}
where the denominator is nonzero, with the usual projective interpretation at its pole
and at infinity. Right lines are essential: the column $(q,1)^T$ is multiplied on the
left by $M$, then normalized on the right by $(cq+d)^{-1}$. For finite affine arguments
in its domain, direct subtraction gives
\begin{equation}\label{squat:eq:mobiusdiff}
 (f(q)-f(p))(cp+d)=(a-f(q)c)(q-p),
\end{equation}
so the algebraic differential is $df_q(h)=(a-f(q)c)\,h\,(cq+d)^{-1}$. Its first factor is
nonzero, since otherwise the row $(a,b)$ would equal $f(q)(c,d)$, contradicting
invertibility of $M$. Consequently
\[
 \abs{df_q(h)}=\frac{\abs{a-f(q)c}}{\abs{cq+d}}\,\abs h,
\]
so these transformations are conformal in this exact algebraic differential sense over
$F$. No real compactification, manifold completeness or contour integration is needed.
Formula \cref{squat:eq:mobiusdiff} also shows why moving a denominator to the wrong side
changes the transformation.

\section{One-sided polynomials: roots, spheres, and multiplicity}\label{squat:sec:polynomials}

A polynomial with quaternion coefficients needs an evaluation convention. We use a
\emph{central} indeterminate $X$ with coefficients on the right:
\[
 P(X)=\sum_{n=0}^dX^na_n\in\HF[X],\qquad P(q)=\sum_{n=0}^dq^na_n .
\]
The product in this ring is written $\star$ when it must be distinguished from pointwise
multiplication of the evaluated functions, which it is not. Centrality of $X$ does not
make evaluation multiplicative. These conventions agree with the polynomial part of slice
regularity \cite{GP}.

\begin{lemma}[Twisted product and left division]\label{squat:lem:twisted}
If $P(q)\ne0$ then
\begin{equation}\label{squat:eq:twisted}
 (P\star Q)(q)=P(q)\,Q\bigl(P(q)^{-1}qP(q)\bigr),
\end{equation}
while if $P(q)=0$ then $(P\star Q)(q)=0$. For any $q$ there is a unique polynomial $Q$
with
\begin{equation}\label{squat:eq:leftdivision}
 P=(X-q)\star Q+P(q),
\end{equation}
so $P(q)=0$ if and only if $P=(X-q)\star R$ for some $R$: a zero gives a \emph{left}
linear factor when coefficients are evaluated on the right.
\end{lemma}
\begin{proof}
Writing $Q=\sum_nX^nb_n$ and regrouping finite sums, the product evaluates to
$\sum_nq^nP(q)b_n$. If $P(q)\ne0$, use $q^nP(q)=P(q)(P(q)^{-1}qP(q))^n$; if $P(q)=0$
every term vanishes. For division, solve coefficients downward: $b_{d-1}=a_d$ and
$b_{n-1}=a_n+qb_n$, with final remainder $a_0+qb_0=\sum_nq^na_n$. This also proves
uniqueness.
\end{proof}

Let $P^c=\sum X^n\bar a_n$ and define the \emph{normal polynomial}, equivalently the
central symmetrization,
\begin{equation}\label{squat:eq:normalpoly}
 N_P=P\star P^c=P^c\star P\in F[X].
\end{equation}
Writing $P=P_0+P_1\ii+P_2\jj+P_3\kk$ with $P_m\in F[X]$, whose components commute, gives
the explicit central form $N_P=P_0^2+P_1^2+P_2^2+P_3^2$. For degree $d$ its degree is
$2d$ with leading coefficient $N(a_d)>0$, and $N_{P\star Q}=N_PN_Q$. At a central $x$,
$N_P(x)=N(P(x))$; at a noncentral quaternion one must not replace the polynomial normal
by that pointwise formula.

For a noncentral class $[a,r]$, $r>0$, set $\Delta_{a,r}(X)=X^2-2aX+a^2+r^2$, which is
irreducible over $F$ and vanishes at every point of that class. Divide centrally:
\begin{equation}\label{squat:eq:sphereremainder}
 P=\Delta_{a,r}Q+XA+B,\qquad A,B\in\HF .
\end{equation}

\begin{lemma}[Remainder on a sphere]\label{squat:lem:sphere-remainder}
On the sphere $[a,r]$ evaluation is $P(q)=qA+B$. Reducing $N_P$ modulo $\Delta_{a,r}$
gives the two central coefficients
\begin{equation}\label{squat:eq:remainder-norm}
 \bigl(2aN(A)+2\Rea(A\bar B)\bigr)X+N(B)-(a^2+r^2)N(A).
\end{equation}
Hence $P$ has a zero on $[a,r]$ if and only if $\Delta_{a,r}\mid N_P$, and then exactly
one of the following holds: $A=B=0$ and the entire sphere is a zero set; or $A\ne0$ and
the unique zero on that sphere is
\begin{equation}\label{squat:eq:sphere-root}
 q=-BA^{-1}=-B\bar A/N(A),\qquad \Rea q=a,\quad N(q)=a^2+r^2 .
\end{equation}
\end{lemma}
\begin{proof}
The characteristic quadratic vanishes on the sphere, which gives the evaluation. Two
independent routes complete the proof, and both were carried out in the source
manuscripts.

\emph{Route one, via the norm of the remainder.} Modulo $\Delta_{a,r}$ the normal
polynomial is congruent to $N(XA+B)=N(A)X^2+2\Rea(A\bar B)X+N(B)$. If
$\Delta_{a,r}\mid N_P$ and $A=0$, this constant remainder forces $N(B)=0$, hence $B=0$.
If $A\ne0$, the degree-two polynomial must equal $N(A)\Delta_{a,r}$, so
$\Rea(A\bar B)=-aN(A)$ and $N(B)=(a^2+r^2)N(A)$; for $q=-BA^{-1}$ these give
$\Rea q=a$ and $N(q)=a^2+r^2$, so $q\in[a,r]$ and it is the required root.

\emph{Route two, by conjugating the remainder.} Conjugate and multiply the remainder and
reduce modulo $\Delta_{a,r}$; the result is exactly \cref{squat:eq:remainder-norm}. If
$\Delta_{a,r}\mid N_P$ both of its central coefficients vanish; if $A=0$ positivity
forces $B=0$, and otherwise $q=-B\bar A/N(A)$ has precisely the scalar part and norm in
\cref{squat:eq:sphere-root}.

Conversely, if $qA+B=0$ at a point of the sphere then either $A=B=0$, or $B=-qA$, in which
case $N(B)=(a^2+r^2)N(A)$ and $\Rea(A\bar B)=-aN(A)$, so
\cref{squat:eq:remainder-norm} vanishes and $\Delta_{a,r}\mid N_P$.
\end{proof}

\begin{theorem}[Complete zero-class criterion]\label{squat:thm:polynomialzeros}
Let $F$ be real closed and $P$ nonzero of degree $d$. For each $a\in F$ and $r>0$, $P$
has a root on $[a,r]$ if and only if $\Delta_{a,r}\mid N_P$. \emph{When this
divisibility holds}, if $A=0$ in
\cref{squat:eq:sphereremainder} then $B=0$ and the whole sphere consists of roots, and if
$A\ne0$ there is exactly one root there, $q=-BA^{-1}$. The central roots of $P$ are
exactly the roots of $N_P$ in $F$. This is an effective algebraic algorithm relative to
exact real-closed-field operations: form $N_P$, factor it over $F$, compute the
two-coefficient remainders, and distinguish spherical from isolated roots.
\end{theorem}
\begin{proof}
\Cref{squat:lem:sphere-remainder} gives every assertion about a fixed sphere. For central
$x$ the last claim is $N_P(x)=N(P(x))$ together with positivity of the norm.
\end{proof}

\begin{theorem}[Constructive fundamental theorem and factorization]\label{squat:thm:fta}
Over any real closed field $F$, every positive-degree one-sided $P\in\HF[X]$ has a zero
in $\HF$ and factors as
\[
 P=(X-q_1)\star(X-q_2)\star\cdots\star(X-q_d)\,a_d,
\]
with the displayed order of factors significant. Its zero set meets at most $d$ conjugacy
classes, and in each noncentral class it is either a single point or the whole sphere. Its
evaluation map $\HF\to\HF$ is surjective.
\end{theorem}
\begin{proof}
The nonconstant central polynomial $N_P$ has a linear or an irreducible quadratic factor
over $F$. A linear factor $X-x$ gives $0=N_P(x)=N(P(x))$, hence $P(x)=0$. Every monic
irreducible quadratic over a real closed field is $X^2-2aX+a^2+r^2$ with $r>0$; apply
\cref{squat:thm:polynomialzeros}. Left division \cref{squat:eq:leftdivision} reduces the
degree, and induction gives the factorization. Every noncentral zero class contributes its
distinct quadratic factor to $N_P$, and every central zero $x$ makes all four $P_m(x)$
vanish, so $(X-x)^2\mid N_P$; distinct classes supply relatively prime factors of degree
two, and $\deg N_P=2d$, so at most $d$ classes occur. Applying the existence result to
$P-b$ for each target $b$ gives surjectivity.
\end{proof}

This is a real-closed-field form of classical quaternionic zero theory --- compare the
slice-algebra treatment in \cite{GP} --- proved with no appeal to topological degree on a
compact sphere, which would need a separate justification over $\No$. It supplies an
actual finite root-recovery procedure from the central norm polynomial. It does not
promise that factoring central polynomials over an arbitrary surreal presentation is
computationally decidable.

\begin{proposition}[Conservation of class multiplicity]\label{squat:prop:rootcount}
For a noncentral zero class $C=[a,r]$ put $m_C=\ord_{\Delta_{a,r}}N_P$, and for a central
root $\lambda$ put $m_\lambda=\tfrac12\ord_{X-\lambda}N_P$, which is an integer because
repeatedly dividing $P$ by the central factor $X-\lambda$ removes a factor
$(X-\lambda)^2$ from its normal. Then for $\deg P=d$,
\[
 \sum_{C\text{ noncentral}}m_C+\sum_{\lambda\text{ central}}m_\lambda=d .
\]
A spherical zero class has $m_C\ge2$, so if $s$ is the number of spherical classes and
$\ell$ the number of remaining individual zeros, including central ones, then
\begin{equation}\label{squat:eq:rootcount}
 2s+\ell\le d .
\end{equation}
\end{proposition}
\begin{proof}
Every irreducible factor of $N_P$ corresponds to a zero by
\cref{squat:thm:polynomialzeros}; the degree of $N_P$ is $2d$ and the indicated
multiplicities count half its factor degrees. If the entire sphere vanishes, the remainder
criterion gives $\Delta_{a,r}\mid P$, and normal multiplicativity then gives
$\Delta_{a,r}^2\mid N_P$. Every other zero class has multiplicity at least one, which
implies \cref{squat:eq:rootcount}.
\end{proof}

\begin{example}[Factor roots are not a list of evaluated roots]\label{squat:ex:poly}
Let $P=(X-\ii)\star(X-\jj)=X^2-X(\ii+\jj)+\kk$. Then
\[
 P(\ii)=0,\qquad P(\jj)=2\kk,\qquad N_P=(X^2+1)^2 .
\]
Modulo $X^2+1$ the remainder is $-X(\ii+\jj)+(\kk-1)$, whose unique root on the only
possible class is $\ii$; the whole zero set is the single point $\ii$, an isolated root of
class multiplicity two. The two written linear factors therefore do not supply two
distinct pointwise zeros. In contrast $X^2+1$ vanishes on all of $\mathbb S(F)$. This
example is a compact regression test that $\star$ is not pointwise multiplication.
\end{example}

\begin{warning}[Two ways a noncommutative polynomial equation can fail to have a root]
\label{squat:warn:interspersed}
\Cref{squat:thm:fta} cannot be restated for expressions with coefficients interspersed
among the occurrences of the variable, and it cannot be restated for arbitrary
noncommutative expressions. Two independent counterexamples:
\begin{enumerate}[label=\textup{(\arabic*)},leftmargin=2.4em]
\item $q+\ii q\ii+1=0$ has no solution, since for $q=a+b\ii+c\jj+d\kk$ its left side is
$1+2c\jj+2d\kk$;
\item $\ii q-q\ii=1$ has no solution, since its left side has zero scalar part.
\end{enumerate}
The surquaternions are a noncommutative division algebra, not an algebraically closed
field, and the one-sided, central-indeterminate convention is a mathematical hypothesis,
not merely notation. Central polynomials behave differently again: there a noncentral root
forces its entire conjugacy sphere to vanish.
\end{warning}

\section{Quaternion-valued Hahn series, valuation, and residue quaternions}\label{squat:sec:hahn}

\subsection{Normal form and multiplication}\label{squat:sub:normalform}
Writing surreal normal forms with $t^\gamma=\omega^{-\gamma}$ turns the
reverse-well-order convention into a well-order convention. Combining the four scalar
Conway normal forms of a surquaternion and collecting coefficients gives the following.

\begin{theorem}[Normal form, multiplication, localization]\label{squat:thm:normal}
Every surquaternion has a unique expression
\begin{equation}\label{squat:eq:qnormal}
 q=\sum_{\gamma\in S}t^\gamma h_\gamma,\qquad h_\gamma\in\HR\setminus\{0\},
\end{equation}
with $S\subset\No$ a set-sized well-ordered support. Addition is coefficientwise and
multiplication is \emph{ordered} quaternion convolution,
\begin{equation}\label{squat:eq:convolution}
 (pq)_\eta=\sum_{\alpha+\beta=\eta}p_\alpha q_\beta ,
\end{equation}
each displayed coefficient sum being finite. Every set of surquaternions lies in one
workspace $D_\Gamma$ with $\Gamma$ set-sized and divisible
(\cref{squat:thm:localization}), and the coefficientwise identification
\begin{equation}\label{squat:eq:hahnbridge}
 D_\Gamma=\mathbb H_{\R((t^\Gamma))}\simeq\HR((t^\Gamma))
\end{equation}
is an isomorphism of algebras with involution. The monomials $t^\gamma$ are
\emph{central}.
\end{theorem}
\begin{proof}
A finite union of well-ordered subsets of a linearly ordered set is well-ordered, and
uniqueness follows in each scalar coordinate. The ordinary Hahn sumset lemma makes the
convolution locally finite, and centrality of the monomials leaves the product of
coefficients in its original order; finite coefficient rearrangement respecting that
order proves associativity and distributivity. Localization is
\cref{squat:thm:localization}, and real closedness of $K_\Gamma$ for divisible $\Gamma$ is
the cited Hahn-field theorem \cite{Alling,EK}. The four-coordinate construction proves the
final isomorphism, which respects conjugation because the structure constants are
integers.
\end{proof}

Consequently \cref{squat:eq:hahnbridge} is \emph{a Hahn series algebra with
noncommutative coefficients, not a series construction with noncommuting exponents}. The
order of $h_\alpha$ and $k_\beta$ in \cref{squat:eq:convolution} is essential; the value
group is never noncentral, and all the noncommutativity sits in the $\HR$-valued
coefficients.

\begin{definition}[Hahn-summable family]\label{squat:def:summable}
A \emph{set-indexed} family $(q_j)_{j\in J}$ is \emph{Hahn-summable}, equivalently
\emph{strongly summable}, if the union of its supports is well-ordered and each exponent
belongs to the supports of only finitely many $q_j$. Its sum is defined coefficientwise,
by adding those finitely many quaternion coefficients at each exponent. Every infinite
sum below means this, unless an ordinary real or complex integral is explicitly named
\cite{BM}. Ordinary convergence of infinitely many real coefficients is not a substitute,
and this definition does not require convergence of partial sums in a topology. Such
convergence may hold in a specified workspace, but is a separate assertion.
\end{definition}

\begin{theorem}[Exact valuation, leading coefficients, and residue]\label{squat:thm:valuation}
For $q\ne0$ put $w(q)=\min\supp(q)$ and $\lc(q)=h_{w(q)}\in\HR^\times$, with
$w(0)=+\infty$. Then for nonzero $p,q$,
\begin{align}
 w(pq)&=w(p)+w(q), &\lc(pq)&=\lc(p)\lc(q),\label{squat:eq:valuation-product}\\
 w(p+q)&\ge\min\{w(p),w(q)\}, &w(\bar q)&=w(q),\\
 w(N(q))&=2w(q), &w(\abs q)&=w(q),\label{squat:eq:normval}
\end{align}
and in coordinates $w(a+b\ii+c\jj+d\kk)=\min\{w(a),w(b),w(c),w(d)\}$. The valuation ring
and its maximal \emph{two-sided} ideal are
\[
 \OO_H=\{q:w(q)\ge0\},\qquad \mm_H=\{q:w(q)>0\},
\]
consisting respectively of the quaternions with finite coordinates and with infinitesimal
coordinates, where ``finite'' means bounded by some positive ordinary real. Coefficient
extraction at exponent zero induces
\begin{equation}\label{squat:eq:residue}
 \st:\OO_H\longrightarrow\HR,\qquad \ker\st=\mm_H,\qquad
 \OO_H/\mm_H\simeq\HR,
\end{equation}
a surjective homomorphism of algebras with involution onto the \emph{noncommutative}
ordinary quaternion division algebra. Every finite $q$ decomposes uniquely as
$q=\st(q)+\eps$ with $\eps\in\mm_H$.
\end{theorem}
\begin{proof}
At the smallest possible exponent in a product the coefficient is the product of the two
leading coefficients, which is nonzero because $\HR$ is a division algebra; this proves
\cref{squat:eq:valuation-product}, and the ultrametric inequality is coefficientwise. The
leading coefficient of $N(q)=q\bar q$ is $N(\lc(q))>0$, giving the first equality in
\cref{squat:eq:normval}, and taking its positive square root gives the second. A negative
leading exponent is infinitely large, a positive one infinitesimal, and exponent zero
carries a real quaternion coefficient followed by infinitesimal terms; on nonnegative
supports only two zero exponents can contribute to the zero coefficient of a product, so
$\st$ is a ring homomorphism with the indicated kernel and image. It commutes with
conjugation because conjugation is coefficientwise. The coordinate formula is immediate
from supports.
\end{proof}

Thus $\OO_H^\times$ consists exactly of the elements with nonzero standard part. The
residue algebra is $\HR$, not $\C$ and not a commutative field: leading-coefficient
noncancellation in \cref{squat:eq:valuation-product} is a division-algebra argument, and
the identity $\lc(pq)=\lc(p)\lc(q)$ is \emph{order-sensitive}. Valuation is invariant
under inner conjugation, while its leading quaternion is conjugated by the leading
quaternion of the conjugating element. The associated graded algebra in a workspace is the
group-graded algebra with one copy of $\HR$ at each $\gamma\in\Gamma$, written
$\HR[\Gamma]$ with central monomials.

\begin{warning}[Valuation is not the modulus]\label{squat:warn:valnotnorm}
The valuation is an additive order-of-magnitude measure, not the modulus: $1$ and
$2+\ii$ have the same valuation and different moduli. For matrices, only the inequality
$w(AB)\ge w(A)+w(B)$ holds; the equality valid for individual nonzero quaternions must
not be copied (see \cref{squat:sub:matrixnorm}).
\end{warning}

\subsection{Multiplicative normal form, the exact geometric inverse, and its certificate}
\label{squat:sub:geominverse}

\begin{proposition}[Exact geometric inverse and error certificate]\label{squat:prop:neumann}
Every nonzero $q$ has a unique decomposition
\begin{equation}\label{squat:eq:leadingdecomp}
 q=t^\gamma h(1+\eps),\qquad
 \gamma=w(q),\ h=\lc(q)\in\HR^\times,\ \eps\in\mm_H,
\end{equation}
and its inverse is the Hahn sum
\begin{equation}\label{squat:eq:neumann}
 q^{-1}=\left(\sum_{n\ge0}(-\eps)^n\right)h^{-1}t^{-\gamma}.
\end{equation}
For nonzero $\eps\in\mm_H$ and every integer $M\ge0$,
\begin{equation}\label{squat:eq:neumannerror}
 w\left((1+\eps)^{-1}-\sum_{n=0}^M(-\eps)^n\right)=(M+1)\,w(\eps).
\end{equation}
The finite unit group $\OO_H^\times$ is a split extension of $\HR^\times$ by $1+\mm_H$,
with the conjugation action displayed by
$h(1+\eps)\cdot k(1+\delta)=hk(1+k^{-1}\eps k)(1+\delta)$.
\end{proposition}
\begin{proof}
Set $\eps=h^{-1}t^{-\gamma}q-1$; the leading term has been removed, so all remaining
exponents are positive. \Cref{squat:lem:neumann} makes the geometric family
Hahn-summable, and multiplication by $1+\eps$ telescopes coefficientwise to $1$ on either
side. The finite remainder is exactly $(-\eps)^{M+1}(1+\eps)^{-1}$, whose final factor has
valuation zero, so the valuation identities give \cref{squat:eq:neumannerror}. Reduction
modulo $\mm_H$ supplies the split group map, and the displayed product exhibits the
action.
\end{proof}

The order in \cref{squat:eq:neumann} is not cosmetic: in general $h$ and $\eps$ do not
commute. The multiplicative structure is a central monomial factor together with a
semidirect product of principal units and ordinary nonzero quaternions, and the principal
units are already nonabelian (\cref{squat:sec:explog}).

\begin{example}[Two scales in one inverse]\label{squat:ex:twoscaleinverse}
Let $t=\omega^{-1}$ and $q=1+t\ii+t^2\jj$. Then
\[
 q^{-1}=\frac{1-t\ii-t^2\jj}{1+t^2+t^4},
\]
and its formal geometric expansion agrees exactly with \cref{squat:eq:neumann}. The terms
of the inverse are not obtained by independently inverting the four coordinate series.
The absence of a $\kk$ term here is the identity $\ii\jj+\jj\ii=0$, not an approximation.
\end{example}

\section{Summability certificates and the topology that must not be confused with them}
\label{squat:sec:topology}

\subsection{The all-lengths positive-support lemma}\label{squat:sub:neumann}

\begin{lemma}[Positive-support (Neumann) lemma, all lengths]\label{squat:lem:neumann}
Let $S$ be a well-ordered subset of the positive cone of an ordered abelian group. Then
the generated additive monoid
\[
 S^*=\{s_1+\cdots+s_n:n\ge0,\ s_m\in S\}
\]
is well-ordered, and for each fixed exponent $\gamma$ only finitely many finite ordered
words $(s_1,\ldots,s_n)$ over $S$ have sum $\gamma$, counting \emph{all lengths at once}.
Consequently every noncommutative formal power series in finitely many variables with
ordinary quaternion coefficients can be evaluated at finitely many elements of $\mm_H$,
retaining the order of all factors.
\end{lemma}
\begin{proof}
A well-ordered alphabet is well-quasi-ordered. Higman's finite-word lemma says that its
finite words are well-quasi-ordered under subword embedding with entrywise domination.
Such an embedding can only increase the sum of a word, and increases it strictly unless no
entries are inserted and no entries increase, because all entries are positive. A strictly
decreasing sequence of word sums would contradict the well-quasi-order property, so the
sums are well-ordered; and distinct words with one fixed sum form an antichain, so there
are only finitely many of them. (The usual choice principle is used to characterize a
non-well-order by a decreasing sequence.)

For evaluation, take the finite union of the supports of the arguments. A contribution at
exponent $\gamma$ is specified by a word over this union, a choice of one of the finitely
many variables at each position, and its coefficient term; there are only finitely many
such contributions, and the possible exponents lie in $S^*$.
\end{proof}

This all-lengths form, rather than a bound for each fixed word length, is the
combinatorial input behind every infinitesimal power-series evaluation below; its use in
surreal summation is discussed in \cite{BM}, and the repository's
\texttt{Surreal/HahnSeries/Neumann.lean} proves the same all-lengths statement
\cite{Repo}. It is important that the proof works for arbitrary ordered abelian groups,
not only rank-one or Archimedean ones.

\begin{warning}[Positive valuation is a summability certificate, not cofinality]
\label{squat:warn:cofinality}
If $w(\eps)=\delta>0$ then $w(\eps^n)=n\delta$, but the exponents $n\delta$ need not be
cofinal in the value group: in the lexicographically ordered group $\Z\times\Z$ the
sequence $n(0,1)$ never exceeds $(1,0)$. The geometric series is Hahn-summable all the
same. Thus a remainder estimate $w(R_n)\ge n\delta$ is \emph{not} automatically the
assertion that $R_n\to0$ in the valuation topology of the workspace, and a valid Hahn
identity need not be the limit of its partial sums. Hahn summation avoids the ambiguity by
specifying its coefficients directly.
\end{warning}

\subsection{The full surreal fine topology}\label{squat:sub:fine}
The \emph{fine topology} on $\SQ$ has basic balls $B(q,r)=\{p:\abs{p-q}<r\}$ for
\emph{every} positive surreal $r$, on the full surreal class. This is a
ball-neighbourhood scheme; no collection whose elements are all open classes is required.

\begin{theorem}[Set discreteness and the failure of sequential completion]\label{squat:thm:discrete}
Every set $A\subset\SQ$ is uniformly separated and closed in the fine topology. Every
set-indexed fine-Cauchy net is eventually constant, and every convergent set-indexed net
is eventually equal to its limit. Every continuous map from an ordinary connected real
interval into $\SQ$ is constant.
\end{theorem}
\begin{proof}
The positive distances between distinct points of $A$ form a set of positive surreals.
Every set of positive surreals has a strictly positive lower bound --- take a surreal
larger than all their reciprocals and invert it; equivalently, the surreal cut property
supplies a positive $r$ with $2r$ smaller than all of them. Thus every ball of radius $r$
contains at most one point of $A$. For $q\notin A$ the same argument applied to
$\{\abs{q-a}:a\in A\}$ supplies a ball at $q$ disjoint from $A$. This proves closedness
and discreteness, the empty and singleton cases being trivial.

A set-indexed net has a set as its image, and the Cauchy condition applied to a uniform
separation radius forces all sufficiently late terms to coincide; if a net converges,
apply the separation argument to its image together with the proposed limit. Finally a
continuous interval map has a set-sized image, discrete in the induced topology, and a
connected subset of a discrete space has at most one point.
\end{proof}

This theorem is about the \emph{full} surreal class and its fine topology. It is false as
a blanket assertion about a set-sized Hahn field with its own \emph{intrinsic} valuation
topology: in $\R((t^\Q))$ one has $t^n\to0$, because the integers are cofinal in $\Q$, and
the geometric partial sums converge there to $(1-t)^{-1}$. Inside $\No$, however, every
$t^n$ exceeds $t^\omega$, and neither assertion of convergence survives the finer
neighbourhood test; the subspace topology that the entire workspace inherits from $\No$ is
discrete, whereas its intrinsic topology is not.

\begin{remark}[A workspace carries three different topologies]\label{squat:rem:threetopologies}
For $\Gamma\ne0$, on $K_\Gamma$ or $D_\Gamma$ one must distinguish the
\emph{intrinsic} valuation topology, equivalently the topology of scalar modulus balls
using radii in that workspace (the order topology on $K_\Gamma$); the discrete subspace
topology inherited from the full fine topology; and the \emph{standard-part} topology on
the finite part, pulled back from $\R$ or $\HR\simeq\R^4$ by
\cref{squat:eq:residue}. The latter is not Hausdorff: $0$ and a nonzero infinitesimal
have the same standard part. For the stated equivalence of intrinsic topologies, choose
$\delta>0$ in $\Gamma$. Monomial radii bound valuation balls, and, for any workspace
radius $r>0$, a valuation ball with threshold greater than $w(r)$ lies inside the
modulus ball of radius $r$; conversely a modulus ball of radius $t^{\gamma+\delta}$
lies inside $\{x:w(x)>\gamma\}$.

The degenerate group $\Gamma=0$ is an exception: $K_0=\R$ and $D_0=\HR$, the trivial
valuation topology and the inherited full fine topology are discrete, whereas the
modulus and standard-part topologies are the usual Hausdorff real or Euclidean ones.
In the full fine topology no sequence of distinct
rational approximations converges finely to $\sqrt2$, and the partial sums of
$\sum\omega^{-n}$ do not converge finely to their Hahn sum. Coefficientwise analysis over
ordinary domains uses its ordinary domain as a geometric base and does not pretend that its
paths are fine-continuous surreal paths.
\end{remark}

\begin{example}[Zero fine differential does not imply constancy]\label{squat:ex:zeroderivative}
On $\SQ$ let $f(q)=0$ when all coordinates of $q$ are finite and $f(q)=1$ otherwise. The
finite-coordinate part is open, since a ball of radius $1$ around one of its points stays
finite; its complement is open, since at an infinite $q$ the ball of radius $\abs q/2$
stays infinite by the reverse triangle inequality. Hence $f$ is locally constant and has
zero fine differential at every point, but it is not globally constant.
\end{example}

The example is not a paradox: a mean-value argument would require a connectedness or
integration theorem that this topology does not provide. For the rest of the article an
infinite expression is marked by its construction --- Hahn summation, ordinary real or
complex analysis on coefficient functions, or a specified surreal operation --- and the
word ``convergent'' is never used as a substitute for one of those choices.

\section{Infinitesimal Lie theory: exponential, logarithm, and rotation groups}\label{squat:sec:explog}

\subsection{The canonical neighbourhood of the identity}\label{squat:sub:explog}
For $x,y\in\mm_H$ define
\begin{equation}\label{squat:eq:explog}
 \Expm(x)=\sum_{n\ge0}\frac{x^n}{n!},\qquad
 \Logm(1+y)=\sum_{n\ge1}\frac{(-1)^{n+1}y^n}{n},
\end{equation}
both Hahn-summable by \cref{squat:lem:neumann}. Each individual series lies in the
commutative subalgebra generated by its one argument together with the central scalars;
the constant coefficient is handled as a single term, not as an ordinary convergent sum of
constants, and the subscript emphasizes Hahn evaluation rather than a limit of partial
sums.

\begin{theorem}[Local exponential--logarithm equivalence]\label{squat:thm:principal-log}
The maps in \cref{squat:eq:explog} are mutually inverse bijections
\[
 \Expm:\mm_H\longleftrightarrow1+\mm_H:\Logm .
\]
They commute with quaternion conjugation and with inner automorphisms, and for $x\ne0$
\[
 w(\Expm(x)-1)=w(x),\qquad \lc(\Expm(x)-1)=\lc(x).
\]
If $x$ and $y$ commute then $\Expm(x+y)=\Expm(x)\Expm(y)$. Their restrictions give a
bijection
\begin{equation}\label{squat:eq:logunits}
 \Expm:\mm_H\cap\Ima\SQ\ \longleftrightarrow\ U_1:\Logm ,
\end{equation}
with $U_1$ as in \cref{squat:eq:unitsreduction}.
\end{theorem}
\begin{proof}
The one-variable formal identities $\log(\exp T)=T$ and $\exp(\log(1+T))=1+T$ may be
substituted at positive valuation: at each exponent only finitely many words contribute by
\cref{squat:lem:neumann}, so formal regrouping becomes valid Hahn regrouping. Each
identity involves powers of a single element, which generate a commutative subalgebra, so
noncommutativity elsewhere is irrelevant. The leading coefficient of $\Expm(x)-1$ comes
only from the linear term. Coefficientwise conjugation preserves these identities, and
inner automorphisms preserve products and valuation; multiplying a Hahn-summable family by
fixed factors $h,h^{-1}$ preserves summability by the sumset lemma, although individual
supports change.

If $x$ is imaginary then $\overline{\Expm(x)}=\Expm(-x)$, so the modulus is $1$.
Conversely if $u\in U_1$ then $\bar u=u^{-1}$, hence
$\overline{\Logm u}=\Logm(u^{-1})=-\Logm u$ and the logarithm is imaginary. The commuting
addition law follows from the formal binomial expansion and the same support regrouping.
\end{proof}

This is a bijection, \emph{not} an assertion that the additive group of imaginary
infinitesimals is the group of infinitesimal rotations: the transported group law is
noncommutative, and the correct law is $\BCH$.

\begin{proposition}[Baker--Campbell--Hausdorff at positive valuation]\label{squat:prop:bch}
For $x,y\in\mm_H$ the expression
$\BCH(x,y)=\Logm\bigl(\Expm(x)\Expm(y)\bigr)$ is defined by Hahn summation, and
\begin{align}\label{squat:eq:bch}
 \BCH(x,y)={}&x+y+\tfrac12[x,y]
 +\tfrac1{12}[x,[x,y]]+\tfrac1{12}[y,[y,x]]\notag\\
 &-\tfrac1{24}[y,[x,[x,y]]]+R_5 ,
\end{align}
where ``degree'' counts occurrences of $x$ and $y$, not powers of a single numerical
scale. If $w(x),w(y)\ge\delta>0$ then the degree-four-and-higher remainder has valuation
at least $4\delta$ and $w(R_5)\ge5\delta$. Writing $\alpha=w(x)$, $\beta=w(y)$ for nonzero
$x,y$, the group commutator satisfies
\begin{equation}\label{squat:eq:groupcomm}
 \Expm x\,\Expm y\,\Expm(-x)\,\Expm(-y)=1+[x,y]+R,
 \qquad w(R)\ge\alpha+\beta+\min\{\alpha,\beta\} .
\end{equation}
Each exponential here is the infinitesimal one of \cref{squat:conv:exps}(a).
\end{proposition}
\begin{proof}
Use the formal exponential and logarithm in the free associative power-series algebra on
two variables, filtered by total word degree. Their compositions have finitely many words
at each degree, and expansion through degree four gives \cref{squat:eq:bch}. At arguments
of positive valuation every word is controlled by the monoid generated by their combined
positive supports, and \cref{squat:lem:neumann} makes evaluation and regrouping legitimate
across all degrees; a common support certificate remains necessary for the whole infinite
expression. Terms of degree at least $n$ have valuation at least $n\delta$. In the group
commutator the degree-two term is $xy-yx$, and every remaining nonzero word involves both
variables and at least one additional occurrence, giving the stated bound. This derivation,
rather than any analytic convergence theorem, is the one used.
\end{proof}

For imaginary $x,y$ one has $[x,y]=2x\times y$. If their leading imaginary coefficients
are not parallel, the commutator has valuation exactly $\alpha+\beta$ and is nonzero.

\begin{example}[Rotations at unequal scales]\label{squat:ex:unequalscales}
For $x=t^\alpha\ii$ and $y=t^\beta\jj$ with $\alpha,\beta>0$,
\[
 \BCH(x,y)=t^\alpha\ii+t^\beta\jj+t^{\alpha+\beta}\kk
 -\tfrac13t^{2\alpha+\beta}\jj-\tfrac13t^{\alpha+2\beta}\ii+\cdots,
\]
and the leading failure of commutativity appears at the sum of the two scales; there is no
requirement that $\alpha$ and $\beta$ be rationally comparable. For $\alpha=1$, $\beta=2$,
\[
 \Expm x\,\Expm y\,\Expm(-x)\,\Expm(-y)=1+2t^3\kk+
 \text{terms of valuation at least }4 .
\]
For a vector $v$ the leading infinitesimal action is
$\Expm(x)v\Expm(-x)=v+2t^\alpha(\ii\times v)+\cdots$. Thus two infinitesimal rotations can
agree with the identity to distinct orders and still produce a third, smaller rotational
scale when commuted. The associated graded bracket is precisely twice the ordinary cross
product, accompanied by addition of exponents. These bounds do not assume $n\delta$
cofinal and do not turn the Hahn identities into limits of a sequence in the fine topology
(\cref{squat:warn:cofinality}).
\end{example}

\begin{proposition}[Adjoint series at infinite arguments]\label{squat:prop:adjoint}
For $x\in\mm_H$ and any fixed $q\in D_\Gamma$,
\begin{equation}\label{squat:eq:adjoint}
 \Expm(x)\,q\,\Expm(-x)=\sum_{n\ge0}\frac{\ad_x^n(q)}{n!},
 \qquad \ad_x(q)=[x,q],
\end{equation}
and the series is Hahn-summable \emph{including when $q$ is infinite}.
\end{proposition}
\begin{proof}
The formal identity follows by collecting terms in $\sum_{r,s}x^rq(-x)^s/(r!s!)$. Its
supports lie in $\supp(q)+S^*$ with $S=\supp(x)$, and the Hahn sumset lemma with
\cref{squat:lem:neumann} gives local finiteness. Noncommutativity is retained in the words
$x^rqx^s$.
\end{proof}

\subsection{The filtered group of infinitesimal rotations}\label{squat:sub:rotationgroup}

\begin{theorem}[Logarithms of infinitesimal rotations]\label{squat:thm:rotation-log}
The maps $\Expm$ and $\Logm$ restrict to mutually inverse bijections between the
imaginary infinitesimals $\Ima(\mm_H)$ and $U_1$. On the left the group law is $\BCH$, not
addition. Every modulus-one $u$ factors uniquely as $u=\st(u)\,\omega_1$ with
$\omega_1\in U_1$, once this multiplication order is fixed; equivalently the sequence
\cref{squat:eq:unitsreduction} splits, and every unit quaternion is finite, so this
describes the whole modulus-one group.
\end{theorem}
\begin{proof}
The bijection is \cref{squat:eq:logunits}. The factorization is $\omega_1=\st(u)^{-1}u$,
which has standard part one and modulus one, and it is unique because $\st$ is a
homomorphism on $\OO_H^\times$. Every coordinate of a unit has absolute value at most
$1$, so units are finite.
\end{proof}

\begin{theorem}[Associated graded rotation algebra]\label{squat:thm:graded}
For $\gamma\in\Gamma_{>0}$ put $U_{\ge\gamma}=\{u\in U_1:w(u-1)\ge\gamma\}$ and define
$U_{>\gamma}$ similarly. Leading coefficients give
\[
 U_{\ge\gamma}/U_{>\gamma}\simeq(\Ima\HR,+).
\]
If $u=\Expm(x)$ and $v=\Expm(y)$ with $x,y$ imaginary infinitesimals of valuations
$\alpha,\beta$ and leading coefficients $X,Y$, then
\[
 uvu^{-1}v^{-1}=1+t^{\alpha+\beta}[X,Y]+\text{terms of valuation greater than }
 \alpha+\beta ,
\]
so its valuation above $1$ is exactly $\alpha+\beta$ when $X\times Y\ne0$, and the induced
graded bracket is $[X,Y]=2X\times Y$.
\end{theorem}
\begin{proof}
The product of $1+a$ and $1+b$ has leading correction $a+b$, since $ab$ has larger
valuation; the modulus-one condition forces the leading correction to be imaginary, and
every such coefficient is realized by $\Expm(t^\gamma X)$. Expanding the group commutator,
its first nonzero formal term is $[x,y]$ by \cref{squat:eq:groupcomm}, and each subsequent
term contains both arguments and at least one additional positive-valuation factor. This
includes the possible cancellation when $[X,Y]=0$.
\end{proof}

\section{Finite angles and polar coordinates at arbitrary scale}\label{squat:sec:angles}

Let $\OO_{\No}$ and $\mm_{\No}$ denote the finite and the infinitesimal surreal scalars.
For a finite surreal $\theta=r+\eta$ with $r\in\R$ and $\eta\in\mm_{\No}$ define
\begin{align}\label{squat:eq:finite-trig}
 \cos\theta&=\cos r\sum_{n\ge0}\frac{(-1)^n\eta^{2n}}{(2n)!}
            -\sin r\sum_{n\ge0}\frac{(-1)^n\eta^{2n+1}}{(2n+1)!},\\
 \sin\theta&=\sin r\sum_{n\ge0}\frac{(-1)^n\eta^{2n}}{(2n)!}
            +\cos r\sum_{n\ge0}\frac{(-1)^n\eta^{2n+1}}{(2n+1)!}.\notag
\end{align}
The ordinary real values are taken \emph{first}, and only the infinitesimal correction is
Hahn-summed; the real Taylor sums are not misrepresented as Hahn-summable families of
constants. This is a restricted analytic extension, consistent with the established
surreal analytic structure \cite{vdDE,EK} and with the repository's trigonometry report
\cite{RepoTrig}, whose scalar results we cite rather than reprove. Formal addition formulas
and the ordinary real addition formulas prove their finite-surreal analogues. In the
vocabulary of \cref{squat:sec:ledger}, $\theta$ is a \emph{finite angle} and $\eta$ its
\emph{infinitesimal phase residue}.

\begin{proposition}[Every direction has a finite angle; the finite phase theorem]
\label{squat:prop:finiteangles}
Fix $I$ with $I^2=-1$. The map $\theta\mapsto\cos\theta+I\sin\theta$ is a surjective
homomorphism from $(\OO_{\No},+)$ onto the unit circle of $\No+\No I$ with kernel
$2\pi\Z$. Equivalently: for $c,s\in\No$ with $c^2+s^2=1$ there is a \emph{finite} surreal
$\theta$, unique modulo $2\pi\Z$, with $c+Is=\cos\theta+I\sin\theta$, and if $s\ge0$ there
is a unique representative $\theta\in[0,\pi]$.
\end{proposition}
\begin{proof}
The addition and norm identities are formal Taylor identities with a common positive
support. The coordinates $c,s$ are finite and their standard parts lie on the ordinary unit
circle; choose an ordinary angle $r$ of that standard part. Multiplication by
$\cos r-I\sin r$ leaves a unit in $1+\mm_H$ lying in the commutative slice $\No+\No I$, and
its principal logarithm is $I\eta$ for a unique infinitesimal $\eta$ by
\cref{squat:thm:principal-log}; then $\theta=r+\eta$ works. If a phase equals $1$, reducing
to its standard part gives $r\in2\pi\Z$, and injectivity of $\Logm$ forces $\eta=0$; this
is both the kernel computation and uniqueness modulo $2\pi\Z$. On the upper semicircle
choose the ordinary representative in $[0,\pi]$; at its endpoints the sign of the
infinitesimal sine correction selects the inward direction, so the corrected angle stays in
the interval. The argument takes standard parts of the two slice coordinates, not of the
symbol $I$, so it also works for a surreal imaginary unit.
\end{proof}

\begin{theorem}[Polar representation at arbitrary surreal radius]\label{squat:thm:polar}
Every nonzero surquaternion is expressible as
\[
 q=\rho(\cos\theta+I\sin\theta),\qquad \rho=\abs q>0,\quad
 \theta\in\OO_{\No},\quad I^2=-1 .
\]
For noncentral $q$ one may uniquely take $I=\Ima q/\abs{\Ima q}$ and $0<\theta<\pi$ with
$\cos\theta=\Rea q/\rho$ and $\sin\theta=\abs{\Ima q}/\rho$. Positive central elements use
$\theta=0$, negative ones $\theta=\pi$ with arbitrary $I$.
\end{theorem}
\begin{proof}
The normalized scalar and imaginary modulus lie on the unit circle with nonnegative second
coordinate; apply \cref{squat:prop:finiteangles}. The sign of the sine selects the open
upper semicircle and the unique representative in $(0,\pi)$, including the cases where the
angle or its defect from $\pi$ is infinitesimal, as is seen from
\cref{squat:eq:finite-trig}.
\end{proof}

An infinite modulus $\rho$ therefore does not require an infinite angle: the surreal scale
resides in the radius, and even an imaginary direction with infinitesimal coordinate
components has a finite angle. Every unit surquaternion is $\cos\theta+I\sin\theta$ with
$\theta\in[0,\pi]$, and formula \cref{squat:eq:rodrigues} becomes rotation about $I$
through the finite angle $2\theta$, as follows from the finite-angle double-angle
identities --- even for vectors of infinite or infinitesimal length.

\section{Three exponentials: the tube, the global radial map, and their laws}\label{squat:sec:exponentials}

The infinitesimal exponential $\Expm$ was constructed in \cref{squat:sec:explog}. This
section constructs the other two objects of \cref{squat:conv:exps} and records exactly
which law each satisfies. A global phase on all surreal scalars requires additional
structure, and we make that structure explicit rather than treating $\sum q^n/n!$ as an
everywhere-defined Hahn sum.

\subsection{The finite-angle tube exponential}\label{squat:sub:tube}
Let
\[
 \mathcal T=\{a+v\in\SQ:a\in\No,\ v\in\Ima\SQ,\ \abs v\text{ finite}\}
\]
and, using Gonshor's real exponential and the finite-angle trigonometry
\cref{squat:eq:finite-trig}, define
\begin{equation}\label{squat:eq:tubeexp}
 \Expt(a+v)=
 \begin{cases}
 \exp_{\No}(a)\left(\cos r+\dfrac{v}{r}\sin r\right),&r=\abs v>0,\\[2pt]
 \exp_{\No}(a),&v=0 .
 \end{cases}
\end{equation}
This is canonical within the chosen standard surreal exponential and restricted analytic
structure. Its domain contains every finite quaternion and all central surreals, but
\emph{not} $\omega\ii$.

\begin{theorem}[Polar logarithms and arbitrary-scale roots]\label{squat:thm:polarlog}
The map $\Expt$ is conjugation-equivariant, satisfies $\abs{\Expt(a+v)}=\exp_{\No}(a)$,
obeys the addition law on commuting pairs in $\mathcal T$, and is surjective onto
$\SQ^\times$: if $q=\rho(\cos\theta+I\sin\theta)$ with $\rho=\abs q>0$ and
$0\le\theta\le\pi$ as in \cref{squat:thm:polar}, then
\begin{equation}\label{squat:eq:polarlog}
 \Expt(\log_{\No}\rho+I\theta)=q ,
\end{equation}
and for every ordinary integer $n\ge1$ one $n$th root of $q$ is
\begin{equation}\label{squat:eq:nthroot}
 \rho^{1/n}\left(\cos(\theta/n)+I\sin(\theta/n)\right).
\end{equation}
\end{theorem}
\begin{proof}
Inner conjugation fixes the scalar part and the modulus of the imaginary part and carries
its direction to the conjugated direction. The modulus formula is
$\cos^2r+\sin^2r=1$. A commuting pair lies in a common slice by
\cref{squat:thm:conjugacy}; there \cref{squat:eq:tubeexp} reads
$\exp_{\No}(a)(\cos b+I\sin b)$ with a signed finite coordinate $b$, and the addition law
is the finite sine and cosine addition formula. Every unit quaternion has the polar
representation of \cref{squat:thm:polar}, which gives \cref{squat:eq:polarlog}; applying
the addition formula $n$ times in the same slice gives \cref{squat:eq:nthroot}.
\end{proof}

There is no single-valued logarithm implicit here. On the negative scalar axis every
$\log_{\No}\rho+\pi I$ with $I\in\mathbb S(\No)$ exponentiates to $-\rho$, and even off
that axis ordinary $2\pi$ periods give other logarithms. Surjectivity of a quaternion
exponential need not entail a preferred logarithm.

On finite quaternions, \cref{squat:eq:tubeexp} agrees with the four-real-variable Taylor
extension of the ordinary quaternion exponential: at a nonzero ordinary imaginary radius
this follows by analytic composition with the square root, and at radius zero the
removable function $\sin r/r$ is evaluated by its even power series.

\begin{warning}[The tube exponential is neither a power series nor a product]
\label{squat:warn:notaylor}
If $q=q_0+\eps$ with $q_0$ an ordinary quaternion and $\eps$ infinitesimal, then
$\Expt(q)$ is \emph{not} in general $\exp(q_0)\Expm(\eps)$: the two summands need not
commute. Nor is \cref{squat:eq:tubeexp} obtained by Hahn-summing $\sum q^n/n!$ for every
$q\in\mathcal T$, and the two failure modes are distinct: at a nonzero ordinary scalar all
those terms have exponent zero, violating finite overlap in
\cref{squat:def:summable}; at the central $q=\omega\in\mathcal T$ their exponents
are unbounded below. The latter failure also occurs for $q=\omega\ii$, which lies
outside the tube. Neither failure is repaired by citing ordinary real convergence.
\end{warning}

\subsection{The Ehrlich--Kaplan scalar normalization}\label{squat:sub:ekphase}
Split a scalar normal form into its purely infinite and finite parts,
\[
 x=x_{\mathrm{PI}}+\fin(x),\qquad \fin(x)=r+\eta\in\OO_{\No}.
\]
The projection $\fin$ is \emph{additive} but not multiplicative, and $\fin(\omega)=0$. Put
\begin{equation}\label{squat:eq:global-phase}
 \Sin x=\sin(\fin x),\qquad \Cos x=\cos(\fin x),
\end{equation}
the canonical scalar phase functions associated with the initial integer part, the omnific
integers
\[
 \Oz=\No_{\mathrm{PI}}+\Z .
\]
Indeed $2\pi\Oz=\No_{\mathrm{PI}}+2\pi\Z$, so reducing modulo this period group amounts
exactly to discarding the purely infinite part and then reducing the finite angle;
equivalently, for each $x\in\No$ one takes the unique $d\in\Oz$ with
$x-2\pi d\in[0,2\pi)$ and evaluates the restricted real values there. This is the
normalization of Ehrlich and Kaplan, Section~11 \cite{EK}. It gives
\begin{equation}\label{squat:eq:omegaphase}
 \Sin\omega=0,\qquad \Cos\omega=1 ,
\end{equation}
and on each complex slice a surjective exponential with period group $2\pi\Oz I$.

These are consequences of a \emph{specified normalization}, not of ordered-field algebra
alone, and the repository's trigonometry report \cite{RepoTrig} presents this character as
one member of a family and explicitly does not claim it as new. We inherit that caveat: a
different choice of global circular character would change the objects computed in
\cref{squat:thm:logfibres,squat:thm:exp-derivative}, and comparing such alternatives is
left open in \cref{squat:sub:directions}. Global surreal trigonometry must not be dismissed
as impossible; it must be named.

\subsection{The global radial exponential}\label{squat:sub:radial}

\begin{definition}[Global radial surquaternion exponential]\label{squat:def:globalexp}
For $q=a+v$ with $r=\abs v$, set
\begin{equation}\label{squat:eq:globalexp}
 \Expr(q)=
 \begin{cases}
 \exp_{\No}(a)\bigl(\Cos r+(v/r)\Sin r\bigr),&r>0,\\
 \exp_{\No}(a),&r=0 .
 \end{cases}
\end{equation}
\end{definition}

The construction is \emph{radial}: it uses the actual modulus $r$ and the actual direction
$v/r$ of the imaginary part, and only then applies the chosen scalar phase to $r$.
Discarding the purely infinite parts of the three vector coordinates first would define a
different and in general non-equivariant operation. Both source manuscripts define exactly
this map; one writes $\Exp_{\mathrm{EK}}$ and the other $\Exp_{\mathrm{can}}$, and
\cref{squat:conv:exps}(c) fixes the single name and symbol used here.

\begin{theorem}[Properties of the global radial exponential]\label{squat:thm:globalexp}
The map \cref{squat:eq:globalexp} is a surjection $\SQ\to\SQ^\times$ satisfying
\begin{gather*}
 \abs{\Expr(q)}=\exp_{\No}(\Rea q),\qquad
 \Expr(-q)=\Expr(q)^{-1},\qquad
 \Expr(\bar q)=\overline{\Expr(q)},\\
 \Expr(hqh^{-1})=h\,\Expr(q)\,h^{-1}\quad(h\ne0).
\end{gather*}
It agrees with the finite-angle $\Expt$ on finite inputs and with $\Expm$ on $\mm_H$, and
it restricts on every slice $\No+\No I$ to the canonical surcomplex exponential, a
surjective homomorphism with kernel $2\pi\Oz I$. It satisfies
$\Expr(p+q)=\Expr(p)\Expr(q)$ whenever $p$ and $q$ commute. It is \emph{not} an
additive-to-multiplicative homomorphism on noncommuting arguments, and consequently its
fibre over $1$ must not be called a kernel. In particular, since $\omega/(2\pi)\in\Oz$,
\begin{equation}\label{squat:eq:expomegai}
 \Expr(\omega\ii)=1 .
\end{equation}
\end{theorem}
\begin{proof}
The modulus, inverse and conjugation formulas follow from the scalar phase identities of
\cref{squat:sub:ekphase}. Inner conjugation preserves $a$ and $r$ and carries the direction
$I$ to $hIh^{-1}$, proving equivariance. On a slice, writing $q=a+Ib$ turns
\cref{squat:eq:globalexp} into $\exp_{\No}(a)(\Cos b+I\Sin b)$, using oddness of $\Sin$
when $b<0$; the slice is commutative, so the scalar phase homomorphism property and
\cref{squat:prop:finiteangles} give the slicewise homomorphism with the stated period
group. Commuting noncentral quaternions lie in one slice by \cref{squat:thm:conjugacy}, and
central cases cause no change. Given $q\ne0$, \cref{squat:thm:polar} gives
$\Expr(\log_{\No}\abs q+I\theta)=q$, which is surjectivity. For infinitesimal $q=a+v$ the
central scalar $a$ commutes with $v$, and the even and odd powers of $v$ are exactly the
Hahn cosine and sine series at $r$, so expanding the commuting product gives the local
series and hence agreement with $\Expm$. Agreement with $\Expt$ on finite inputs holds
because $\fin$ is the identity there. Failure of the global homomorphism law is witnessed by
\cref{squat:ex:radialnotproduct}; the additive quaternion group is abelian while the
multiplicative group is not, so no global homomorphism onto $\SQ^\times$ is possible, and
a noncommuting addition law is neither expected nor asserted. Finally
$\fin(\omega)=0$ gives \cref{squat:eq:omegaphase} and hence \cref{squat:eq:expomegai}.
\end{proof}

Equation \cref{squat:eq:expomegai} does not identify a divergent or nonsummable Taylor
series with $1$. It computes a value of a separately defined global operation, whose
domain and law are those recorded in \cref{squat:conv:exps}(c).

\begin{theorem}[Complete logarithm fibres]\label{squat:thm:logfibres}
For noncentral $q=\rho(\cos\theta+I\sin\theta)$ in the canonical polar form with
$0<\theta<\pi$,
\begin{equation}\label{squat:eq:logfibre-nonreal}
 \Expr^{-1}(q)=\{\log_{\No}\rho+I(\theta+2\pi z):z\in\Oz\} .
\end{equation}
For $\rho>0$ the central fibres are
\begin{align}
 \Expr^{-1}(\rho)&=\{\log_{\No}\rho+I(2\pi z):I^2=-1,\ z\in\Oz\},
 \label{squat:eq:positive-fibre}\\
 \Expr^{-1}(-\rho)&=\{\log_{\No}\rho+I\pi(2z+1):I^2=-1,\ z\in\Oz\}.
 \label{squat:eq:negative-fibre}
\end{align}
There is no preimage of zero.
\end{theorem}
\begin{proof}
The modulus formula determines the scalar part of every preimage. The image of a
quaternion lies in its own slice; if the image is noncentral, its centralizer is that
unique slice by \cref{squat:thm:conjugacy}, so the preimage has the form
$\log_{\No}\rho+Ib$ for a scalar $b$, and \cref{squat:prop:finiteangles} with
\cref{squat:eq:global-phase} makes its phase fibre exactly $b\in\theta+2\pi\Oz$. When the
image is central its imaginary part vanishes, so its radial phase is $1$ or $-1$ and every
axis is allowed, which gives \cref{squat:eq:positive-fibre,squat:eq:negative-fibre},
including the repetition of the scalar logarithm at $z=0$ in the positive case. Moduli are
strictly positive, excluding zero.
\end{proof}

For positive and negative central values the fibres therefore sweep over \emph{all}
imaginary units $I$. These fibre descriptions do not settle the robustness questions for
alternative surcomplex exponentials posed in \cite{EK}.

\begin{example}[Infinite radius is not coordinatewise phase reduction]
\label{squat:ex:radialnotproduct}
With $q=\omega\ii+\pi\jj$,
\[
 r=\sqrt{\omega^2+\pi^2}=\omega+\frac{\pi^2}{2\omega}+\mathrm O(\omega^{-3}),
 \qquad\text{so}\qquad
 \Expr(q)=1+\frac{\pi^2}{2\omega}\ii+\mathrm O(\omega^{-2}),
\]
whereas $\Expr(\omega\ii)\Expr(\pi\jj)=-1$. Here and below, order symbols for Hahn
expansions bound the valuation of the remainder; they are not sequence limits. The example
shows both that the direction and radius must be formed before the purely infinite phase is
discarded, and that the homomorphism law genuinely fails at noncommuting arguments.
Compare \cref{squat:sub:consistency}: on either single slice the law does hold.
\end{example}

\section{The differential of the radial exponential and its critical spheres}\label{squat:sec:differential}

Differentiating a quaternionic function with respect to four scalar coordinates is not the
same operation as applying a scalar-field derivation, and neither is the coefficientwise
slice derivative. We use the following, which is \cref{squat:conv:derivs}(c).

\begin{definition}[Fine differential]\label{squat:def:fine-derivative}
An $\No$-linear map $L$ is the differential of $f$ at $q$ when for every surreal
$\epsilon>0$ there is a surreal $\delta>0$ with
\[
 0<\abs h<\delta\ \Longrightarrow\ \abs{f(q+h)-f(q)-L(h)}<\epsilon\abs h .
\]
This is meaningful as a class scheme even though set-indexed fine limits are degenerate by
\cref{squat:thm:discrete}. It does not assert a real Banach-space structure, and a
nonsingular fine differential is never used below as an inverse function theorem; the
positive existence tool is \cref{squat:thm:lifting}.
\end{definition}

\begin{theorem}[Radial derivative and critical spheres at infinite radius]
\label{squat:thm:exp-derivative}
Let $q=a+rI$ with $r>0$, and write a variation as $h=s+\alpha I+u$, where
$s,\alpha\in\No$ and $u$ is imaginary and perpendicular to $I$. Put
$E=\exp_{\No}(a)$, $C=\Cos r$, $S=\Sin r$. Then
\begin{equation}\label{squat:eq:exp-derivative}
 d\Expr(q)[h]=E\left(s(C+IS)+\alpha(-S+IC)+\frac{S}{r}\,u\right),
\end{equation}
and at a central point $a$ the differential is $h\mapsto\exp_{\No}(a)h$. As a
four-dimensional $\No$-linear map,
\begin{equation}\label{squat:eq:exp-jacobian}
 \det d\Expr(a+rI)=\exp_{\No}(4a)\left(\frac{\Sin r}{r}\right)^2\qquad(r>0),
\end{equation}
which is singular precisely for positive $r\in\pi\Oz$, and has rank exactly two there.
\end{theorem}
\begin{proof}
Fix $q$. For a variation $u'=\alpha I+u$ small relative to $r$, positive-square-root
expansion gives
\[
 \abs{rI+u'}=r+\alpha+\mathrm O(\abs{u'}^2/r),\qquad
 \frac{rI+u'}{\abs{rI+u'}}=I+\frac{u}{r}+\mathrm O(\abs{u'}^2/r^2),
\]
Taylor identities on infinitesimal relative increments, justified by
\cref{squat:lem:neumann}; alternatively the modulus expansion follows by rationalizing the
difference of two square roots. The constants in their bounds may depend on the fixed
surreal $r$. For infinitesimal scalar $\eta$ the additive phase law gives
$\Cos(r+\eta)=C-S\eta+\mathrm O(\eta^2)$ and $\Sin(r+\eta)=S+C\eta+\mathrm O(\eta^2)$
\emph{even when $r$ is infinite}, because $\fin$ is additive; likewise
$\exp_{\No}(a+s)=E(1+s+\mathrm O(s^2))$. Combining the three expansions proves
\cref{squat:eq:exp-derivative} with a remainder bounded by $M_q\abs h^2$ on a sufficiently
small ball for some positive surreal $M_q$, and shrinking the ball below
$\epsilon/(M_q+1)$ gives the differential in the sense of
\cref{squat:def:fine-derivative}. At $r=0$ use the Hahn exponential at the small increment
$h$ together with the commuting central scalar $a$.

On the scalar--radial plane the matrix is
$E\left(\begin{smallmatrix}C&-S\\S&C\end{smallmatrix}\right)$, of determinant $E^2$, and
on the two transverse coordinates it is $(ES/r)$ times the identity; this gives
\cref{squat:eq:exp-jacobian}. Finally $\Sin r=0$ precisely when $\fin r\in\pi\Z$,
equivalently $r\in\pi\Oz$, and the scalar--radial block remains invertible, so a singular
differential has rank exactly two.
\end{proof}

\begin{example}[A critical sphere of infinite radius]\label{squat:ex:criticalsphere}
At $q=\omega I$ the normalized phase has $C=1$, $S=0$, so
\[
 d\Expr(\omega I)[s+\alpha I+u]=s+\alpha I .
\]
Every transverse first-order variation vanishes while the scalar and radial ones survive.
For a central formal small parameter $\eta$ and $J\perp I$,
\[
 \Expr(\omega I+\eta J)=1+\frac{\eta^2}{2\omega}I+\frac{\eta^3}{2\omega^2}J
 +\text{terms of higher valuation}.
\]
This is not obtained by exponentiating each imaginary coordinate independently; it is a
radial critical-point phenomenon at an infinite scale.
\end{example}

The determinant formula extends the familiar finite-radius criticality of the quaternion
exponential, but the allowed radii now include a proper class of infinite integer-part
periods, the infinite radius $\omega$ among them. That this happens at all is a consequence
of the chosen normalization of \cref{squat:sub:ekphase}, not of ordered-field algebra.

\section{Derivations and the obstruction to an oscillatory calculus}\label{squat:sec:derivations}

\subsection{Every algebraic derivation has a scalar and an inner part}\label{squat:sub:derivclass}
Let $\delta:F\to F$ be a field derivation, meaning an additive map with
$\delta(xy)=x\delta(y)+\delta(x)y$.

\begin{theorem}[Derivation decomposition]\label{squat:thm:derivations}
The formula
\begin{equation}\label{squat:eq:coeffderiv}
 \delH(a+b\ii+c\jj+d\kk)=\delta(a)+\delta(b)\ii+\delta(c)\jj+\delta(d)\kk
\end{equation}
defines the unique derivation of $\HF$ extending $\delta$ and annihilating $\ii,\jj,\kk$.
It commutes with conjugation, and
\begin{equation}\label{squat:eq:normderiv}
 \delta N(q)=2\Rea\bigl(\delH(q)\bar q\bigr),\qquad
 \delta\abs q=\frac{\Rea(\delH(q)\bar q)}{\abs q}\quad(q\ne0).
\end{equation}
Every ring derivation $D$ of $\HF$ has the unique form
\begin{equation}\label{squat:eq:derivdecomp}
 D(q)=\delH(q)+[w,q],\qquad w\in\Ima\HF,
\end{equation}
where $\delta$ is the restriction of $D$ to the center, and the inner representative is
explicitly
\begin{equation}\label{squat:eq:innerrep}
 w=-\tfrac14\bigl(E(\ii)\ii+E(\jj)\jj+E(\kk)\kk\bigr),\qquad E=D-\delH .
\end{equation}
\end{theorem}
\begin{proof}
Hamilton multiplication has integer structure constants, which proves the Leibniz rule for
the coefficientwise extension and commutation with conjugation; differentiating
$N(q)=q\bar q$ and then $\abs q^2=N(q)$ gives \cref{squat:eq:normderiv}.

Differentiating $aq-qa=0$ for central $a$ shows that $D(a)$ commutes with every
quaternion, so $D$ restricts to a derivation of $F$; subtracting the componentwise
extension of that restriction leaves an $F$-linear derivation $E$. Two independent routes
now identify $E$, and both were carried out in the source manuscripts.

\emph{Route one.} For imaginary $u$, differentiate $u^2=-N(u)$. Since $E$ vanishes on $F$,
this forces $E(u)$ to have zero scalar part and to be orthogonal to $u$ (first for
$u\ne0$). Applying this to $u+v$ shows that $E$ restricted to $V=\Ima\HF$ is
skew-symmetric.

\emph{Route two.} Differentiate $\ii^2=\jj^2=\kk^2=-1$ to see that
$E(\ii),E(\jj),E(\kk)$ are imaginary and perpendicular to the corresponding basis vectors,
then differentiate the anticommutation relations to see that the $3\times3$ matrix of $E$
on $V$ is skew-symmetric.

Either way, every skew-symmetric $3\times3$ matrix is uniquely $u\mapsto2w\times u$ for a
vector $w\in V$, by its three independent entries, and by \cref{squat:prop:center} that map
is $[w,\,\cdot\,]$; both maps vanish on $F$, which proves
\cref{squat:eq:derivdecomp,squat:eq:innerrep}. Uniqueness follows because an imaginary
element commuting with all of $\HF$ is zero.
\end{proof}

This is a concrete quaternionic form of the inner-derivation property of central simple
algebras, and no infinite-dimensional algebra theorem is needed. For the componentwise
extension one has
\begin{align}\label{squat:eq:derivformulas}
 \delH(q^{-1})&=-q^{-1}\delH(q)q^{-1},
 &\delH(q^n)&=\sum_{m=0}^{n-1}q^m(\delH q)q^{n-1-m}.
\end{align}
The last formula becomes $nq^{n-1}\delH q$ \emph{only} when the relevant factors commute.
The inner part in \cref{squat:eq:derivdecomp} gives an algebraic analogue of a change of
quaternionic frame; its commutator cannot be absorbed into a scalar derivative of the
coefficients.

\subsection{The Berarducci--Mantova extension}\label{squat:sub:bm}
Fix the Berarducci--Mantova derivation $\partial$ on $\No$ normalized by
$\partial\omega=1$. We use the established properties that it is strongly additive,
commutes with the real surreal exponential, has constant field exactly $\R$, is
surjective, and maps infinitesimals to infinitesimals \cite[Theorems A and B]{BM}. These
are external scalar results, not consequences of the quaternion construction. Let $\dH$
denote the componentwise extension to $\SQ$, with $\dH\ii=\dH\jj=\dH\kk=0$.

\begin{corollary}[Constants, surjectivity, antiderivatives]\label{squat:cor:bmconstants}
$\ker\dH=\HR$, so the constant algebra of $\dH$ is exactly the ordinary quaternions;
$\dH$ is surjective, and every surquaternion has a coefficientwise antiderivative;
and $\dH$ respects Hahn-summable families coefficientwise.
\end{corollary}
\begin{proof}
An element is annihilated exactly when each of its four coordinates lies in
$\ker\partial=\R$. Choosing one scalar antiderivative for each of the four coordinates
gives surjectivity, and strong additivity gives the last assertion.
\end{proof}

Surjectivity of $\dH$ is \emph{not} a theorem that arbitrary nonlinear quaternionic
differential equations have solutions, and a set-sized Hahn workspace need not be stable
under $\partial$: statements about a differential workspace require closure under that
derivation or an explicit enlargement.

For an infinitesimal $x$, strong additivity of the \emph{chosen Berarducci--Mantova}
derivation justifies differentiating its Hahn exponential term by term, giving
\begin{equation}\label{squat:eq:dexp}
 \dH\Expm(x)=\sum_{n\ge1}\frac1{n!}\sum_{m=0}^{n-1}x^m\,\dH(x)\,x^{n-1-m},
\end{equation}
whose right-hand family is support-controlled: it is built from words in the positive
support of $x$ with one distinguished factor from $\supp\dH(x)$, and
\cref{squat:lem:neumann} with the finite antidiagonals of sums of well-ordered sets
justifies the regrouping. An arbitrary scalar derivation $\delta$ need not be strongly
additive, so this passage does not follow merely from the algebraic extension theorem.
For completeness, write $y=\dH(x)$. The formal noncommutative exponential identity gives
\[
 \Expm(-x)\dH\Expm(x)-y
 =-\tfrac12[x,y]+\sum_{n\ge2}\frac{(-1)^n}{(n+1)!}\ad_x^n(y).
\]
The same support argument justifies this evaluation. If $[x,y]\ne0$, every term in
the sum has valuation at least $w([x,y])+w(x)>w([x,y])$, so cannot cancel the first
term. Thus the simpler expression $\dH\Expm(x)=\Expm(x)\dH(x)$ holds exactly when
$[x,\dH(x)]=0$ (the case $x=0$ is immediate). This local assertion supplies no global
commuting chain rule.

\begin{proposition}[Angular velocity identity]\label{squat:prop:angular}
Let $u\in\HF$ satisfy $N(u)=1$ and set $\xi=\delH(u)u^{-1}$. Then $\xi$ is imaginary, and
for every $x\in\HF$
\begin{equation}\label{squat:eq:angular-velocity}
 \delH(uxu^{-1})=[\xi,uxu^{-1}]+u\,\delH(x)\,u^{-1}.
\end{equation}
For constant imaginary $x$ the derivative of its rotated value is
$2\xi\times(uxu^{-1})$.
\end{proposition}
\begin{proof}
Differentiate $u\bar u=1$ to get $\xi+\bar\xi=0$. The inverse rule is
$\delH(u^{-1})=-u^{-1}\delH(u)u^{-1}$; apply Leibniz to $uxu^{-1}$ and collect terms.
\end{proof}

This formula applies to surreal differential coefficients without pretending that the
parameter is an ordinary real-time path into the fine topology. It provides an algebraic
angular-velocity equation, not an existence theorem for trajectories.

\subsection{No constant-frequency oscillation, and no commuting chain rule}
\label{squat:sub:nooscillation}

\begin{lemma}[Finite inputs have infinitesimal derivatives]\label{squat:lem:small}
If $q\in\OO_H$ then $\dH q\in\mm_H$.
\end{lemma}
\begin{proof}
Write $q=\st(q)+\eps$. The standard part lies in $\HR$ and is killed by $\dH$. Each
coordinate of $\eps$ is infinitesimal, and its derivative is infinitesimal by the
smallness property of $\partial$. The four-coordinate assertion follows.
\end{proof}

\begin{theorem}[No constant-frequency quaternionic oscillation]\label{squat:thm:nooscillation}
Let $I\in\mathbb S(\R)$ and $c\in\R\setminus\{0\}$. Neither equation
\begin{equation}\label{squat:eq:oscillation}
 \dH y=cIy\qquad\text{nor}\qquad \dH y=ycI
\end{equation}
has a nonzero solution in $\SQ$. In particular $\dH y=\ii y$ has only the zero solution,
even inside a fixed surcomplex slice.
\end{theorem}
\begin{proof}
For the left equation,
\[
 \dH N(y)=\overline{\dH y}\,y+\bar y\,\dH y=-\bar y\,cIy+\bar y\,cIy=0,
\]
and the right equation gives the same cancellation, using centrality of $N(y)$. Thus
$N(y)$ lies in the constant field $\R$. If $y\ne0$, this is a positive ordinary real
number, so $\abs y=\sqrt{N(y)}$ is a positive ordinary real and every coordinate of $y$ is
finite. By \cref{squat:lem:small}, $\dH y$ has infinitesimal coordinates, so its modulus is
infinitesimal or zero. But the right-hand side of either equation has modulus
$\abs c\,\abs y$, a positive ordinary real. This is impossible.
\end{proof}

\begin{corollary}[No global commuting chain rule]\label{squat:cor:nochain}
There is no map $E:\SQ\to\SQ^\times$ satisfying
\[
 \dH(E(q))=E(q)\,\dH q
\]
for every $q$ with $[q,\dH q]=0$, and the same conclusion holds with the factors on the
right-hand side reversed. In particular the usual commuting exponential chain rule cannot
hold at $q=\omega\ii$.
\end{corollary}
\begin{proof}
At $q=\omega\ii$ one has $\dH q=\ii$ and $[q,\dH q]=0$. A nonzero $E(q)$ would solve one of
the equations forbidden by \cref{squat:thm:nooscillation}.
\end{proof}

For the global normalization of \cref{squat:eq:globalexp} the contradiction is visible as
an exact identity:
\begin{equation}\label{squat:eq:chainwitness}
 \dH\Expr(\omega\ii)=\dH1=0,
 \qquad \Expr(\omega\ii)\,\dH(\omega\ii)=\ii .
\end{equation}
The normalization is still a valid global algebraic exponential, and by
\cref{squat:thm:globalexp} a genuine homomorphism on each commuting slice. It simply does
not have this differential compatibility.

The proof of \cref{squat:thm:nooscillation} uses the constant field and the smallness
property, not any assumption about the global phase. It therefore rules out an entire
compatibility \emph{requirement}, rather than exposing a defect of one phase choice: any
proposed unified surquaternion calculus must respect this boundary.

\begin{remark}[Relation to the companion results, and what is not being claimed]
\label{squat:rem:oscillation-relations}
Three separations must be kept.
\begin{enumerate}[label=\textup{(\arabic*)},leftmargin=2.4em]
\item \Cref{squat:thm:nooscillation} is the componentwise quaternionic
\emph{extension} of ``the missing oscillator'' of \cite{RepoDiff}, where the scalar and
surcomplex statement --- that $\partial^2y+y=0$ has only $y=0$, and that $\sin\omega$ is
not a Berarducci--Mantova differential element --- is proved three independent ways. We
cite that result and do not restate its proofs, and we do not present the quaternionic
version as an independent discovery.
\item These conclusions do not conflict with the existence of real antiderivatives or with
the Liouville-closed real differential-field properties of \cite{BM}: a constant imaginary
frequency is not a scalar coefficient in that real differential field. Adjoining formal
oscillations would require a larger differential algebra or a different derivation, and
ordinary quaternion scalar extension by itself does not produce them.
\item \Cref{squat:cor:nochain} concerns the single derivation $\dH$ and says nothing about
the multiplicative behaviour of $\Expr$. It is therefore consistent both with the
slicewise homomorphism property of \cref{squat:thm:globalexp} --- hence with the
commutative surcomplex homomorphism claim of \cite{RepoPhys} --- and with the fine
differential computed in \cref{squat:thm:exp-derivative}. See
\cref{squat:sub:consistency}.
\end{enumerate}
\Cref{squat:thm:nooscillation,squat:cor:nochain} are proved here from the cited properties
of the derivation. We make no claim that this formulation has never appeared in the
literature; their role is to identify a concrete boundary.
\end{remark}

\section{Quaternionic linear algebra and multiscale spectral stability}\label{squat:sec:spectral}

The complex and real closed case of this material is already committed, in the companion
report \emph{Surcomplex Spectral Theory: Singular Values, Infinitesimal Rank, and
Multiscale Stability} at
\texttt{docs/surcomplex/\allowbreak spectral-theory/\allowbreak article.tex}
\cite{RepoSpec}. That report proves Hermitian and
normal diagonalization, positive square roots, the singular value decomposition, polar
decomposition and the variational principles with attained extrema over a real closed
field and its complexification, and over a Hahn workspace it proves gap-sensitive spectral
projectors and subspace perturbation. We cite that report for those statements and print
here only what noncommutativity changes.

What changes is the setting itself. On the \emph{right} $\HF$-module $\HF^n$ use
\[
 \angles{x,y}=\sum_{m=1}^n\bar x_my_m,\qquad \abs x=\sqrt{\angles{x,x}}\in F_{\ge0},
\]
with matrices acting on the \emph{left}, the adjoint $A^*=\bar A^{\,T}$ the conjugate
transpose, and $U$ unitary when $U^*U=I$. These conventions matter whenever eigenvalues or
coefficients are moved across vectors: an eigenvalue must be written on the right,
$Ax=x\lambda$, and it is a theorem, not a convention, that a Hermitian matrix has its
eigenvalues in the \emph{center} $F$. Orthogonalization uses finite sums and positive
square roots, so quaternionic Gram--Schmidt is a finite algebraic procedure over a real
closed field.

\begin{theorem}[Hermitian spectral theorem over $\HF$]\label{squat:thm:spectral}
Let $F$ be real closed and $A=A^*\in M_n(\HF)$. There are a quaternionic unitary $U$ and
central scalars $\lambda_1,\ldots,\lambda_n\in F$ with
\[
 U^*AU=\diag(\lambda_1,\ldots,\lambda_n),
\]
and $A$ is positive definite exactly when every $\lambda_m>0$. Every positive definite
Hermitian matrix has a unique positive definite Hermitian square root; every matrix admits
a singular value decomposition with nonnegative central singular values; and every
invertible square matrix has a polar decomposition $A=UP$ with $U$ unitary and
$P=(A^*A)^{1/2}$ positive definite Hermitian, uniquely.
\end{theorem}
\begin{proof}
Write $A=Z+W\jj$ with matrices over $C_F=F[\ii]$ and extend \cref{squat:eq:rho} in blocks,
\[
 \rho(A)=\begin{pmatrix}Z&W\\-\bar W&\bar Z\end{pmatrix},
\]
which is Hermitian over the algebraically closed field $C_F$ when $A$ is. Its
characteristic polynomial has a root $\lambda\in C_F$ with a nonzero eigenvector $v$; since
$v^*v>0$ and $v^*\rho(A)v$ is fixed by conjugation, $\lambda=(v^*\rho(A)v)/(v^*v)$ lies in
$F$. Two equivalent routes return from $\rho(A)$ to $A$: the identification
$x+y\jj\mapsto(x,-\bar y)^T$ intertwines the action of $\rho(A)$ with that of $A$; or,
equivalently, every column $v\in C_F^{2n}$ is obtained by stacking the first columns of
$\rho(x_m)$ for a unique $x\in\HF^n$. Because $\lambda$ is central, the matrix
eigen-equation becomes $Ax=x\lambda$, a right eigenvector with the same central
eigenvalue. Normalize $x$ by the positive square root of its norm square. Its quaternionic
orthogonal complement is $A$-invariant: if $\angles{x,y}=0$ then
$\angles{x,Ay}=\angles{Ax,y}=\lambda\angles{x,y}=0$, using $A=A^*$ and centrality of
$\lambda$. Quaternionic Gram--Schmidt produces an orthonormal basis of that complement and
induction on $n$ gives the diagonalization. The positivity criterion follows by evaluating
the diagonal quadratic form, including on basis vectors.

For a positive definite matrix, replace each positive diagonal eigenvalue by its positive
square root; uniqueness holds because a positive square root commutes with its square,
preserves its eigenspaces, and on each eigenspace can only have the positive eigenvalue
$\sqrt\lambda$. Diagonalizing $A^*A$ gives singular values; for each nonzero singular
value $\sigma$, normalize $Av$ by $\sigma$ and complete orthonormal bases of the remaining
spaces, which proves the singular value decomposition. For invertible $A$, writing
$A=(A(A^*A)^{-1/2})(A^*A)^{1/2}$ gives the polar factors.
\end{proof}

No appeal to compactness of a surreal unit sphere occurs in this proof: it uses only
finite linear algebra and algebraic closedness of the complexification of the center.
Polar decomposition here is algebraic and does not assert compactness of $S^3(\No)$. No
infinite-dimensional Hilbert-space or proper-class spectral theorem is asserted anywhere.

\subsection{Matrix modulus and its valuation}\label{squat:sub:matrixnorm}
Define $\norm A$ to be the largest singular value of $A$. The singular value decomposition
proves $\abs{Ax}\le\norm A\abs x$ with equality for a suitable unit vector, and for a
square $n\times n$ matrix
\begin{equation}\label{squat:eq:matrixnorm}
 \max_{m,m'}\abs{a_{mm'}}\le\norm A\le n\max_{m,m'}\abs{a_{mm'}},
\end{equation}
whence for a nonzero surquaternion matrix
\begin{equation}\label{squat:eq:matrixvaluation}
 w(\norm A)=\min_{m,m'}w(a_{mm'}).
\end{equation}
Matrix multiplication may cancel leading terms, so only $w(AB)\ge w(A)+w(B)$ is automatic;
the equality of \cref{squat:eq:valuation-product}, valid for individual nonzero
quaternions, must not be copied to matrices.

\begin{warning}[Two determinant arguments that are unavailable here]\label{squat:warn:nodet}
The companion report \cite{RepoSpec} proves a determinantal singular-scale theorem, in
which the $k$th determinantal valuation recovers the singular scales, and it uses a
positive Cauchy--Binet expansion of the minors of a product. These proofs use the
ordinary determinant of matrices over a commutative field; their minor identities cannot
be transferred unchanged to $M_n(\HF)$. Noncommutative determinants and the complex
matrix realization provide other tools, but relating them to quaternionic singular
scales would require a separate argument. No such determinant theorem is proved here;
the multiscale results of this section use the min--max and projector arguments.
That report's explicit warning about Gram squaring applies here
verbatim and is inherited, not restated.
\end{warning}

\subsection{Weyl's estimate without compactness}\label{squat:sub:weyl}
Order the eigenvalues of a Hermitian matrix increasingly. Finite-dimensional spectral
algebra gives the min--max formula
\[
 \lambda_k(A)=\min_{\dim W=k}\ \max_{x\in W,\ \abs x=1}\angles{x,Ax},
\]
where the span of the first $k$ eigenvectors realizes the upper bound, and any
$k$-dimensional right subspace meets the span of eigenvectors $k,\ldots,n$ nontrivially,
proving the lower bound. The inner maximum exists by applying \cref{squat:thm:spectral} to
the compressed Hermitian form on $W$. This notation therefore represents \emph{attained
extrema justified by linear algebra}, not by a general extreme-value theorem over $F$.
Since $\abs{\angles{x,Ex}}\le\norm E$ for unit $x$,
\begin{equation}\label{squat:eq:weyl}
 \abs{\lambda_k(A+E)-\lambda_k(A)}\le\norm E\qquad(A=A^*,\ E=E^*).
\end{equation}
If $A$ has finite entries, its eigenvalues are finite by \cref{squat:eq:matrixnorm}, and
reducing a unitary diagonalization by standard part shows that their ordered standard parts
are exactly the ordered eigenvalues of $\st(A)$, with multiplicities.

\begin{theorem}[Spectral-projector estimate at arbitrary surreal gap scale]
\label{squat:thm:projector}
Let $A=A^*\in M_n(\HF)$ with $n\ge2$ have a simple eigenvalue $\lambda$, and let
\[
 g=\min_{\mu\ne\lambda}\abs{\mu-\lambda}>0,\qquad e=\norm E<g/2,\qquad E=E^* .
\]
There is exactly one eigenvalue $\lambda'$ of $A+E$ in $(\lambda-g/2,\lambda+g/2)$,
counted with quaternionic multiplicity, and if $P,P'$ are the associated rank-one
orthogonal projections then
\begin{equation}\label{squat:eq:projectorbound}
 \norm{P'-P}\le\frac{e}{g-e}.
\end{equation}
For surreal coefficients the sufficient condition $w(E)>w(g)$ gives
\begin{equation}\label{squat:eq:projectorval}
 w(P'-P)\ge w(E)-w(g),
\end{equation}
zero matrices being interpreted as having valuation $+\infty$.
\end{theorem}
\begin{proof}
\Cref{squat:eq:weyl} places the matched ordered eigenvalue within $e$ of $\lambda$ and
every other perturbed eigenvalue at least $g-e>g/2$ from $\lambda$, which proves
uniqueness and simplicity. Let $v'$ be a unit vector for $\lambda'$ and write
$v'=Pv'+y$ with $y=(I-P)v'$. Applying $I-P$ to $(A+E-\lambda')v'=0$ gives
$(A-\lambda')y=-(I-P)Ev'$. On $\operatorname{im}(I-P)$ every spectral value of
$A-\lambda'$ has absolute value at least $g-\abs{\lambda'-\lambda}\ge g-e$, so
\cref{squat:thm:spectral} gives $\abs y\le e/(g-e)$.

For two rank-one orthogonal projections, $\norm{P'-P}=\abs{(I-P)v'}$. To verify this over
$\HF$, multiply one generating vector on the right by a unit quaternion so that its inner
product with the other is a nonnegative central scalar, and restrict to their
two-dimensional span; the resulting real symmetric $2\times2$ matrix has eigenvalues
$\pm\sqrt{1-\abs{\angles{v,v'}}^2}$, and the orthogonal complement contributes zero. This
proves \cref{squat:eq:projectorbound}. If $w(E)>w(g)$ then $e/g$ is infinitesimal, so
$e<g/2$ and $w(g-e)=w(g)$; taking valuations and using
\cref{squat:eq:matrixvaluation} yields \cref{squat:eq:projectorval}.
\end{proof}

The estimate is useful exactly when the standard-part matrix has a multiple eigenvalue: an
infinitesimal gap can still resolve its eigenvectors at a finer scale, and replacing $g$ by
an ordinary standard-part gap would discard precisely that information.

\begin{example}[A sharp loss of gap scale]\label{squat:ex:spectral}
Let $t=\omega^{-1}$, $0<\alpha<\beta$, and
\[
 A=\begin{pmatrix}0&t^\beta\ii\\-t^\beta\ii&t^\alpha\end{pmatrix},
\]
a perturbation of $\diag(0,t^\alpha)$ of modulus $t^\beta$ with gap $t^\alpha$. Its
eigenvalues are exactly
\begin{equation}\label{squat:eq:twoeigen}
 \lambda_\pm=\frac{t^\alpha\pm\sqrt{t^{2\alpha}+4t^{2\beta}}}{2},
\end{equation}
and the positive-order binomial expansion gives
\begin{align*}
 \lambda_-&=-t^{2\beta-\alpha}+t^{4\beta-3\alpha}
   +\text{terms of valuation at least }6\beta-5\alpha,\\
 \lambda_+&=t^\alpha+t^{2\beta-\alpha}-t^{4\beta-3\alpha}
   +\text{terms of valuation at least }6\beta-5\alpha .
\end{align*}
The lower spectral projector is $P_-=(\lambda_+I-A)/(\lambda_+-\lambda_-)$, whose
off-diagonal entry is
\[
 -\frac{t^\beta\ii}{\sqrt{t^{2\alpha}+4t^{2\beta}}}
 =-t^{\beta-\alpha}\ii+\text{terms of strictly higher valuation},
\]
while its diagonal changes have valuation $2(\beta-\alpha)$. Thus
$w(P_--\diag(1,0))=\beta-\alpha$, realizing the loss in
\cref{squat:eq:projectorval} \emph{exactly}, whereas the eigenvalue displacement occurs at
the finer scale $2\beta-\alpha$. The standard-part matrix is degenerate, so an
ordinary-gap version of the estimate would have discarded the information entirely.
\end{example}

\begin{example}[An infinitesimal avoided crossing]\label{squat:ex:crossing}
For central $a,\delta\in F$ and $z\in\HF$, the Hermitian matrix
\[
 A=\begin{pmatrix}a+\delta&z\\\bar z&a-\delta\end{pmatrix}
\]
satisfies $(A-aI)^2=(\delta^2+N(z))I$, so its eigenvalues are
$a\pm\sqrt{\delta^2+N(z)}$, the coincident case included. Taking $a=\omega$,
$\delta=\omega^{-2}$ and $z=\omega^{-3}\jj$ gives the exact gap
$2\omega^{-2}\sqrt{1+\omega^{-2}}$. The infinite common offset and the much smaller mixing
scale cause no ambiguity in the algebraic spectrum. The phrase is borrowed; nothing is
asserted here about atomic or molecular spectra.
\end{example}

Valuation error bounds of this kind are exact symbolic data. Truncating a Neumann inverse
after degree $M$ certifies \cref{squat:eq:neumannerror}, and a spectral perturbation can
carry both $w(E)$ and the gap valuation $w(g)$ to certify
\cref{squat:eq:projectorval}. These are stronger than a floating-point residual when the
important effect lives below every ordinary real scale. They do not by themselves certify a
limit in any topology, least of all one whose value group has larger scales than all
multiples of the current step (\cref{squat:warn:cofinality}). Substituting a small positive
\emph{real} $t$ in \cref{squat:ex:spectral} should be treated as an illustration, not as
the definition of the surreal problem.

\section{Support-controlled analysis: slice series, coherent families, Cauchy and Fueter}
\label{squat:sec:analysis}

A satisfactory analytic theory must relate the many complex slices without confusing the
four operators of \cref{squat:conv:derivs}: the coefficientwise slice derivative $\dsl$,
the field derivation $\dH$, the fine differential of
\cref{squat:def:fine-derivative}, and the Fueter operator $\DD$. They have different
domains and different defining properties. This section constructs a controlled extension
of classical slice analysis \cite{GP}; every infinite expression uses Hahn supports or
ordinary coefficient analysis, and no assertion about arbitrary class functions on $\SQ$ is
intended.

\subsection{Admissible series on an infinitesimal disk}\label{squat:sub:formalseries}
Consider right-coefficient series
\begin{equation}\label{squat:eq:formalseries}
 f(x)=\sum_{n\ge0}x^na_n .
\end{equation}
If every $a_n\in\HR$ and $x\in\mm_H$, this is Hahn-summable regardless of the ordinary
real radius of convergence of the coefficients: infinitesimal support, not a real analytic
estimate, does the work. More generally one may take $a_n\in\OO_D$ under the
\emph{common-support hypothesis} that all $\supp(a_n)$ lie in one well-ordered set
$S\subseteq\Gamma_{\ge0}$.

\begin{proposition}[Uniform-support evaluation]\label{squat:prop:evaluation}
Under the common-support hypothesis, \cref{squat:eq:formalseries} is Hahn-summable at every
$x\in\mm_H$ in a common workspace, and so is its formal derivative
$\dsl f=\sum_{n\ge1}nx^{n-1}a_n$. Products and formal differentiation are compatible with
Hahn regrouping whenever the same common-support conditions hold for the participating
series.
\end{proposition}
\begin{proof}
Let $T=\supp(x)$, which is positive. Every summand is supported in $S+nT$. The monoid
$T^*$ is well-ordered by \cref{squat:lem:neumann}, and $S+T^*$ is well-ordered. For a
fixed exponent there are only finitely many pairs consisting of an exponent of $S$ and one
of $T^*$, and only finitely many words representing the latter across all lengths, so only
finitely many summands contribute. Shifting the index by one proves the derivative
assertion, and the same word count proves the regroupings and the formal product rule.
\end{proof}

\begin{warning}[The slice derivative is not the differential]\label{squat:warn:slicederiv}
The derivative in \cref{squat:prop:evaluation} is the \emph{slice or formal} derivative
$\dsl$ of \cref{squat:conv:derivs}(b). It is not a claim that $df_x(h)=h\,\dsl f(x)$ for
arbitrary noncommuting $h$: even the polynomial $x^2$ has full algebraic differential
$h\mapsto xh+hx$ rather than $h\mapsto2hx$ in all directions, as
\cref{squat:prop:sylvester} makes exact.
\end{warning}

\begin{warning}[The common-support hypothesis is substantive, and so is the coefficient class]
\label{squat:warn:commonsupport}
Two independent failures must be recorded.
\begin{enumerate}[label=\textup{(\arabic*)},leftmargin=2.4em]
\item Knowing $w(a_n)\ge0$ individually is \emph{not} a substitute for a common $S$.
For a counterexample with integral coefficients, take $\Gamma=\Q\times\Q$ in
lexicographic order, realized inside $\No$ by $(a,b)\mapsto a\omega+b$, and put
$x=t^{(0,1)}$ and $a_n=t^{(1,-n^2)}$ for $n\ge1$. Every $a_n$ has positive valuation,
but $x^na_n=t^{(1,n-n^2)}$ has a strictly decreasing sequence of exponents for $n\ge1$,
so the family is not Hahn-summable.

A separate failure, allowing nonintegral coefficients, illustrates dependence on the
workspace. In $\Gamma=\Q$, the series $\sum_{n\ge0}t^{-n^2}x^n$ is not Hahn-summable at
any nonzero argument: if $w(x)=\delta\in\Q$ then the leading exponents $-n^2+n\delta$ are
unbounded below. Enlarging the workspace can change this: at the surreal argument
$x=t^\omega$ those leading exponents increase and the monomial family is Hahn-summable. A
failure in a fixed value group is not automatically a failure at every surreal scale.
\item Ordinary quaternion coefficients and Hahn-valued coefficients are different classes
and must not be interchanged. A nonzero germ with coefficients in $\HR$ has a first nonzero
coefficient $a_m$, and then $f(q)=q^m(a_m+\eps)$ with $\eps\in\mm_H$ for every nonzero
infinitesimal $q$, so $f$ does not vanish; Hahn-valued coefficients allow genuinely
displaced infinitesimal zeros, such as the zero of $q-t^\gamma\ii$.
\end{enumerate}
\end{warning}

A central translation and scaling replace $x$ by $(q-a)/r$ with $a,r\in\No$, $r\ne0$,
giving controlled infinitesimal disks even around an infinite central center. It does not
supply arbitrary expansions about noncentral centers: the order of increments and center
terms then matters and must be retained in noncommutative words.

\begin{proposition}[Slice representation, Cauchy--Riemann, twisted product]
\label{squat:prop:representation}
For an admissible series $f$, central infinitesimal $x,y$ and any $I,J\in\mathbb S(\No)$ in
the workspace,
\begin{equation}\label{squat:eq:representation-series}
 f(x+yJ)=\tfrac12(1-JI)f(x+yI)+\tfrac12(1+JI)f(x-yI),
\end{equation}
and on a fixed slice $(\partial_x+I\partial_y)f(x+yI)=0$ coefficientwise. The formal
convolution product obeys the twisted evaluation law \cref{squat:eq:twisted} whenever the
value used for conjugation is invertible, and gives zero when that value is zero.
\end{proposition}
\begin{proof}
For central $x,y$ the powers of $x+yI$ have the form $A_n(x,y)+IB_n(x,y)$ with real
integer-coefficient polynomials $A_n,B_n$ independent of $I$; summing gives
$f(x+yI)=A+IB$, and solving the two equations with $I$ and $-I$ produces
\cref{squat:eq:representation-series}. The Cauchy--Riemann and derivative identities hold
for every monomial and then by common-support regrouping. For the product formula, repeat
the proof of \cref{squat:lem:twisted}; inner conjugation leaves the valuation and modulus
of the infinitesimal argument unchanged, and a support certificate follows from the fixed
factors $f(q)^{\pm1}$ together with the supports of the original evaluated terms.
\end{proof}

\subsection{Common-domain coefficient families: two formalizations}\label{squat:sub:coherent}
For analytic continuation and contours one needs a \emph{single} ordinary domain shared by
all coefficients. The two source manuscripts set this up in two genuinely different ways,
and both are kept here because they support different theorems.

\paragraph{Formalization A: Hahn stems with an auxiliary central unit.}
Let $\Omega\subset\C$ be an ordinary connected conjugation-invariant open set meeting the
real axis, and introduce an \emph{auxiliary central} complex unit $\iota$, distinct from
every quaternion unit. An ordinary \emph{stem coefficient} is a holomorphic map
\[
 F_\gamma:\Omega\to\HR\otimes_\R\C\quad\text{with}\quad
 F_\gamma(\bar z)=\overline{F_\gamma(z)}^{\,\iota},
\]
where the overline changes $\iota$ to $-\iota$ and leaves ordinary quaternion coefficients
unchanged; the target is a complexified quaternion algebra, not the division algebra
$\SQ$. A \emph{common-domain Hahn stem} is
\begin{equation}\label{squat:eq:hahnstem}
 F(z)=\sum_{\gamma\in S}t^\gamma F_\gamma(z),
\end{equation}
with $S$ well-ordered and the same $\Omega$ for every coefficient. At a finite
$z=z_0+\eta$ with $z_0=\st(z)\in\Omega$, evaluate by
\begin{equation}\label{squat:eq:stemeval}
 F(z)=\sum_{\gamma\in S}\sum_{n\ge0}t^\gamma\frac{F_\gamma^{(n)}(z_0)}{n!}\eta^n,
\end{equation}
the factors $\eta$ being central in this complexification; the support proof of
\cref{squat:prop:evaluation} makes the double family Hahn-summable, requiring no common
bound on the magnitudes of all ordinary derivatives, since finiteness at each exponent
replaces an analytic sum over $\gamma$. Writing $F(a+\iota b)=A(a,b)+\iota B(a,b)$ with
$A,B\in\SQ$, the induced \emph{slice function} is
\begin{equation}\label{squat:eq:slicefunction}
 f(a+Ib)=A(a,b)+IB(a,b),\qquad I\in\mathbb S(\No),
\end{equation}
for finite $a,b$ with $\st(a+\iota b)\in\Omega$. The parities $A(a,-b)=A(a,b)$ and
$B(a,-b)=-B(a,b)$ make this independent of the representation by $(-I,-b)$, and at $b=0$
the vanishing of $B$ makes it independent of $I$ altogether. Values may be infinite, since
$S$ may contain negative exponents, even though the argument domain here is a finite
standard-part thickening. The Cauchy--Riemann identities for the stem are
\begin{equation}\label{squat:eq:CR}
 A_a=B_b,\qquad A_b=-B_a,
\end{equation}
so on a fixed slice $\tfrac12(\partial_a+I\partial_b)f=0$. These variable derivatives mean
ordinary coefficient derivatives followed by the same Taylor evaluation; they are not $\dH$
acting on the surreal coefficients as numbers.

\paragraph{Formalization B: coherent families over an ordinary circular quaternionic domain.}
Fix instead an ordinary circular quaternionic domain $\Omega\subseteq\HR$ and a
well-ordered $S\subseteq\Gamma$, and consider
\begin{equation}\label{squat:eq:coherent}
 f=\sum_{\gamma\in S}t^\gamma f_\gamma,\qquad
 f_\gamma:\Omega\to\HR\ \text{ordinary stem-induced slice-regular functions},
\end{equation}
with $\Omega$ common to all coefficients; slice regularity is in the stem-function sense on
a conjugation-invariant complex domain, $f_\gamma(x+Iy)=A_\gamma(x,y)+IB_\gamma(x,y)$ with
quaternion-valued components satisfying \cref{squat:eq:CR} and the even/odd parity that
makes the definition independent of the sign of $I$. For finite $x,y$ with standard part in
that domain, evaluate $A_\gamma,B_\gamma$ by their ordinary Taylor germs at the standard
point, substitute the positive Hahn corrections, and then sum over $\gamma$; the support
certificate is $S+T^*$ with $T$ the union of the supports of the finitely many scalar
increments, and the extra factor $I$ is a fixed finite quaternion contributing only its own
well-ordered support. Uniqueness of ordinary Taylor germs ensures that shrinking a
neighbourhood does not change the value.

Both formalizations are closed under the slice product induced by multiplying stem
functions, and under coefficientwise differentiation, with support sets transformed by
finite sums and sumsets. In neither is the slice product the pointwise product of
quaternion values. A family of analytic germs with no common ordinary domain is a
\emph{different space} and supports neither contour theorem below.

\begin{theorem}[Representation formula and coherent identity principle]
\label{squat:thm:representation}
For every slice function of the form \cref{squat:eq:slicefunction} and all
$I,J\in\mathbb S(\No)$,
\begin{equation}\label{squat:eq:representation}
 f(a+Jb)=\tfrac12\bigl[(1-JI)f(a+Ib)+(1+JI)f(a-Ib)\bigr].
\end{equation}
If a common-domain Hahn slice function vanishes on an ordinary nonempty real interval in
$\Omega$, then it vanishes throughout its induced domain.
\end{theorem}
\begin{proof}
Substituting $A+IB$ and $A-IB$ on the right of \cref{squat:eq:representation} gives
$A+JB$, no coefficient having been moved across $I$ or $J$. At an ordinary real point $s$
the function is $\sum_\gamma t^\gamma F_\gamma(s)$ with ordinary quaternion coefficients;
if it is zero then uniqueness of Hahn coefficients makes every $F_\gamma(s)$ zero.
Vanishing on a real interval together with the ordinary complex identity theorem, applied
to its four holomorphic coordinates, makes each $F_\gamma$ zero on connected $\Omega$, and
evaluation gives zero everywhere.
\end{proof}

This identity principle is deliberately coefficientwise. It does not claim that an
arbitrary fine-differentiable class function with many zeros is identically zero;
\cref{squat:ex:zeroderivative} already shows the danger of such a claim.

\subsection{Products, reciprocals, and zero spheres}\label{squat:sub:starfunctions}
If two stems have values $A+\iota B$ and $C+\iota D_1$, their product induces
\begin{equation}\label{squat:eq:starfunction}
 (f\star g)(a+Ib)=(AC-BD_1)+I(AD_1+BC),
\end{equation}
the slice product, which need not equal $(A+IB)(C+ID_1)$. Supports of the product lie in
the sum of two well-ordered sets with finite coefficient decompositions, so the class is
closed under it. At a point where $f(q)\ne0$ the same twist as in
\cref{squat:eq:twisted} holds:
\begin{equation}\label{squat:eq:twistedfunction}
 (f\star g)(q)=f(q)\,g\bigl(f(q)^{-1}qf(q)\bigr),
\end{equation}
since conjugating $q=a+Ib$ replaces $I$ by $I'=f(q)^{-1}If(q)$ while leaving $a,b$ fixed,
and then $f(q)(C+I'D_1)=f(q)C+If(q)D_1$ expands to \cref{squat:eq:starfunction}; if
$f(q)=0$, substitution gives $(f\star g)(q)=0$.

Quaternion-conjugating the coefficients of a stem defines $f^c$, and its normal
$f^s=f\star f^c=f^c\star f$ has central-complex stem values,
\[
 (A+\iota B)(\bar A+\iota\bar B)=N(A)-N(B)+2\iota\Rea(A\bar B),
\]
so $f^s$ is \emph{slice-preserving}: at $a+Ib$ its value lies in $\No+\No I$, not
necessarily in the center $\No$. Wherever the normal is nonzero and a reciprocal stem is
admitted locally, the slice inverse evaluates as
\begin{equation}\label{squat:eq:starinverse}
 f^{-\star}(q)=f^s(q)^{-1}f^c(q),
\end{equation}
and the displayed order is part of the formula. Local admission requires support control:
on an ordinary neighbourhood where a leading central coefficient of the normal has no
zeros, factor out that coefficient and monomial and invert the positive-order remainder by
\cref{squat:prop:neumann}. A pointwise nonzero normal is by itself no proof that one global
common-domain reciprocal family exists without shrinking or changing domains.

On one conjugacy sphere the value is $A+IB$: if $B=0$, either the entire sphere vanishes or
none of it does; if $B\ne0$, at most one point vanishes, the only candidate being
$I=-AB^{-1}$, which must also satisfy $I^2=-1$. This is the analytic counterpart of
\cref{squat:lem:sphere-remainder}, and it gives the local shape of the zero set, not a
finiteness theorem for the total number of zeros of every analytic function.

\subsection{Two coefficientwise Cauchy formulas}\label{squat:sub:cauchy}
Both manuscripts prove a coefficientwise Cauchy formula, and the two results are genuinely
different --- different kernels, different generality in the argument, different hypotheses
--- so both are printed. In each, the integral of a common-support Hahn family along an
ordinary contour is \emph{defined} by integrating each ordinary quaternion coefficient
function; this is a definition of a functional, not a fine-topological limit of quaternion
Riemann sums.

\begin{theorem}[Common-domain Cauchy formula at arbitrary slice direction]
\label{squat:thm:cauchy-general}
Fix an \emph{ordinary} $I_0\in\mathbb S(\R)$ and a bounded ordinary
conjugation-invariant domain $U$ whose closure lies in $\Omega$ and whose positively
oriented boundary is piecewise smooth; write $U_{I_0}$ for its copy in $\R+\R I_0$. Put
\begin{equation}\label{squat:eq:cauchykernel}
 K(s,q)=-\bigl(q^2-2\Rea(s)q+\abs s^2\bigr)^{-1}(q-\bar s).
\end{equation}
Then for a common-domain Hahn slice function as in Formalization~A and any finite $q$ in
the standard-part thickening of the circularization of $U$,
\begin{equation}\label{squat:eq:cauchy}
 f(q)=\frac1{2\pi}\int_{\partial U_{I_0}}^{\!H}K(s,q)\,(-I_0\,ds)\,f(s),
\end{equation}
the superscript $H$ denoting the coefficientwise construction just specified.
\end{theorem}
\begin{proof}
If the standard-part conjugacy sphere of $q$ comes from a point of $U$, the quadratic
denominator in \cref{squat:eq:cauchykernel} is nonzero for $s\in\partial U_{I_0}$, and the
symmetry of $U$ ensures that both ordinary slice representatives lie in the interior rather
than on the boundary.

First take ordinary $q\in U_{I_0}$. The commuting factors in its slice make
$K(s,q)=(s-q)^{-1}$. Each coefficient of $f$ restricted to that slice is an ordinary
left-complex-holomorphic quaternion-valued function; splitting its values into two complex
components gives the ordinary Cauchy formula with multiplication on the displayed sides,
and all ordinary derivatives of that formula hold as well, so Taylor evaluation proves it
at infinitesimally displaced points of this fixed slice.

For a fixed general finite $q$, expand the rational kernel about $q_0=\st(q)$. On the
ordinary compact boundary its zeroth-order denominator has no zero, so its reciprocal
admits a Neumann expansion in the finitely many infinitesimal coordinates of $q-q_0$ by
\cref{squat:prop:neumann}. All kernel coefficients are ordinary smooth functions on the
contour, and all supports lie in the positive monoid generated by those infinitesimal
supports; multiplying by the common support $S$ of $f(s)$ produces a well-ordered support
with finite overlap, uniformly along the contour as a support certificate. Thus the stated
coefficientwise integral exists and the coefficient identities may be integrated.

Finally $K(s,q)$ is a slice function of $q$, its quadratic factor having real coefficients
and its numerator being affine with the constant on the right. Apply
\cref{squat:eq:representation} to both sides of \cref{squat:eq:cauchy} and use the
fixed-slice result at $a\pm I_0b$ to obtain the equality at $a+Jb$ for every
$J\in\mathbb S(\No)$.
\end{proof}

\begin{theorem}[Coefficientwise Cauchy formula on a fixed slice]\label{squat:thm:cauchy-slice}
Fix an ordinary imaginary unit $I$ and an ordinary positively oriented piecewise smooth
Jordan contour $C$ whose interior and boundary lie in $\Omega\cap(\R+\R I)$, so that each
$f_\gamma$ of \cref{squat:eq:coherent} is holomorphic on a neighbourhood of that closed
interior. If $q_0$ is an ordinary interior point and $q=q_0+\eta$ with infinitesimal
$\eta\in K_\Gamma+IK_\Gamma$, then
\begin{equation}\label{squat:eq:cauchy-slice}
 f(q)=\frac1{2\pi}\int_C(s-q)^{-1}\,ds_I\,f(s),\qquad ds_I=-I\,ds,
\end{equation}
the integral and its infinitesimal kernel expansion being defined coefficientwise.
\end{theorem}
\begin{proof}
On the slice, split an ordinary quaternion-valued holomorphic function into two
complex-valued holomorphic components with a fixed perpendicular unit on the right. The
ordinary complex Cauchy formula proves \cref{squat:eq:cauchy-slice} for each $f_\gamma$ at
$q_0$ and gives all its derivatives there. Expand
$(s-q_0-\eta)^{-1}=\sum_{n\ge0}\eta^n(s-q_0)^{-n-1}$; the contour stays a positive ordinary
distance from $q_0$, so each coefficient integral is an ordinary well-defined integral
equal to the corresponding Taylor coefficient. Supports of the total family lie in
$S+\supp(\eta)^*$ and are locally finite by \cref{squat:lem:neumann}. Thus coefficientwise
integration reconstructs exactly the defined value $f(q)$.
\end{proof}

Only ordinary contours in coefficient space have been used. The compactness appealed to
concerns the real contour, not a surreal unit sphere, and neither formula is an integral
along a fine-continuous surquaternion path. The kernel
\cref{squat:eq:cauchykernel} is the classical slice Cauchy kernel \cite{GP}; the additional
work here supplies the Hahn support certificate and clarifies what the integral means over
surreal coefficients. No claim is made that testing an arbitrary surquaternion function at
ordinary points determines it: such uniqueness is a property of the specified coherent
coefficient representation, not of all functions on $\SQ$.

\begin{example}[Coefficient germs without a common domain]\label{squat:ex:nocommondomain}
A common domain cannot be removed from the hypotheses merely by saying that every
coefficient is an analytic germ. For $\eta>0$,
\begin{equation}\label{squat:eq:nocommondomain}
 F(z)=\sum_{r\ge1}\frac{t^{r\eta}}{1-rz}
\end{equation}
has a convergent germ for each coefficient, but its coefficient poles $1/r$ approach $0$:
there is no ordinary disk about $0$ on which all of them are holomorphic, and no common
small circle on which the holomorphic-inside-contour argument applies. Formal evaluation on
a suitable infinitesimal disk can still make sense --- that is a different function class,
precisely the distinction emphasized in the repository's document catalogue
\cite{RepoDocs}.
\end{example}

\subsection{Fueter regularity is a different condition}\label{squat:sec:fueter}
For quaternion coordinates $q=a+x_1\ii+x_2\jj+x_3\kk$ define the left Fueter operator and
the four-dimensional Laplacian by
\[
 \DD=\partial_a+\ii\partial_{x_1}+\jj\partial_{x_2}+\kk\partial_{x_3},
 \qquad
 \Delta_4=\partial_a^2+\partial_{x_1}^2+\partial_{x_2}^2+\partial_{x_3}^2,
\]
with coefficients multiplying on the left. Again the partial derivatives are
\emph{variable} operators on ordinary coefficient functions and their Taylor extensions,
never $\dH$. A function is left Fueter regular when $\DD f=0$. Let
$r=\sqrt{x_1^2+x_2^2+x_3^2}$ and $I=(x_1\ii+x_2\jj+x_3\kk)/r$ away from the real axis. For
an axial slice function $f=A(a,r)+IB(a,r)$ satisfying \cref{squat:eq:CR}, direct coordinate
differentiation gives
\begin{align}
 \DD f&=\Bigl(A_a-B_r-\frac{2B}{r}\Bigr)+I(B_a+A_r)=-\frac{2B}{r},
 \label{squat:eq:fueterslice}\\
 \Delta_4f&=\frac{2A_r}{r}+I\left(\frac{2B_r}{r}-\frac{2B}{r^2}\right),
 \label{squat:eq:fuetertransform}
\end{align}
so a slice-regular function need not be Fueter-regular: the coordinate function $f(q)=q$
is slice regular with $\DD f=-2$. The two theories are not interchangeable.

\begin{theorem}[Fueter transfer: polynomials and coherent families]\label{squat:thm:fueter}
For every right-coefficient slice polynomial $P$ over a real closed field,
\begin{equation}\label{squat:eq:fueterpoly}
 \DD\Delta_4P=0,
\end{equation}
and the same identity holds coefficientwise for a common-domain Hahn slice function
constructed from holomorphic stems, in either formalization of
\cref{squat:sub:coherent}:
\begin{equation}\label{squat:eq:fuetertheorem}
 \DD\Delta_4f=0 .
\end{equation}
The formula holds away from the real axis and extends across that axis wherever the stem is
defined on an ordinary neighbourhood there, the extension being coefficientwise analytic.
\end{theorem}
\begin{proof}
Set $C=2A_r/r$ and $E=2B_r/r-2B/r^2$, so that $\Delta_4f=C+IE$ by
\cref{squat:eq:fuetertransform}, using that $A$ and $B$ are harmonic in $(a,r)$. The Fueter
operator of an arbitrary axial function is
\[
 \DD(C+IE)=(C_a-E_r-2E/r)+I(E_a+C_r),
\]
and using $A_a=B_r$, $B_a=-A_r$ and their differentiated forms,
\begin{align*}
 C_a-E_r-\frac{2E}{r}
 &=\frac{2B_{rr}}r-\left(\frac{2B_{rr}}r-\frac{4B_r}{r^2}+\frac{4B}{r^3}\right)
   -\left(\frac{4B_r}{r^2}-\frac{4B}{r^3}\right)=0,\\
 E_a+C_r&=\frac{2B_{ar}}r-\frac{2B_a}{r^2}+\frac{2A_{rr}}r-\frac{2A_r}{r^2}=0 .
\end{align*}
Hence $\DD\Delta_4f=0$ for $r>0$. The parity of a stem makes $A$ even and $B$ odd in $r$,
so the apparent divisions in \cref{squat:eq:fuetertransform} have removable analytic
singularities at $r=0$; for example $B_r/r-B/r^2$ cancels the contribution of the linear
term in $B$. For a polynomial the identity extends over the real axis as a polynomial
identity, and, read as an identity in the four scalar coordinates, it consequently holds
over every real closed field, which is \cref{squat:eq:fueterpoly}. For an ordinary
slice-regular coefficient the components are real analytic, so the identity extends across
real points by analyticity. Differentiation of coefficient functions does not enlarge their
exponent support, and Taylor evaluation is Hahn-summable, so all identities extend
coefficientwise to the whole family.
\end{proof}

This is the classical Fueter transformation placed in a specified surreal setting, not a
new claim about the classical theorem \cite{GP}. Its derivation also makes the convention
visible: right-handed regularity would require the corresponding right-handed operators and
product order. For concrete checks,
\[
 \DD q=-2,\qquad \Delta_4(q^2)=-4,\qquad
 \Delta_4(q^3)=-4\bigl(3a+x_1\ii+x_2\jj+x_3\kk\bigr),
\]
and applying $\DD$ to either right-hand side gives zero; multiplying such a slice
polynomial by $t^\gamma$ or taking a common-support Hahn combination preserves the
coefficientwise theorem. The recorded suites check $\DD\Delta_4(q^n)=0$ for
$n=0,\ldots,7$, which is a genuinely quaternionic signature and not a surcomplex fact.

\subsection{What has and has not been constructed}\label{squat:sub:analysis-scope}
We now have formal infinitesimal functions, two classes of common-domain coefficient
families, their slice products, a coherent identity principle, two Cauchy functionals and a
Fueter transformation. The full surreal quaternion class does not thereby become an
ordinary Banach-space manifold. We have not asserted an unrestricted maximum principle, a
global residue theorem on arbitrary surreal paths, or the existence of a Taylor expansion
for every fine-differentiable function; the identity principle is deliberately
coefficientwise, and coefficientwise integration gives no uniqueness for arbitrary
functions on $\SQ$. For an application requiring such a statement, its function class and
its integration or support notion must be specified first.

\section{Support-controlled implicit lifting and square-root conditioning}\label{squat:sec:lifting}

This section gives the positive existence mechanism for solving perturbed quaternionic
equations. Its essential hypothesis is nonsingularity of the \emph{full real-coordinate
Jacobian}; a nonzero derivative on a single complex slice is not enough. For a scalar
vector $x=(x_1,\ldots,x_d)$ put $w(x)=\min_mw(x_m)$ with $w(0)=+\infty$, and write
$\OO_K,\mm_K$ for the valuation ring and its maximal ideal in $K_\Gamma$.

\begin{theorem}[Support-controlled implicit lifting]\label{squat:thm:lifting}
Let $P\in\OO_K[X_1,\ldots,X_d]^d$ and let $x_0\in\R^d$ satisfy
\[
 \st P(x_0)=0,\qquad
 J_0=\left(\frac{\partial\,\st P_m}{\partial X_{m'}}(x_0)\right)_{m,m'}\in\GL_d(\R).
\]
Then there is exactly one $h\in\mm_K^d$ with $P(x_0+h)=0$. Let
$S\subset\Gamma_{>0}$ be the union of the supports of the positive-order parts of the
finitely many coefficients of $P(x_0+X)$. Then
\[
 \supp(h_m)\subseteq S^*\setminus\{0\}\qquad(1\le m\le d),
\]
and if $P(x_0)\ne0$ then $w(h)=w(P(x_0))$. Uniqueness holds among \emph{all} infinitesimal
vectors. No countable convergence or cofinality assumption on $\Gamma$ is required.
\end{theorem}
\begin{proof}
Write the finitely many coefficient perturbations as $e_1,\ldots,e_{m_0}\in\mm_K$ and
replace them by independent central variables $Y_1,\ldots,Y_{m_0}$. This gives a polynomial
system $\mathcal P(X,Y)$ over $\R$ with $\mathcal P(0,0)=0$ and $X$-Jacobian $J_0$ at
$(0,0)$.

There is a unique formal solution $H(Y)\in(Y)\R[[Y]]^d$, by the following complete
recursion. If all homogeneous terms of $H$ below degree $n$ have been determined, the
degree-$n$ part of $\mathcal P(H(Y),Y)$ is $J_0H_n+R_n$, where $R_n$ depends only on the
already determined terms; set $H_n=-J_0^{-1}R_n$. Each degree involves only finitely many
monomials, which determines a formal solution and proves its formal uniqueness.

Now evaluate $Y_m=e_m$. By \cref{squat:lem:neumann} the family of all evaluated monomials
is Hahn-summable: its supports lie in $S^*$ and each exponent has only finitely many
contributing words. Thus $h=H(e_1,\ldots,e_{m_0})$ exists as a Hahn vector, has the
asserted support, and satisfies $P(x_0+h)=0$ because finite polynomial operations respect
this evaluation.

To prove uniqueness among all infinitesimal vectors, not merely among those obtained by
this substitution, let $h,h'\in\mm_K^d$. Polynomial divided differences give
\[
 P(x_0+h)-P(x_0+h')=(J_0+E)(h-h'),\qquad E\in M_d(\mm_K),
\]
and the determinant of $J_0+E$ has nonzero real constant coefficient, so the matrix is
invertible and two roots coincide.

Finally put $\varrho=P(x_0)$. The equation for $h$ is
$\varrho+J_0h+E_1h+R_{\ge2}(h)=0$ with $E_1$ of positive order and all coefficients of
$R_{\ge2}$ integral. If $h\ne0$, both error terms have valuation strictly greater than
$w(h)$; an invertible real matrix preserves vector valuation, so cancellation at the least
exponent forces $w(\varrho)=w(h)$. If $\varrho=0$, uniqueness gives $h=0$.
\end{proof}

This is a support-explicit form of nonsingular Hensel lifting. Its extra usefulness here is
that quaternion multiplication is polynomial in four scalar coordinates, so any finite
system of quaternion equations --- including expressions with coefficients inserted on both
sides of the unknowns --- becomes such a scalar system, and
\cref{squat:thm:lifting} applies as soon as its real-coordinate Jacobian is nonsingular.
This does \emph{not} give a root theorem for every noncommutative polynomial expression:
\cref{squat:warn:interspersed} still applies, and the interspersed equation
$q+\ii q\ii+1=0$ has no solution at all.

\subsection{The exact Sylvester operator for squaring}\label{squat:sub:sylvester}
Let $q_0=a+rI\ne0$ with $r>0$ and $I^2=-1$. Every increment has the unique orthogonal
decomposition
\[
 h=h_{\parallel}+h_{\perp},\qquad
 h_{\parallel}=\frac{h-IhI}{2}\in F+FI,\qquad
 h_{\perp}=\frac{h+IhI}{2},
\]
where $h_\perp$ is imaginary and perpendicular to $I$.

\begin{proposition}[Square-root linearization and condition numbers]\label{squat:prop:sylvester}
The differential of $q\mapsto q^2$ at $q_0$ is
\begin{equation}\label{squat:eq:sylvester}
 L_{q_0}(h)=q_0h+hq_0=2q_0h_{\parallel}+2a\,h_{\perp},
\end{equation}
with determinant, as an $F$-linear map on $F^4$,
\[
 \det L_{q_0}=16a^2N(q_0).
\]
For $a\ne0$ its inverse is
\begin{equation}\label{squat:eq:sylvester-inverse}
 L_{q_0}^{-1}(u)=(2q_0)^{-1}u_{\parallel}+(2a)^{-1}u_{\perp},
\end{equation}
its singular values are $2\abs{q_0}$ twice and $2\abs a$ twice, and if $a=0$ its kernel is
exactly the two-dimensional perpendicular plane.
\end{proposition}
\begin{proof}
The parallel component commutes with $I$ and the perpendicular component anticommutes with
it, which proves \cref{squat:eq:sylvester} by expansion. On the parallel two-plane, left
multiplication by $2q_0$ is a similarity of scale $2\abs{q_0}$ and determinant $4N(q_0)$;
on the perpendicular two-plane the map is multiplication by $2a$. The remaining assertions
follow.
\end{proof}

A purely imaginary square root can therefore have a perfectly nonzero \emph{slice}
derivative $2q_0$ while its four-dimensional differential has a two-dimensional kernel.
This is the linear-algebraic explanation for the spheres of square roots of a negative
central scalar in \cref{squat:prop:sqrt}, and the reason one cannot import a complex
simple-root test without modification. Note that the ``condition number'' language here is
an exact valuation-theoretic statement about $L_{q_0}^{-1}$, not a numerical-analysis
heuristic.

\begin{example}[A noncommuting displaced square root]\label{squat:ex:displacedsqrt}
Let $t$ be a positive infinitesimal monomial and $q_0=1+\ii$. There is a unique root close
to $q_0$ of $q^2=q_0^2+t\jj$, whose first terms are
\begin{equation}\label{squat:eq:sqrt-example}
 q=1+\ii+\frac{t}{2}\jj+\frac{t^2}{16}(1-\ii)-\frac{t^3}{32}\jj+\mathrm O(t^4),
\end{equation}
where $\mathrm O(t^4)$ denotes a Hahn remainder of valuation at least $4w(t)$, not an
ordinary limiting process. Substituting $q=q_0+\sum_{n\ge1}t^nh_n$ gives the exact
recursion
\[
 L_{q_0}(h_1)=\jj,\qquad
 L_{q_0}(h_n)=-\sum_{m=1}^{n-1}h_mh_{n-m}\quad(n\ge2),
\]
with support in $\N w(t)$; \cref{squat:thm:lifting} justifies the whole series rather than
only the displayed truncation.
\end{example}

\subsection{A quantified near-singular lifting region}\label{squat:sub:conditioned}
The next result keeps track of the small perpendicular singular value, and is useful even
when $q_0$ itself has no nonsingular real reduction.

\begin{theorem}[A sufficient conditioned square-root threshold]\label{squat:thm:conditioned}
Let $a\in\mm_K$ be nonzero, put $\kappa=w(a)>0$ and $q_0=a+\ii$. For $e\in D_\Gamma$ with
$\sigma=w(e)>2\kappa$ there is exactly one solution of
\[
 q^2=q_0^2+e
\]
with $w(q-q_0)>\kappa$, and writing $h=q-q_0$,
\begin{equation}\label{squat:eq:conditioned-bounds}
 w(h)\ge\sigma-\kappa,\qquad
 w\bigl(h-L_{q_0}^{-1}e\bigr)\ge2\sigma-3\kappa .
\end{equation}
The condition $\sigma>2\kappa$ is asserted as \emph{sufficient}, not as an optimal
threshold for every direction of $e$.
\end{theorem}
\begin{proof}
Write $h=az$. The equation $L_{q_0}h+h^2=e$ becomes
$z+aL_{q_0}^{-1}(z^2)=a^{-1}L_{q_0}^{-1}e$. The real-coordinate matrix of
$aL_{q_0}^{-1}$ has integral entries: its parallel block is multiplication by
$a/(2q_0)$ and its perpendicular block is $1/2$; moreover every entry of $L_{q_0}^{-1}$
has valuation at least $-\kappa$ by \cref{squat:prop:sylvester}. The right-hand side
therefore has valuation at least $\sigma-2\kappa>0$. The system in $z$ has real reduction
vanishing at $0$ with Jacobian the identity, so \cref{squat:thm:lifting} gives a unique
infinitesimal $z$ with $w(z)\ge\sigma-2\kappa$, which is the first estimate for $h$. The
exact identity $h-L_{q_0}^{-1}e=-L_{q_0}^{-1}(h^2)$ and multiplicativity of the valuation
give the second. The condition $w(h)>\kappa$ is precisely $z\in\mm_K$, which also proves
the stated uniqueness region.
\end{proof}

The theorem identifies the relevant loss as that of the \emph{inverse linearized
operator}. The determinant alone records a product of four scales and obscures which
perturbation directions are ill-conditioned: parallel errors and perpendicular errors can
behave very differently.

\section{Exact computation and formalization routes}\label{squat:sec:computation}

A surquaternion is represented by four scalar objects together with their common field or
workspace certificate, and addition, conjugation, multiplication, norm square and
inversion reduce to finitely many scalar operations. The difficulty is never Hamilton
multiplication; it is what the scalar backend can prove about surreal coefficients and
infinite supports. No finite-string format can encode every member of the proper class
$\No$, and even within a set-sized Hahn field an arbitrary support or coefficient sequence
need not be computable.

\subsection{A representation hierarchy}\label{squat:sub:hierarchy}

\begin{center}\small
\begin{tabular}{L{0.22\textwidth}L{0.34\textwidth}L{0.34\textwidth}}
\toprule
Scalar model & Exact quaternion operations & Additional requirements\\
\midrule
Rational functions in central monomials & Four-coordinate arithmetic, conjugation,
inverse, normal polynomials, Cayley and stereographic charts & Ordered monomial group and
scalar equality\\
Real algebraic extensions of that field & Modulus, square roots, finite real-closed
algebra, spectral equations & Certified algebraic roots and order\\
Effective Hahn series with support certificates & Neumann inverse, infinitesimal
exponential and logarithm, finite coefficient extraction, implicit lifting to a stated
degree & Decidable support queries and finite-overlap witnesses\\
Common-domain analytic coefficient families & Slice product, Taylor evaluation,
coefficientwise Cauchy and Fueter operators & Domain certificates and exact coefficient
analytic operations\\
Unrestricted symbolic surreal expressions & Finite denotation of some further values & No
blanket decidability or evaluation guarantee\\
\bottomrule
\end{tabular}
\end{center}

The division formula $q^{-1}=\bar q/N(q)$ should normally be the first implementation
path: it needs one central inverse, not a noncommutative linear-system solver. Positive
roots and moduli should be delegated to the scalar real-closed model, and an
implementation must not replace the $F$-valued modulus by a machine floating-point norm. A
rational Hahn model is not closed under every square root; ramification or algebraic
extension may be required.

For Hahn objects, maintain well-ordered support data, a leading exponent and an ordinary
quaternion leading coefficient; inversion first removes the leading monomial and
coefficient and then uses the positive-order geometric series
\cref{squat:eq:neumann}. For an exact infinite object, a support certificate and a
coefficient rule are part of its meaning, and a finite truncation should state the
valuation of its omitted remainder, as in \cref{squat:eq:neumannerror}. Such truncations
need not provide all coefficients below an arbitrary cutoff in finitely many steps: an
infinite support can have infinitely many exponents below that cutoff.

A series evaluator must return a value together with a summability certificate, or reject
the request as unsupported in that representation. It must not silently reinterpret
$\sum(\omega\ii)^n/n!$ as a global exponential (\cref{squat:warn:notaylor}). Likewise a
polynomial implementation must keep the $\star$ product separate from pointwise function
multiplication; \cref{squat:ex:poly} is a compact regression test for exactly that
distinction. Polynomial root computation has a precise algebraic route --- form $N_P$,
factor it over the real closed scalar field into linear and irreducible quadratic factors,
and compute the two-term remainder modulo each quadratic, which decides point versus
sphere and recovers an isolated root as $-BA^{-1}$ --- but its effectiveness depends on the
scalar field's factorization facilities. For implicit lifting, the finite formal recursion
of \cref{squat:thm:lifting} is implementable up to a chosen formal degree; if supports have
several independent scales, a formal-degree cutoff and a valuation cutoff are different
specifications, and a correct program records both the computed expression and the support
or valuation information proving what has been omitted.

\subsection{Two proposed formalization layerings}\label{squat:sub:lean}
The historical repository snapshot cited by the source manuscripts uses generic algebra
and Hahn-series infrastructure rather than a
completed surreal construction, and its README explicitly says that its generic Lean
prerequisites do \emph{not} yet construct the surreal field or the normal-form bridge
\cite{Repo}. We neither identify those prerequisites with a formalization of this article
nor infer a theorem from a successful build reported elsewhere. Mathlib already provides a
generic quaternion algebra with conjugation, norm square and division-ring instances under
ordered-field hypotheses \cite{Mathlib}; an initial formalization should instantiate and
extend those objects rather than duplicate their multiplication proofs, and in particular a
commutative complexification abstraction alone does not encode the noncommuting quaternion
units. The following are \emph{proposed layers, not claims that corresponding modules
exist or compile}. The current theorem-by-theorem status is maintained separately in
\texttt{docs/FORMALIZATION.md}; at this review all $48$ indexed statements of this
report remain pending. Both source manuscripts proposed a layering, and they differ in
granularity and in where they place the surreal bridge; both orderings are recorded.

\paragraph{Layering I, five layers, bridge last.}
\emph{Layer A: finite quaternion algebra.} Work over a generic ordered or real closed
coefficient field, formalizing the norm, inverse, centralizer, conjugacy criterion,
central-polynomial division and the normal-polynomial zero algorithm \emph{before}
introducing surreal numbers, with all statements quantified over a set-sized type.
\emph{Layer B: geometry and matrices.} Add the $F$-valued modulus, unit quaternions, the
rotation action, Hermitian diagonalization, Weyl's inequality and the projector estimate,
avoiding an ordinary real-normed-space abstraction that would incorrectly force the modulus
into $\R$. \emph{Layer C: quaternion Hahn series.} Prove the four-coordinate equivalence
$\mathbb H_{\R((t^\Gamma))}\simeq\HR((t^\Gamma))$ for a type-sized ordered abelian group
$\Gamma$, reusing the repository's support, regrouping and standard-part theorems; the new
obligations are ordered noncommutative coefficient convolution, leading-coefficient
noncancellation, the norm valuation identity, and the infinitesimal
exponential/logarithm substitution lemmas. \emph{Layer D: the actual surreal bridge.}
Supply a separately justified surreal construction and normal-form interpretation with an
explicit universe or class discipline, and only then specialize the generic results and
prove localization; generic Hahn lemmas alone do not establish this bridge, and no
set-sized type should be advertised as containing every surreal number at the same
foundational level. \emph{Layer E: analytic structures and their incompatibilities.}
Formalize common-domain stems and the coefficientwise integral using an ordinary complex
analytic library, and independently formalize the chosen surreal derivation and the
properties used in \cref{squat:thm:nooscillation}; the no-oscillation argument is short
once those foundations exist. A global trigonometric normalization belongs in another
explicitly named structure, and one must not give it an unjustified instance of a
chain-rule axiom.

\paragraph{Layering II, six layers, with lifting and coherent analysis separated.}

\begin{center}\small
\begin{tabular}{L{0.30\textwidth}L{0.62\textwidth}}
\toprule
Proposed layer & Main statements and dependencies\\
\midrule
Quaternion algebra & Conjugacy, slices, one-sided polynomial reduction, finite spectral
theory; generic ordered or real closed scalar field.\\
Quaternion Hahn series & Coefficientwise equivalence, multiplicative valuation, residue
algebra; four-coordinate transport of support lemmas.\\
Infinitesimal Lie theory & Exponential, logarithm, $\BCH$ and unit-group filtration;
Neumann finiteness for all word lengths.\\
Support-controlled lifting & Formal multivariable recursion and its evaluation;
real-coordinate Jacobian certificate.\\
Coherent slice analysis & Common ordinary domains, coefficientwise contour operations,
representation and Fueter identities.\\
Full surreal bridge & Class or universe discipline, normal forms, scalar exponential,
finite part, and the chosen global phase.\\
\bottomrule
\end{tabular}
\end{center}

\noindent The first five layers of Layering~II can be meaningfully proved in set-sized
generic workspaces; the last is not obtained by declaring a small type whose elements are
all surreal numbers. Both layerings agree on the essential point: the surreal bridge comes
\emph{after} the generic Hahn work, and the ordered magnitude calls for care, since the
usual real-valued normed-space interfaces are not interchangeable with an $F$-valued
magnitude.

\subsection{Scope of the executable checks}\label{squat:sub:checks}
Two independent exact symbolic suites accompany this report, preserved unchanged in
\texttt{code/} with their recorded outputs in \texttt{data/}; see
\cref{squat:app:reproduce} for the inventory and for their independently reproduced
results. Both use exact SymPy arithmetic on indeterminates and rational numbers, with no
floating-point sampling and no network access.

The first suite checks the Hamilton table, generic associativity and norm identities, the
complex matrix representation, Rodrigues' formula, the derivation reconstruction formula
\cref{squat:eq:innerrep}, the polynomial example \cref{squat:ex:poly}, a Cayley unit,
truncated noncommutative exponential, logarithm and $\BCH$ identities, the two-scale
matrix of \cref{squat:ex:spectral} through order $t^9$, and the Fueter polynomial
identities. Its recorded run passed all $69$ check groups, comprising $252$ scalar
identities. The second suite checks generic associativity in twelve scalar variables,
conjugation reversing a generic product, norm-square multiplicativity, the two-sided
inverse, the commutator as twice the cross product, matrix representation
multiplicativity, determinant and adjoint, unnormalized Rodrigues, the stereographic
chart's exact unit norm, the full square-map Jacobian determinant of
\cref{squat:prop:sylvester}, the parallel/perpendicular inverse of the square
linearization, the root asymmetry of \cref{squat:ex:poly} with its central norm polynomial,
the interspersed no-root example of \cref{squat:warn:interspersed}, $\log(\exp x)=x$ and
$\exp(\log(1+x))=1+x$ through degree six, $\BCH$ through degree four, the adjoint series
through degree five, the graded group commutator leading term $2t^3\kk$, the displaced
square root \cref{squat:eq:sqrt-example} through degree seven, the radial exponential
differential determinant \cref{squat:eq:exp-jacobian}, and
$\DD\Delta_4(q^n)=0$ for $n=0,\ldots,7$. Its recorded run passed all $31$ named checks.

These checks detect algebraic sign errors, misplaced factors and expansion mistakes in the
displayed examples. They do \emph{not} verify set or class foundations, the Hahn support
lemma, real-closed-field transfer, the analytic hypotheses, or the Berarducci--Mantova
construction, and they do not establish any general analytic or spectral theorem. The
general proofs remain the mathematical arguments in this article, relying on the explicitly
cited foundational results.

\section{Limitations and non-claims}\label{squat:sec:nonclaims}

Every honest limitation recorded by either source manuscript is reproduced here. They are
numbered for reference and none of them is weakened.

\begin{enumerate}[label=\textup{(N\arabic*)},leftmargin=3.2em]
\item\label{squat:nc:refereed} This is not a refereed article, not a priority claim, and
not a claim to solve a named published open problem. Of
\cref{squat:thm:nooscillation,squat:cor:nochain} in particular: we make no claim that this
formulation has never appeared in the literature.
\item\label{squat:nc:literature} Literature inspection was targeted, not an exhaustive
priority search. Classical quaternion results and existing scalar-surreal constructions are
attributed rather than presented as new, and no first-publication or open-conjecture claim
is made anywhere.
\item\label{squat:nc:classical} Much of
\cref{squat:sec:algebra,squat:sec:slices,squat:sec:polynomials} is acknowledged to be
classical quaternion algebra over a real closed field \cite{Voight}, written out to expose
hypotheses; the slice framework and the Fueter transformation have classical antecedents
\cite{GP}; and the scalar phase normalization is due to Ehrlich and Kaplan \cite{EK}. The
genuinely surreal content is the class/set boundary, support control, multiscale estimates,
and the incompatibility results.
\item\label{squat:nc:nolean} No Lean implementation or formal verification of this
article is claimed. The archives contain no purported Lean proof; both layerings in
\cref{squat:sub:lean} are proposed modules only, and we refuse to infer a theorem from a
successful build reported elsewhere.
\item\label{squat:nc:checks1} The $69$ symbolic check groups and $252$ scalar identities of
the first suite verify finite polynomial and rational identities and truncated formal
expansions only. They do not construct the surreals, do not prove the Hahn support lemma,
do not verify real-closed-field transfer, do not verify the Berarducci--Mantova
construction, and do not establish any general analytic or spectral theorem.
\item\label{squat:nc:checks2} The $31$ checks of the second suite are finite identity and
truncation checks only --- not a verification of theorems about proper classes, arbitrary
real closed fields, or infinite well-ordered supports. Its truncation order is an explicit
finite parameter, and no finite truncation is presented as a proof of an infinite-series
identity.
\item\label{squat:nc:nophysics} Neither source manuscript makes any claim about any
empirical subject, and none is made here. There is no assertion about spacetime,
relativity, quantum mechanics, gauge theory, renormalization, measurement, or physical
spin. The physics-adjacent vocabulary that does appear is strictly mathematical: ``spin''
refers to the algebraic double cover of \cref{squat:eq:spin3}; ``rotation'' always means an
element of $\SO(3,F)$ or $\SO(4,F)$ acting on $\Ima\HF$, never a physical motion; the
``avoided crossing'' of \cref{squat:ex:crossing} is finite Hermitian $2\times2$ linear
algebra; the ``angular velocity'' identity \cref{squat:eq:angular-velocity} is a derivation
identity in the surreal coefficients and expressly not a time evolution or an existence
theorem for trajectories; and ``condition number'' in
\cref{squat:prop:sylvester,squat:thm:conditioned} is an exact valuation-theoretic statement,
not a numerical heuristic.
\item\label{squat:nc:covering} The two-to-one map $S^3(F)\to\SO(3,F)$ is stated only as a
map on algebraic group points, \emph{not} as a covering-space theorem for the fine
topology. Neither presentation of the rational chart in \cref{squat:prop:charts}
claims that $\Sp(1,F)$ is topologically a real three-sphere, and the polar decomposition
\cref{squat:eq:polargroup} is algebraic and does not assert compactness of $S^3(\No)$.
\item\label{squat:nc:onesided} The root theorem covers \emph{only} one-sided polynomials
with a central formal indeterminate and coefficients on the right. Arbitrary
interspersed-coefficient expressions can have no root at all, by both counterexamples of
\cref{squat:warn:interspersed}; $\SQ$ is a noncommutative division algebra, not an
algebraically closed field.
\item\label{squat:nc:finitespectral} All spectral theory here is finite-dimensional matrix
algebra. No Hilbert-space, infinite-dimensional or proper-class spectral theorem is
asserted, and no compactness of a surreal unit sphere is used.
\item\label{squat:nc:nodeterminant} The determinantal singular-scale theorem and the
positive Cauchy--Binet argument of the companion spectral-theory report \cite{RepoSpec} are
commutative determinant arguments and cannot be transferred unchanged to quaternionic
matrices; no quaternionic determinant version is claimed here
(\cref{squat:warn:nodet}). Commutative results are used only with their scalar or
slice hypotheses established.
\item\label{squat:nc:matrixval} For matrices only $w(AB)\ge w(A)+w(B)$ holds. The equality
\cref{squat:eq:valuation-product}, valid for individual nonzero quaternions, must not be
copied to matrices.
\item\label{squat:nc:cluster} The cluster and block multiscale spectral calculus is a
\emph{direction}, not a theorem: only one rank-one instance is proved
(\cref{squat:thm:projector,squat:ex:spectral}), with no optimal block constant.
\item\label{squat:nc:expmnotiso} $\Expm$ gives a \emph{bijection} between the imaginary
infinitesimals and $U_1$, \emph{not} an isomorphism of groups: the transported group law is
noncommutative and is $\BCH$.
\item\label{squat:nc:globalexp} The global radial exponential depends on a \emph{chosen}
scalar normalization, that of Ehrlich and Kaplan; it is not an everywhere Hahn-summed power
series and not an additive-to-multiplicative homomorphism on the whole additive algebra, so
its fibre over $1$ is not a kernel. There is no preferred logarithm: on the negative
central axis every $\log_{\No}\rho+\pi I$ works, and ordinary $2\pi$ periods give further
logarithms elsewhere. The fibre computations of \cref{squat:thm:logfibres} do not settle the
robustness questions for alternative surcomplex exponentials posed in \cite{EK}.
\item\label{squat:nc:tube} The tube exponential $\Expt$ is not a Hahn sum of $q^n/n!$ and
is not $\exp(q_0)\Expm(\eps)$ in general (\cref{squat:warn:notaylor}), and its domain
excludes $\omega\ii$.
\item\label{squat:nc:chainrule} \Cref{squat:cor:nochain} forbids a global commuting
exponential chain rule for the componentwise Berarducci--Mantova derivation $\dH$ only. It
is not a statement about the multiplicative behaviour of any exponential, it does not
conflict with the Liouville-closed real differential-field results of \cite{BM}, and it does
not contradict the commutative surcomplex homomorphism claim of \cite{RepoPhys}; see
\cref{squat:sub:consistency,squat:rem:oscillation-relations}.
\item\label{squat:nc:bmsurjective} Surjectivity of $\dH$ is not a solvability theorem for
nonlinear quaternionic differential equations; it gives coefficientwise antiderivatives
only.
\item\label{squat:nc:workspaceclosure} A set-sized Hahn workspace need not be closed under
Gonshor's exponential or under a chosen derivation. Such closure is an additional
hypothesis wherever those operations are used, and localization
(\cref{squat:thm:localization}) does not supply it, nor does it identify a workspace's
intrinsic topology with its subspace topology in $\No$.
\item\label{squat:nc:summability} Positive valuation is a summability certificate, not
cofinality: a valid Hahn identity need not be the limit of its partial sums in any
valuation topology (\cref{squat:warn:cofinality}). The word ``convergent'' is never used as
a substitute for a named construction.
\item\label{squat:nc:fine} Fine differentiation does not produce ordinary continuous paths
and is explicitly not used to justify a general inverse-function theorem;
\cref{squat:thm:lifting} is the existence tool. A zero fine differential does not imply
constancy (\cref{squat:ex:zeroderivative}), and the full class does not become a
Banach-space manifold.
\item\label{squat:nc:contours} Coefficientwise contour integration is over \emph{ordinary}
contours in an ordinary common domain, not over fine-continuous surreal contours, and it
gives no uniqueness for arbitrary functions on $\SQ$. No unrestricted maximum principle and
no global residue theorem on arbitrary surreal paths is asserted, and the identity
principle is deliberately coefficientwise.
\item\label{squat:nc:commonsupport} Admissible-series and coherent-function statements
require a \emph{common} support or a \emph{common} ordinary domain. Knowing
$w(a_n)\ge0$ pointwise is not a substitute, a family with no common domain is a different
space (\cref{squat:ex:nocommondomain}), and ordinary-quaternion and Hahn-valued coefficient
classes must not be interchanged (\cref{squat:warn:commonsupport}).
\item\label{squat:nc:threshold} The $\sigma>2\kappa$ conditioned square-root threshold of
\cref{squat:thm:conditioned} is asserted as sufficient only, not as optimal for every
perturbation direction.
\item\label{squat:nc:effective} All algorithms are effective \emph{relative to} an exact
scalar oracle. Arbitrary surreal coefficients have no finite computable encoding; root
finding needs factorization over the real closed scalar field; infinite supports can have
infinitely many exponents below a cutoff, so truncation is not generally finitely
computable; and the unrestricted-symbolic-expression row of
\cref{squat:sub:hierarchy} carries no decidability or evaluation guarantee.
\item\label{squat:nc:octonions} The Cayley--Dickson construction over $\No$ also produces
an octonion algebra with positive sum-of-eight-squares norm and hence no zero divisors, but
it is \emph{alternative} rather than associative, and the centralizer, linear-factor and
matrix proofs in this article cannot be copied without revisiting their uses of
associativity. Later Cayley--Dickson algebras already have zero divisors over $\R$
\cite{Baez}, and those examples persist under scalar extension to $\No$; infinite
coefficients do not repair an algebraic zero-divisor identity. This is an observation,
not an octonionic theory.
\item\label{squat:nc:directionsopen} The research directions of
\cref{squat:sub:directions} are explicitly open. None of them is claimed solved by the
constructions above.
\item\label{squat:nc:sha} The first source archive's README lists a file
\texttt{SHA256SUMS.txt} that is \emph{absent} from the delivered tree. The checksums it
would have contained are therefore unavailable, and no integrity claim is made on their
basis.
\item\label{squat:nc:selfwriting} Each accompanying script writes its own result file
beside itself by default, so re-running it in place destroys the delivered record. The
independent reproduction described in \cref{squat:app:reproduce} was therefore performed on
copies; a byte-identity check after an in-place run would compare two equally modified
copies and prove nothing.
\end{enumerate}

\section{Research directions}\label{squat:sub:directions}

Six directions are left open, three from each source manuscript. None is claimed solved.

\emph{Collision-stable spectral subspaces.} The rank-one estimate
\cref{squat:thm:projector} can be developed into cluster-projector estimates with a gap
between two eigenvalue groups and no gap within either group. The natural quaternionic
proof solves an off-diagonal Sylvester equation after spectral diagonalization and
estimates denominators by the intercluster gap; iterating across a hierarchy of valuations
should produce a block-diagonal multiscale spectral calculus. \Cref{squat:ex:spectral}
proves one rank-one instance, not that iterative theorem
(\labelcref{squat:nc:cluster}).

\emph{Support-aware eigenspace perturbation through collisions.} Repeated quaternionic
eigenvalues call for a support-aware perturbation theory of eigenspaces and spectral
projections. \Cref{squat:thm:spectral} supplies existence but not a uniform series for an
arbitrarily chosen eigenbasis through a collision.

\emph{Finite analytic zero algebras over quaternion coefficients.} The sphere remainder
algorithm completely solves finite-degree one-sided polynomials. For common-domain analytic
stems, a useful next target is a support-certified local preparation theorem tracking
isolated and spherical zeros under infinitesimal perturbation. Such a theorem must specify
whether it takes place in a fixed-domain ring, a common-domain germ ring, or a ring of
unrelated coefficient germs: \cref{squat:ex:nocommondomain} shows why conflating these
rings can invalidate the proof before noncommutativity even enters.

\emph{Deformation of zero classes from spheres to points.} Quaternionic zero classes can
change from whole spheres to isolated points under small coefficient perturbations. A
deformation theory should retain the entire finite conjugacy-class algebra, not merely a
list of central roots; the explicit remainder $XA+B$ of
\cref{squat:lem:sphere-remainder} and the singular Sylvester operator of
\cref{squat:prop:sylvester} identify the data such a theory must control.

\emph{Differential extensions containing oscillations.}
\Cref{squat:thm:nooscillation} is a definite obstruction rather than a vague warning:
constant-frequency rotations are absent from the componentwise Berarducci--Mantova
quaternionic differential algebra. One concrete direction is to adjoin solutions of
$\dH y=Iy$ with a controlled involution and norm and study what replaces the finite-element
smallness property; another is to retain the surquaternions and choose a different
derivation, making its constants and order compatibility explicit; a third is to use only
variable derivatives of common-domain stems and never identify them with a field
derivation. These are different mathematical choices, not interchangeable notations for a
single calculus.

\emph{Comparison of alternative global radial phases.} A comparison of global radial phase
extensions should separate their shared infinitesimal Lie theory from their infinite-radius
critical sets. The character family of \cite{RepoTrig,EK} makes clear that
\cref{squat:thm:logfibres,squat:thm:exp-derivative} describe one chosen normalization;
which of their features are robust across the family is open.

\subsection{The resulting picture}\label{squat:sub:synthesis}
Surquaternions form an associative division algebra with an ordered-field-valued modulus,
algebraically closed slices, exact rotation geometry at every scale, and a robust finite
spectral and polynomial theory. Their normal forms carry the \emph{noncommutative} residue
algebra $\HR$, while the value group records genuinely surreal scales. Positive-support
calculus supplies precise infinitesimal exponentials, logarithms, perturbation expansions
and implicit lifting, and ordinary coefficient analysis supplies coherent slice and Fueter
operations.

The structures fit together provided their boundaries remain visible. The global phase of
an infinite imaginary argument is not determined by a raw power series. Hahn summation is
not a universal topological limit. The field derivation is not differentiation of an
analytic variable, and neither is the fine differential or the Fueter operator. The
no-oscillation theorem proves that at least one tempting identification of these structures
is impossible, while the constructive results show how much useful theory remains once that
identification is abandoned --- and how much of it must be reproved rather than imported,
because the algebra does not commute.

\appendix
\section{Notation ledger and conventions}\label{squat:app:ledger}

This table is the ledger of \cref{squat:sec:ledger} in full, extended to cover every
object either source manuscript introduced.

\begingroup
\setlength{\parskip}{0pt}
\begin{longtable}{L{0.21\textwidth}L{0.71\textwidth}}
\toprule
Notation or construction & Exact convention\\
\midrule
\endhead
\midrule
\endfoot
\bottomrule
\endlastfoot
$F$, $\No$, $\HR$ & A real closed ordered field; the surreal proper class; the ordinary
quaternions over $\R$.\\
$\HF$, $\SQ$ & Hamilton algebra over $F$, and over $\No$. Scalar coefficients are
central.\\
$\bar q$, $N(q)$, $\abs q$ & Quaternion conjugation (an anti-involution), reduced norm
$q\bar q$, and its positive square root in the ordered center. Not real-valued and not
floating point.\\
$\Rea q$, $\Ima q$ & Central scalar part and pure imaginary part; not maps to ordinary real
coefficients unless the input is ordinary.\\
$\mathbb S(F)$ & Imaginary square roots of $-1$; $F_I=F+FI$ is a complex slice,
algebraically closed for real closed $F$.\\
$t^\gamma$ & Conway monomial $\omega^{-\gamma}$, \emph{not} an unqualified exponential
power and not $\exp(-\gamma\log\omega)$.\\
$K_\Gamma$, $D_\Gamma$ & $\R((t^\Gamma))$ and $\mathbb H_{K_\Gamma}\simeq\HR((t^\Gamma))$,
for a divisible set-sized ordered subgroup $\Gamma\subseteq(\No,+)$.\\
$\supp(q)$, $w(q)$, $\lc(q)$ & Well-ordered set of increasing exponents; least exponent,
with $w(0)=+\infty$; leading ordinary quaternion coefficient. Source~11 writes $\nu$ for
$w$.\\
$\OO_H$, $\mm_H$, $\st$ & Finite coordinates, infinitesimal coordinates (a two-sided
maximal ideal), and coordinate-zero extraction onto the noncommutative $\HR$. In a
workspace these are written $\OO_D,\mm_D$ and $\OO_K,\mm_K$ for scalars.\\
$S^*$ & All finite sums from a positive well-ordered support $S$, including $0$.\\
Hahn-summable & Set-indexed, well-ordered union of supports, each exponent in finitely many
members; summed coefficientwise. Never a limit of partial sums.\\
$P\star Q$, $P^c$, $N_P$ & Product with a central indeterminate and right-coefficient
evaluation, \emph{not} pointwise multiplication; coefficient conjugation; central
symmetrization $P\star P^c$, also written $P^s$ in source~11.\\
$\Delta_{a,r}$, $[a,r]$ & $X^2-2aX+a^2+r^2$; the conjugacy sphere $a+r\mathbb S(F)$.\\
$S^3(F)$, $\Sp(1,F)$, $U$, $U_1$ & The modulus-one group; the finite units of modulus one;
its principal subgroup $U\cap(1+\mm_H)$.\\
$\Expm$, $\Logm$ & Infinitesimal Hahn exponential and logarithm at positive valuation;
$\Exp_{\mm}$ in source~10 and $\exp_{\mathrm H}$ in source~11.\\
$\Expt$ & Finite-angle tube exponential on $\mathcal T$: Gonshor scalar exponential times
canonical finite angle; finite imaginary part, unrestricted scalar part; $\Exp_0$ in
source~10.\\
$\exp_{\No}$, $\log_{\No}$ & Gonshor's scalar exponential on $\No$ and its inverse.\\
$\fin x$, $\Oz$ & Real-plus-infinitesimal part of the scalar normal form (additive, not
multiplicative, $\fin\omega=0$); the omnific integers $\No_{\mathrm{PI}}+\Z$.\\
$\Sin$, $\Cos$ & Global integer-part-normalized phase, $\sin(\fin x)$ and $\cos(\fin x)$;
distinct from the finite-angle $\sin,\cos$ of \cref{squat:eq:finite-trig}.\\
$\Expr$ & Global radial quaternion exponential \cref{squat:eq:globalexp} using the
Ehrlich--Kaplan normalization; $\Exp_{\mathrm{EK}}$ in source~10 and
$\Exp_{\mathrm{can}}$ in source~11.\\
$\theta$, $\eta$, $B$ & Finite angle; infinitesimal phase residue; accumulated phase. The
three senses fixed in \cite{RepoDiff}; no fourth sense is used.\\
$\partial$, $\dH$ & Berarducci--Mantova scalar derivation with $\partial\omega=1$; its
componentwise extension to $\SQ$ ($D$ in source~10, $\partial_H$ in source~11).\\
$\delta$, $\delH$, $\ad_w$ & An arbitrary scalar field derivation; its componentwise
extension ($D_0$ in source~10); the inner part of a general derivation.\\
$\dsl$ & Coefficientwise slice (formal) derivative of a right-coefficient series; not the
full differential.\\
$df_q$ & Fine ($\No$-linear Fr\'echet) differential of
\cref{squat:def:fine-derivative}; not a Banach derivative.\\
$\partial_a,\partial_b,\partial_r$ & Variable derivatives of ordinary coefficient functions
and their Taylor extensions; never $\dH$.\\
$\iota$ & Auxiliary central complex unit for stems; distinct from every quaternion unit.\\
$\DD$, $\Delta_4$ & Left Fueter operator and four-coordinate variable Laplacian.\\
$\int^H$ & Ordinary contour integration of individual Hahn coefficient functions, with a
uniform support certificate; not a fine-topological limit of Riemann sums.\\
$L_{q_0}$ & The Sylvester operator $h\mapsto q_0h+hq_0$ of
\cref{squat:prop:sylvester}.\\
Fine topology & Balls of every positive surreal radius on the full class; distinct from a
workspace's intrinsic valuation topology and from the coarse standard-part topology
(\cref{squat:rem:threetopologies}).\\
\end{longtable}
\endgroup

\section{Dependencies beyond the displayed proofs}\label{squat:app:dependencies}

\begin{center}\small
\begin{tabular}{L{0.36\textwidth}L{0.54\textwidth}}
\toprule
Results & Main dependencies beyond their displayed proofs\\
\midrule
Division algebra, centralizers, conjugacy & Ordered-field algebra; positive square roots
where normalization is used.\\
Polynomial root algorithm, spectral theorem & Real closedness of $F$ and algebraic
closedness of $F[\ii]$.\\
Surreal realization and localization & Conway normal-form theorem and the real-closed
Hahn-field theorem \cite{Conway,Alling,EK}.\\
Hahn inverse, principal logarithm, $\BCH$, adjoint series, implicit lifting &
Positive-support Neumann lemma (\cref{squat:lem:neumann}, via Higman's finite-word lemma)
and coefficientwise strong regrouping.\\
Finite angles and global quaternion exponential & Gonshor's real exponential and the
Ehrlich--Kaplan global circular character \cite{Gonshor,EK}; the scalar statements are
inherited from \cite{RepoTrig}.\\
No-oscillation theorem and chain-rule corollary & Berarducci--Mantova constant field,
smallness, and $\partial\omega=1$ \cite{BM}; the scalar precursor is inherited from
\cite{RepoDiff}.\\
Cauchy and Fueter results & Ordinary coefficient analysis on one common domain, plus Hahn
Taylor support certificates; classical slice theory \cite{GP}.\\
Fixed-dimension algebraic transfer & First-order real-closed-field transfer, subject to
expressibility \cite{Swan}.\\
\bottomrule
\end{tabular}
\end{center}

\section{Worked exact calculations for reuse}\label{squat:app:calculations}

\subsection{Conjugating two imaginary units}
For $I=\ii$ and $J=\jj$, the conjugating element of \cref{squat:thm:conjugacy} is
$h=1-\jj\ii=1+\kk$. Normalization gives $u=(1+\kk)/\sqrt2$, and
\[
 u\ii u^{-1}=\jj,\qquad u\jj u^{-1}=-\ii,\qquad u\kk u^{-1}=\kk .
\]
Replacing the constant direction by a surreal one does not change the formula: for example
$J=(\ii+t\jj)/\sqrt{1+t^2}$ produces an exact infinitesimal rotation through the same
construction, with no numerically approximated angle.

\subsection{An inverse and its error certificate}
For $x=t\ii$ one has $x^2=-t^2$, so
\[
 (1+x)^{-1}=\frac{1-t\ii}{1+t^2}=1-t\ii-t^2+t^3\ii+t^4-\cdots .
\]
The difference from the first $M+1$ powers $\sum_{n=0}^M(-x)^n$ is exactly
$(-x)^{M+1}(1+x)^{-1}$, of valuation $M+1$. This is an equality and a certificate, even in
the full fine topology where the displayed partial sums do not converge as an ordinary
sequence (\cref{squat:warn:cofinality}).

\subsection{The first noncommutative correction}
For $x=t\ii$ and $y=t\jj$ the degree-two term of $\BCH(x,y)$ is $t^2\kk$, and its
degree-three term is
\[
 \tfrac1{12}\bigl([t\ii,[t\ii,t\jj]]+[t\jj,[t\jj,t\ii]]\bigr)=-\tfrac{t^3}{3}(\ii+\jj),
\]
so
\[
 \BCH(t\ii,t\jj)=t(\ii+\jj)+t^2\kk-\frac{t^3}{3}(\ii+\jj)
 +\text{terms of valuation at least }4 .
\]
A scalar exponential implementation that simply adds the exponents would lose the entire
$t^2\kk$ rotation. Both verification suites check this coefficient calculation by truncated
noncommutative multiplication, independently of the displayed bracket simplification; see
also the unequal-scale form in \cref{squat:ex:unequalscales}.

\subsection{The two-scale spectrum at integer exponents}
For $\alpha=1$, $\beta=2$, \cref{squat:eq:twoeigen} specializes to
\[
 \lambda_-=\frac{t-t\sqrt{1+4t^2}}2=-t^3+t^5-2t^7+5t^9+\cdots,
 \qquad
 \lambda_+=t+t^3-t^5+2t^7-5t^9+\cdots,
\]
and the off-diagonal entry of its lower projector is
\[
 -\frac{t\ii}{\sqrt{1+4t^2}}=-t\ii+2t^3\ii-6t^5\ii+20t^7\ii+\cdots .
\]
The ordinary standard-part matrix is zero, yet a meaningful eigenspace is resolved at order
$t$, with its perturbation explicitly visible. This is one reason ordinary floating-point
substitution at a small positive real $t$ should be treated as an illustration rather than
as the definition of the surreal problem.

\subsection{The radial exponential at an infinite radius}
The two computations to keep side by side are
\[
 \Expr(\omega\ii+\pi\jj)=1+\frac{\pi^2}{2\omega}\ii+\mathrm O(\omega^{-2}),
 \qquad
 \Expr(\omega\ii)\,\Expr(\pi\jj)=-1,
\]
from \cref{squat:ex:radialnotproduct}, together with the critical-point expansion
\[
 \Expr(\omega I+\eta J)=1+\frac{\eta^2}{2\omega}I+\frac{\eta^3}{2\omega^2}J+\cdots
 \qquad(J\perp I)
\]
of \cref{squat:ex:criticalsphere}, whose vanishing first-order transverse term is the
rank-two degeneracy of \cref{squat:eq:exp-jacobian} at $r=\omega\in\pi\Oz$. The recorded
suites check the determinant formula symbolically.

\subsection{The displaced square root and its recursion}
\Cref{squat:eq:sqrt-example} gives
$q=1+\ii+\tfrac t2\jj+\tfrac{t^2}{16}(1-\ii)-\tfrac{t^3}{32}\jj+\mathrm O(t^4)$ as the
unique root of $q^2=(1+\ii)^2+t\jj$ near $1+\ii$, obtained from the exact Sylvester
recursion of \cref{squat:sub:sylvester}. The coefficients $\jj/2$, $(1-\ii)/16$ and
$-\jj/32$ are checked symbolically through degree seven.

\section{Reproducibility, provenance, and archive inventory}\label{squat:app:reproduce}

This source is a standalone \LaTeX{} article with an embedded bibliography, requiring no
external figures, bibliography database or proprietary fonts. Build it with
\begin{verbatim}
latexmk -pdf -interaction=nonstopmode article.tex
\end{verbatim}
and inspect \texttt{article.log} for errors, undefined references, undefined citations,
multiply-defined labels and duplicate destinations before accepting the output.

\paragraph{Source manuscripts.} This report is the merge of two independently written
manuscripts, \emph{noncommutative-hahn-\allowbreak obstructions} and
\emph{radial-exponential-\allowbreak support-lifting}, neither of which is reproduced
here. The first was delivered as a $34$-page article with $27$ numbered statements; the
second as a $32$-page article. The disposition of every result either of them reached is
recorded in the provenance material that remains: the scope and provenance of
\cref{squat:sec:scope}, the ledger of \cref{squat:sec:ledger}, and the limitations of
\cref{squat:sec:nonclaims}.

\paragraph{Code and data.} Two independent suites are preserved unchanged in
\texttt{code/}, with their recorded outputs and metadata in \texttt{data/}. Both pin
\texttt{sympy==1.14.0}. The first suite's record states $69$ of $69$ check groups passed,
$0$ failed, $252$ scalar identities, with no failing entry, under Python $3.13.5$ and SymPy
$1.14.0$; its build metadata records $34$ pages, $27$ theorem-type statements, $27$ proof
environments, zero \LaTeX{} errors, warnings, overfull or underfull boxes, zero unresolved
reference markers, and \texttt{formal\_theorem\_proof\_verification: false}. The second
suite's record states $31$ of $31$ checks passed with \texttt{all\_passed: true}, under the
same versions, and describes itself as ``exact finite algebra and truncated-series checks,
not formal theorem verification''.

\paragraph{Independent reproduction.} Both suites were re-run independently on \emph{copies}
of the trees, under Python $3.14.4$ with SymPy $1.14.0$, and both exited zero, reporting
$252$ identities in $69$ groups and $31$ checks respectively. Running on copies is not a
formality: each script writes its own result file beside itself by default, so an in-place
re-run overwrites the delivered evidence, and a byte-identity check performed afterwards
would compare two equally modified copies and prove nothing
(\labelcref{squat:nc:selfwriting}). Nothing under \texttt{code/} or
\texttt{data/} was modified in producing this report.

\paragraph{Provenance and status.} Both source manuscripts recorded the same repository
snapshot, \texttt{39f2be6667ade51bca2b45daa47e289d69c09764}, inspected 21 September 2026,
as orientation only \cite{Repo,RepoDocs,RepoTrig}; the first records
\texttt{repository\_modified: false} and \texttt{repository\_build\_run: false} in its
manifest. Neither treated the repository's build as a formal verification of its theorems,
and neither claimed to have added or compiled any Lean module. Both are AI-assisted
research drafts with written proofs and finite symbolic checks; neither is refereed or
machine-checked, and this merged report inherits exactly that status. Every limitation of
either manuscript is reproduced in \cref{squat:sec:nonclaims}.

\paragraph{Reading routes.} For an algebra-first reading follow
\cref{squat:sec:algebra,squat:sec:slices,squat:sec:polynomials,squat:sec:hahn} and then
\cref{squat:sec:computation}. For infinitesimal geometry follow
\cref{squat:sec:hahn,squat:sec:topology,squat:sec:explog,squat:sec:angles}. For the
exponential at infinite arguments read
\cref{squat:sec:ledger,squat:sec:angles,squat:sec:exponentials,squat:sec:differential}
together: omitting the scalar normalization of \cref{squat:sub:ekphase} changes the
subject. For analysis and formalization read
\cref{squat:sec:topology,squat:sec:analysis,squat:sec:lifting,squat:sec:computation} with
the four operators of \cref{squat:conv:derivs} kept apart.

\begingroup
\small
\raggedright
\begin{thebibliography}{99}
\setlength{\itemsep}{0.35em}

\bibitem{Conway}
J. H. Conway, \emph{On Numbers and Games}, second edition, A K Peters, 2001
(first edition, Academic Press, 1976).

\bibitem{Gonshor}
H. Gonshor, \emph{An Introduction to the Theory of Surreal Numbers}, London Mathematical
Society Lecture Note Series 110, Cambridge University Press, 1986.
\href{https://doi.org/10.1017/CBO9780511629143}{doi:10.1017/CBO9780511629143}.

\bibitem{Alling}
N. L. Alling, \emph{Foundations of Analysis over Surreal Number Fields}, North-Holland,
1987.

\bibitem{vdDE}
L. van den Dries and P. Ehrlich, Fields of surreal numbers and exponentiation,
\emph{Fundamenta Mathematicae} \textbf{167} (2001), 173--188;
erratum, \textbf{168} (2001), 295--297.

\bibitem{BM}
A. Berarducci and V. Mantova, Surreal numbers, derivations and transseries,
\emph{Journal of the European Mathematical Society} \textbf{20} (2018), 339--390.
\href{https://doi.org/10.4171/JEMS/769}{doi:10.4171/JEMS/769}.
Preprint: \href{https://arxiv.org/abs/1503.00315}{arXiv:1503.00315}, especially
Theorems A and B.

\bibitem{EK}
P. Ehrlich and E. Kaplan, Surreal ordered exponential fields,
\emph{The Journal of Symbolic Logic} \textbf{86} (2021), no.~3, 1066--1115.
Preprint: \href{https://arxiv.org/abs/2002.07739}{arXiv:2002.07739}, especially
Sections 2, 3 and 11.

\bibitem{Voight}
J. Voight, \emph{Quaternion Algebras}, Graduate Texts in Mathematics 288, Springer, 2021.
Open-access book and corrections: \url{https://jvoight.github.io/quat.html}.

\bibitem{GP}
R. Ghiloni and A. Perotti, Slice regular functions on real alternative algebras,
\emph{Advances in Mathematics} \textbf{226} (2011), 1662--1691.
\href{https://arxiv.org/abs/1008.4318}{arXiv:1008.4318}.

\bibitem{Baez}
J. C. Baez, The octonions, \emph{Bulletin of the American Mathematical Society}
\textbf{39} (2002), 145--205.
\href{https://arxiv.org/abs/math/0105155}{arXiv:math/0105155}.

\bibitem{Swan}
R. G. Swan, \emph{Tarski's Principle and the Elimination of Quantifiers}, expository
notes. \url{https://www.math.uchicago.edu/~swan/expo/Tarski.pdf}.

\bibitem{Mathlib}
The mathlib community, \emph{Mathlib.Algebra.Quaternion}, online documentation,
consulted 21 September 2026.
\url{https://leanprover-community.github.io/mathlib4_docs/Mathlib/Algebra/Quaternion.html}.

\bibitem{Repo}
V. Reshetnikov, \emph{Surreal}, repository README, module inventory and
\texttt{Surreal/HahnSeries/Neumann.lean}, commit
\texttt{39f2be6667ade51bca2b45daa47e289d69c09764}, accessed September 21, 2026.
\url{https://github.com/VladimirReshetnikov/Surreal}.

\bibitem{RepoDocs}
V. Reshetnikov, \emph{The surreal and surcomplex reports}, \texttt{docs/README.md} in the
same repository and commit.

\bibitem{RepoTrig}
\emph{Trigonometry on the Surcomplex Plane: Canonical Phases, Arbitrary-Scale Geometry,
and Infinitesimal Degeneration}, unpublished merged report in the same repository,
\texttt{docs/surcomplex/trigonometry/article.tex}.

\bibitem{RepoSpec}
\emph{Surcomplex Spectral Theory: Singular Values, Infinitesimal Rank, and Multiscale
Stability}, unpublished merged report in the same repository,
\texttt{docs/surcomplex/spectral-theory/article.tex}.

\bibitem{RepoDiff}
\emph{Differential Algebra and Differential Equations over the Surreal and Surcomplex
Numbers: The Finite-Primitive Criterion, the Phase Obstruction, and Picard--Vessiot
Theory}, unpublished merged report in the same repository,
\texttt{docs/surcomplex/differential-equations/article.tex}.

\bibitem{RepoPhys}
\emph{Surreal Scalars and Spacetime}, unpublished merged report in the same repository,
\texttt{docs/physics/surreal-scalars-and-spacetime/}, cited here only for its
\emph{commutative} surcomplex Ehrlich--Kaplan exponential with kernel $2\pi\ii\Oz$.

\end{thebibliography}
\endgroup
\end{document}
